% Economic growth and Development(continue) — Economics Upper Sixth — Cours signé OnBuch+
% Official MINESEC syllabus. Compiler avec : tectonic ECO08-economic-growth-and-development-continue.tex
\newcommand{\DOCMATIERE}{Economics}
\newcommand{\DOCNIVEAU}{Upper Sixth}
\newcommand{\DOCMODULE}{Upper Sixth — Economics}
\newcommand{\DOCLECON}{8}
\newcommand{\DOCTITRE}{Economic growth and Development(continue)}
% OnBuch+ — préambule « cours signés » ENRICHI (compilé avec tectonic / XeLaTeX)
% Système visuel « Pop » d'OnBuch+ : fond crème (jamais blanc), bordures encre
% épaisses, ombres « dures » décalées, titres Archivo Black en capitales.
% Version MATHÉMATIQUES : graphes (pgfplots/tikz), boîtes pédagogiques
% (définition, propriété, méthode, attention, le savais-tu, exercice résolu,
% exercice, TP, bilan), page de garde et signature OnBuch+ ancrée en bas.
%
% Macros attendues dans main.tex AVANT \input{preamble} :
%   \DOCMATIERE  \DOCNIVEAU  \DOCMODULE  \DOCLECON (numéro)  \DOCTITRE
\documentclass[11pt]{article}

\usepackage{fontspec}
\usepackage[english]{babel}
\usepackage[a4paper,margin=2cm,top=3.4cm,bottom=2.8cm,headheight=1.4cm,footskip=1.3cm]{geometry}
\usepackage{amsmath,amssymb,amsfonts}
\usepackage{mathtools} % \xmapsto, \coloneqq... souvent utilisés par les modèles
\usepackage{cancel} % simplifications barrées (bilans de forces)
% \abs{z} (module, valeur absolue) et \norm{u} (norme), utilisés par les modèles
\providecommand{\abs}{}\let\abs\relax\DeclarePairedDelimiter{\abs}{\lvert}{\rvert}
\providecommand{\norm}{}\let\norm\relax\DeclarePairedDelimiter{\norm}{\lVert}{\rVert}
\usepackage[table]{xcolor} % [table] : fournit aussi \rowcolors (alternance de lignes)
\usepackage{pagecolor}
\usepackage{tikz}
\usetikzlibrary{shapes.symbols,circuits.logic.US,circuits.logic.IEC,arrows.meta,positioning,calc,shapes.geometric,shapes.misc,shapes.arrows,decorations.pathmorphing,decorations.pathreplacing,decorations.markings,patterns,shadows,mindmap,trees,fit,backgrounds,matrix,intersections,angles,quotes}
\usepackage{pgfplots}
\usepackage[version=4]{mhchem} % équations nucléaires \ce{^{235}_{92}U}
\usepackage[european,siunitx]{circuitikz} % schémas électriques
\pgfplotsset{compat=1.18}
\usepgfplotslibrary{fillbetween,groupplots,polar,statistics,dateplot} % aires entre courbes (calcul intégral)
\usepackage{enumitem}
\usepackage{fancyhdr}
% [most] charge toutes les bibliothèques tcolorbox (listings, minted, raster,
% documentation...) et gonfle inutilement la mémoire TeX sur un document long
% (~300 boîtes) : on ne charge que ce qui est réellement utilisé.
\usepackage[skins,breakable]{tcolorbox}
\usepackage{graphicx}
\usepackage{colortbl}
\usepackage{booktabs}
\usepackage{tabularx}
\usepackage{diagbox} % case d'en-tête diagonale des tableaux à double entrée
% Colonnes extensibles centrée / gauche / droite (C, L, R), souvent utilisées par les modèles
\newcolumntype{C}{>{\centering\arraybackslash}X}
\newcolumntype{L}{>{\raggedright\arraybackslash}X}
\newcolumntype{R}{>{\raggedleft\arraybackslash}X}
\usepackage{array}
\usepackage{multicol}
\usepackage{siunitx}
% parse-numbers=false : accepte aussi \SI{2,5 \times 10^{-3}}{...}
\sisetup{locale=UK,output-decimal-marker={.},group-minimum-digits=4,parse-numbers=false}
\DeclareSIUnit\ohm{\ensuremath{\symup{\Omega}}} % U+2126 absent de la police
\DeclareSIUnit\division{div} % réglages d'oscilloscope (ms/div, V/div)
\DeclareSIUnit\year{yr}\DeclareSIUnit\annum{yr}\DeclareSIUnit\Ma{Ma}\DeclareSIUnit\tr{tr} % tours (vitesse de rotation, ex. \SI{1500}{\tr\per\minute})
\usepackage{qrcode}
% \degree : commande fréquemment utilisée par les modèles pour un angle ou une
% température en dehors de \SI{}{} (non fournie par les paquets chargés).
\providecommand{\degree}{^{\circ}}
% \qed (fin de démonstration), utilisé par les modèles sans amsthm
\providecommand{\qed}{\hfill$\square$}
% \barre : double barre des valeurs interdites dans un tableau de variations (tkz-tab)
\providecommand{\barre}{\ensuremath{\|}}
% environnement proof (amsthm non chargé) utilisé par les modèles
\ifdefined\proof\else\newenvironment{proof}[1][Proof]{\par\noindent\textit{#1.}\ }{\qed\par}\fi
\usepackage{lastpage}
\usepackage{hyperref}
\hypersetup{hidelinks}

% ============================================================
% Palette « Pop » OnBuch+ — mêmes hex que la classe P (plus_ui.dart)
% ============================================================
\definecolor{poporange}{HTML}{F59321}
\definecolor{popdark}{HTML}{E07A0C}
\definecolor{popcream}{HTML}{FFF6E8}
\definecolor{popink}{HTML}{1C1714}
\definecolor{poppurple}{HTML}{7A5AE0}
\definecolor{popgreen}{HTML}{2E9E6A}
\definecolor{poppink}{HTML}{E0457B}
\definecolor{popblue}{HTML}{2F7DE1}
\definecolor{popgold}{HTML}{C98A12}
\definecolor{popmuted}{HTML}{8A7F76}
\definecolor{popline}{HTML}{EADFD2}
\definecolor{popgray}{HTML}{8A7F76}
\colorlet{popgrey}{popgray}
% Alias tolérants (le modèle nomme parfois les couleurs différemment)
\colorlet{popwhite}{white}
\colorlet{popblack}{popink}
\colorlet{popred}{poppink}
\colorlet{popyellow}{popgold}
% Teintes claires pour fonds de tableaux / zones
\colorlet{popblueL}{popblue!10!white}
\colorlet{poporangeL}{poporange!14!white}
\colorlet{popgreenL}{popgreen!12!white}
\colorlet{poppurpleL}{poppurple!12!white}
\colorlet{poppinkL}{poppink!12!white}
\colorlet{popgoldL}{popgold!14!white}

\pagecolor{popcream}

% ============================================================
% Typographie
% ============================================================
\defaultfontfeatures{Ligatures=TeX}
\setmainfont{PlusJakartaSans-Medium.ttf}[Path=../../../fonts/,BoldFont=PlusJakartaSans-Bold.ttf]
% Maths en sans-serif (Fira Math) avec chiffres et lettres droites de Plus
% Jakarta, pour que les formules se fondent dans le texte.
\usepackage[math-style=ISO,bold-style=ISO]{unicode-math}
\setmathfont{FiraMath-Regular.otf}[Path=../../../fonts/]
\setmathfont{PlusJakartaSans-Medium.ttf}[Path=../../../fonts/,range={up/num,up/{Latin,latin},\mathup}]
% (icomma retiré : décimales avec point en anglais)
% Caractères mathématiques Unicode absents des polices de texte (Plus Jakarta,
% Archivo Black) : rendus par leur équivalent mathématique, en texte comme en
% formule — sinon ils s'affichent en carré vide (ex. « Arithmétique dans ℤ »).
\usepackage{newunicodechar}
\newunicodechar{ℝ}{\ensuremath{\mathbb{R}}}
\newunicodechar{ℤ}{\ensuremath{\mathbb{Z}}}
\newunicodechar{ℂ}{\ensuremath{\mathbb{C}}}
\newunicodechar{ℕ}{\ensuremath{\mathbb{N}}}
\newunicodechar{ℚ}{\ensuremath{\mathbb{Q}}}
\newunicodechar{𝔻}{\ensuremath{\mathbb{D}}}
\newunicodechar{α}{\ensuremath{\alpha}}
\newunicodechar{β}{\ensuremath{\beta}}
\newunicodechar{γ}{\ensuremath{\gamma}}
\newunicodechar{δ}{\ensuremath{\delta}}
\newunicodechar{ε}{\ensuremath{\varepsilon}}
\newunicodechar{θ}{\ensuremath{\theta}}
\newunicodechar{λ}{\ensuremath{\lambda}}
\newunicodechar{μ}{\ensuremath{\mu}}
\newunicodechar{σ}{\ensuremath{\sigma}}
\newunicodechar{τ}{\ensuremath{\tau}}
\newunicodechar{φ}{\ensuremath{\varphi}}
\newunicodechar{ψ}{\ensuremath{\psi}}
\newunicodechar{ω}{\ensuremath{\omega}}
\newunicodechar{Σ}{\ensuremath{\Sigma}}
% majuscules produites par \MakeUppercase dans les titres (α-aminés -> Α)
\newunicodechar{Α}{\ensuremath{\alpha}}
\newunicodechar{Β}{\ensuremath{\beta}}
\newunicodechar{Γ}{\ensuremath{\Gamma}}
\newunicodechar{Δ}{\ensuremath{\Delta}}
\newunicodechar{Ε}{\ensuremath{\varepsilon}}
\newunicodechar{Θ}{\ensuremath{\Theta}}
\newunicodechar{Λ}{\ensuremath{\Lambda}}
\newunicodechar{Μ}{\ensuremath{\mu}}
\newunicodechar{Τ}{\ensuremath{\tau}}
\newunicodechar{Φ}{\ensuremath{\Phi}}
\newunicodechar{Ψ}{\ensuremath{\Psi}}
\newunicodechar{Ω}{\ensuremath{\Omega}}
\newunicodechar{↦}{\ensuremath{\mapsto}}
\newunicodechar{∈}{\ensuremath{\in}}
\newunicodechar{∉}{\ensuremath{\notin}}
\newunicodechar{∧}{\ensuremath{\wedge}}
\newunicodechar{⇔}{\ensuremath{\Leftrightarrow}}
\newunicodechar{⟺}{\ensuremath{\Longleftrightarrow}}
\newunicodechar{⇒}{\ensuremath{\Rightarrow}}
\newunicodechar{⟹}{\ensuremath{\Longrightarrow}}
\newunicodechar{∘}{\ensuremath{\circ}}
\newunicodechar{∪}{\ensuremath{\cup}}
\newunicodechar{∩}{\ensuremath{\cap}}
\newunicodechar{≡}{\ensuremath{\equiv}}
\newunicodechar{⊂}{\ensuremath{\subset}}
\newunicodechar{⊥}{\ensuremath{\perp}}
\newunicodechar{∃}{\ensuremath{\exists}}
\newunicodechar{∀}{\ensuremath{\forall}}
\newunicodechar{∛}{\ensuremath{\sqrt[3]{\,}}}
\newunicodechar{∜}{\ensuremath{\sqrt[4]{\,}}}
\newunicodechar{⃗}{\ensuremath{{}^{\to}}}
\newunicodechar{ⁿ}{\ifmmode^{n}\else\textsuperscript{\ensuremath{n}}\fi}
\newunicodechar{ˣ}{\ifmmode^{x}\else\textsuperscript{\ensuremath{x}}\fi}
\newunicodechar{ᵇ}{\ifmmode^{b}\else\textsuperscript{\ensuremath{b}}\fi}
\newunicodechar{ʳ}{\ifmmode^{r}\else\textsuperscript{\ensuremath{r}}\fi}
\newunicodechar{ᵘ}{\ifmmode^{u}\else\textsuperscript{\ensuremath{u}}\fi}
\newunicodechar{ʸ}{\ifmmode^{y}\else\textsuperscript{\ensuremath{y}}\fi}
\newunicodechar{ᵃ}{\ifmmode^{a}\else\textsuperscript{\ensuremath{a}}\fi}
\newunicodechar{ᵉ}{\ifmmode^{e}\else\textsuperscript{\ensuremath{e}}\fi}
\newunicodechar{ᵏ}{\ifmmode^{k}\else\textsuperscript{\ensuremath{k}}\fi}
\newunicodechar{ᵖ}{\ifmmode^{p}\else\textsuperscript{\ensuremath{p}}\fi}
\newunicodechar{ᴺ}{\ifmmode^{N}\else\textsuperscript{\ensuremath{N}}\fi}
\newunicodechar{ᶜ}{\ifmmode^{c}\else\textsuperscript{\ensuremath{c}}\fi}
\newunicodechar{ʰ}{\ifmmode^{h}\else\textsuperscript{\ensuremath{h}}\fi}
\newunicodechar{ᵠ}{\ifmmode^{q}\else\textsuperscript{\ensuremath{q}}\fi}
\newunicodechar{ₙ}{\ifmmode_{n}\else\textsubscript{\ensuremath{n}}\fi}
\newunicodechar{ₐ}{\ifmmode_{a}\else\textsubscript{\ensuremath{a}}\fi}
\newunicodechar{ᵢ}{\ifmmode_{i}\else\textsubscript{\ensuremath{i}}\fi}
\newunicodechar{ᵣ}{\ifmmode_{r}\else\textsubscript{\ensuremath{r}}\fi}
\newunicodechar{ₖ}{\ifmmode_{k}\else\textsubscript{\ensuremath{k}}\fi}
\newunicodechar{ᵦ}{\ifmmode_{\beta}\else\textsubscript{\ensuremath{\beta}}\fi}
\newunicodechar{ₓ}{\ifmmode_{x}\else\textsubscript{\ensuremath{x}}\fi}
\newfontfamily\popdisplay{ArchivoBlack-Regular.ttf}[Path=../../../fonts/]
\newfontfamily\poplogo{SpaceGrotesk-Bold.ttf}[Path=../../../fonts/]
\newfontfamily\popsemi{PlusJakartaSans-SemiBold.ttf}[Path=../../../fonts/]
\newfontfamily\popxbold{PlusJakartaSans-ExtraBold.ttf}[Path=../../../fonts/]
\newfontfamily\popxboldls{PlusJakartaSans-ExtraBold.ttf}[Path=../../../fonts/,LetterSpace=3.0]

\setlist[enumerate]{leftmargin=1.7em,itemsep=5pt,topsep=4pt}
\setlist[itemize]{leftmargin=1.7em,itemsep=4pt,topsep=4pt,label={\color{poporange}\textbullet}}
\setlength{\parindent}{0pt}
\setlength{\parskip}{5pt}
\renewcommand{\arraystretch}{1.25}

% pgfplots : style Pop par défaut
\pgfplotsset{
  every axis/.append style={axis line style={popink,line width=1pt},tick style={popink},
    grid style={popline,line width=0.5pt},label style={font=\small},tick label style={font=\footnotesize}},
  popcurve/.style={poporange,line width=1.8pt,no markers,smooth},
}
% Même style disponible dans un \draw TikZ simple (hors pgfplots)
\tikzset{popcurve/.style={poporange,line width=1.8pt}}
\tikzset{popgrid/.style={popline,very thin}}
\colorlet{popcurve}{poporange} % le nom du style est parfois utilisé comme couleur

% ============================================================
% Pill / squiggle
% ============================================================
\newcommand{\poppill}[3]{%
  \tikz[baseline=-0.55ex]\node[draw=popink,line width=1.2pt,rounded corners=2.6pt,%
    fill=#2,inner xsep=6.5pt,inner ysep=3pt]{%
    \popxboldls\fontsize{7}{7}\selectfont\color{#3}\MakeUppercase{#1}};%
}
\newcommand{\popsquiggle}[1]{%
  \tikz[baseline=-0.4ex]{
    \draw[#1,line width=2.6pt,line cap=round]
      plot [smooth] coordinates {(0,0.09) (0.34,0.01) (0.68,0.21) (1.02,0.01) (1.36,0.10)};
  }%
}
% Mot-clé surligné (marqueur orange sous le texte)
\newcommand{\cle}[1]{{\popxbold\color{popdark}#1}}

% ============================================================
% En-tête / pied
% ============================================================
\pagestyle{fancy}
\fancyhf{}
\renewcommand{\headrulewidth}{0pt}
\renewcommand{\footrulewidth}{0pt}
\lhead{\raisebox{-0.15ex}{\poplogo\fontsize{13}{13}\selectfont\color{popink}OnBuch\color{poporange}+}}
\rhead{\poppill{\DOCMATIERE\ \textbullet\ \DOCNIVEAU\ \textbullet\ Chapter \DOCLECON}{white}{popink}}
\lfoot{\popxbold\fontsize{7.6}{7.6}\selectfont\color{popmuted}\MakeUppercase{Signed course OnBuch+}\ \textbullet\ {\popsemi\DOCTITRE}}
\rfoot{\tikz[baseline=-0.6ex]\node[rounded corners=8.5pt,fill=popink,minimum height=17pt,minimum width=17pt,inner xsep=5pt,inner sep=0]{\popxbold\fontsize{8}{8}\selectfont\color{popcream}\thepage\,/\,\pageref*{LastPage}};}

% ============================================================
% Titres
% ============================================================
\newcommand{\chaptitre}[2]{%
  \begin{tcolorbox}[enhanced,colback=poporange,colframe=popink,boxrule=2.6pt,arc=11pt,
      drop shadow=popink,left=15pt,right=15pt,top=12pt,bottom=11pt,before skip=3pt,after skip=18pt]
    \poppill{#2}{popink}{popcream}\par\vspace{9pt}
    {\raggedright\hyphenpenalty=10000\exhyphenpenalty=10000\popdisplay\fontsize{22}{24}\selectfont\color{popcream}\MakeUppercase{#1}\par}
  \end{tcolorbox}
}
% Section numérotée : pastille orange avec le numéro + titre Archivo
\newcounter{popsec}
\newcommand{\coursec}[1]{%
  \stepcounter{popsec}\setcounter{popsub}{0}%
  \par\vspace{16pt}\noindent
  \begin{minipage}{\linewidth}
  \tikz[baseline=-0.7ex]\node[draw=popink,line width=1.6pt,fill=poporange,rounded corners=4pt,minimum size=22pt,inner sep=0]{\popdisplay\fontsize{12}{12}\selectfont\color{popcream}\thepopsec};%
  \hspace{8pt}{\popdisplay\fontsize{15}{17}\selectfont\color{popink}\MakeUppercase{#1}}\par
  \vspace{4pt}\noindent{\color{poporange}\rule{36pt}{3.6pt}}
  \end{minipage}\par\nopagebreak\vspace{8pt}
}
\newcounter{popsub}
\newcommand{\courssub}[1]{%
  \stepcounter{popsub}%
  \par\vspace{9pt}\noindent{\popxbold\fontsize{11.5}{13}\selectfont\color{popink}{\color{poporange}\thepopsec.\thepopsub}\hspace{6pt}#1}\par\nopagebreak\vspace{4pt}
}

% ============================================================
% Boîtes pédagogiques — même recette : fond blanc, bordure épaisse colorée,
% ombre dure encre, badge plein. Titre optionnel [#1].
% ============================================================
\tcbset{popbase/.style={breakable,enhanced,colback=white,boxrule=2.3pt,arc=9pt,
  shadow={3pt}{-3pt}{0pt}{popink},left=14pt,right=14pt,top=11pt,bottom=11pt,before skip=11pt,after skip=15pt,
  fontupper=\popsemi\fontsize{10.6}{14.5}\selectfont\color{popink},
  fontlower=\popsemi\fontsize{10.6}{14.5}\selectfont\color{popink}}}
% Coche dessinée (le glyphe ✓ n'existe pas dans les polices du document)
\newcommand{\popcheck}[1]{\tikz[baseline=-0.5ex]\draw[#1,line width=1.6pt,line cap=round,line join=round](-0.13,0.0)--(-0.04,-0.09)--(0.13,0.1);}
\providecommand{\coche}{\popcheck{popgreen}}
\providecommand{\resultat}[1]{\par\smallskip\noindent\fcolorbox{popgreen}{popgreenL}{\parbox{\dimexpr\linewidth-2\fboxsep-2\fboxrule\relax}{\textbf{Result.} #1}}\par\smallskip}
\newcommand{\popboxhead}[3]{\poppill{#1}{#2}{white}\ifx\relax#3\relax\else\hspace{6pt}{\popxbold\fontsize{10.6}{12}\selectfont\color{popink}#3}\fi\par\vspace{7pt}}

\newtcolorbox{aretenir}[1][]{popbase,colframe=popgreen,before upper={\popboxhead{Key points}{popgreen}{#1}}}
\newtcolorbox{exemplebox}[1][]{popbase,colframe=poporange,before upper={\popboxhead{Exemple}{poporange}{#1}}}
\newtcolorbox{definition}[1][]{popbase,colframe=popblue,before upper={\popboxhead{Definition}{popblue}{#1}}}
\newtcolorbox{propriete}[1][]{popbase,colframe=poppurple,before upper={\popboxhead{Property}{poppurple}{#1}}}
\newtcolorbox{methode}[1][]{popbase,colframe=popgold,before upper={\popboxhead{Method}{popgold}{#1}}}
\newtcolorbox{attention}[1][]{popbase,colframe=poppink,before upper={\popboxhead{Warning}{poppink}{#1}}}
\newtcolorbox{experience}[1][]{popbase,colframe=popblue,colback=popblueL,before upper={\popboxhead{Experiment}{popblue}{#1}}}
% « Le savais-tu ? » : carte violette pleine (contexte, culture, Cameroun)
\newtcolorbox{savaistu}[1][]{popbase,colback=poppurple,colframe=popink,
  fontupper=\popsemi\fontsize{10.4}{14.2}\selectfont\color{white},
  before upper={\poppill{Did you know?}{white}{poppurple}\ifx\relax#1\relax\else\hspace{6pt}{\popxbold\color{white}#1}\fi\par\vspace{7pt}}}
% Exercice résolu : énoncé en haut, \tcblower puis la solution sur fond crème
\newcounter{popexo}\newcounter{popexr}
\newtcolorbox{exoresolu}[1][]{popbase,colframe=poporange,colbacklower=popcream,
  segmentation style={popink,line width=1.2pt,dashed},
  before upper={\stepcounter{popexr}\popboxhead{Worked example \thepopexr}{poporange}{#1}},
  before lower={\poppill{Solution}{popink}{popcream}\par\vspace{6pt}}}
% Exercice (non résolu en place) — la correction est dans la partie Corrigés
\newtcolorbox{exercice}[1][]{popbase,colframe=popink,
  before upper={\stepcounter{popexo}\popboxhead{Exercise \thepopexo}{popink}{#1}}}
% Corrigé
\newtcolorbox{corrige}[1][]{popbase,colframe=popgreen,colback=popgreenL,
  before upper={\popboxhead{Solution}{popgreen}{#1}}}
% Objectifs / prérequis / bilan
\newtcolorbox{objectifs}{popbase,colframe=popink,colback=poporangeL,
  before upper={\popboxhead{Chapter objectives}{popink}{}}}
\newtcolorbox{prerequis}{popbase,colframe=popink,colback=popblueL,
  before upper={\popboxhead{Prerequisites}{popblue}{}}}
\newtcolorbox{bilan}{popbase,colframe=popink,colback=white,boxrule=2.8pt,
  before upper={\popboxhead{Fiche bilan}{poporange}{L'essentiel à connaître pour l'examen}}}
% Figure encadrée avec légende
\newtcolorbox{popfigure}[1][]{enhanced,breakable=false,colback=white,colframe=popline,boxrule=1.4pt,arc=8pt,
  left=10pt,right=10pt,top=10pt,bottom=8pt,before skip=10pt,after skip=12pt,halign=center,
  lower separated=false,halign lower=center,
  fontlower=\popsemi\fontsize{9}{11.5}\selectfont\color{popmuted},
  #1}
\newcommand{\legende}[1]{\par\vspace{6pt}{\popsemi\fontsize{9}{11.5}\selectfont\color{popmuted}#1}}

% ============================================================
% Signature OnBuch+
% ============================================================
\newcommand{\poplogoinline}[1]{{\poplogo\fontsize{#1}{#1}\selectfont\color{popink}OnBuch\color{poporange}+}}
\newcommand{\signatureonbuch}{%
  \begin{tcolorbox}[enhanced,colback=popink,colframe=popink,boxrule=2.2pt,arc=10pt,
      drop shadow=poporange,left=14pt,right=14pt,top=10pt,bottom=10pt,before skip=0pt,after skip=0pt]
    \tikz[baseline=-0.7ex]\node[circle,fill=popgreen,draw=popcream,line width=1.3pt,minimum size=22pt,inner sep=0]{\popcheck{white}};%
    \hspace{9pt}\begin{minipage}[c]{0.62\linewidth}
      {\popxbold\fontsize{12}{13}\selectfont\color{popcream}Signed\ \textbullet\ {\poplogo OnBuch\color{poporange}+}}\par\vspace{2pt}
      {\raggedright\popsemi\fontsize{8}{10.5}\selectfont\color{popline}\MakeUppercase{OnBuch+ teaching team}\\\MakeUppercase{Follows the official MINESEC syllabus}\par}
    \end{minipage}\hfill
    {\poplogo\fontsize{22}{22}\selectfont\color{popcream}OnBuch\color{poporange}+}
  \end{tcolorbox}%
}

% ============================================================
% Page de garde
% #1 titre · #2 matière · #3 niveau · #4 module · #5 n° leçon · #6 durée ·
% #7 description / objectifs (texte ou itemize)
% ============================================================
\newcommand{\pagegarde}[7]{%
  \thispagestyle{empty}
  \newgeometry{a4paper,margin=1.7cm,top=1.6cm,bottom=1.4cm}%
  \begin{tikzpicture}[remember picture,overlay]
    \fill[poporange!18] ([xshift=-3.2cm,yshift=-3.6cm]current page.north east) circle (4.6cm);
    \fill[poppurple!14] ([xshift=2.4cm,yshift=7.6cm]current page.south west) circle (3.8cm);
  \end{tikzpicture}%
  {\poplogo\fontsize{30}{30}\selectfont\color{popink}OnBuch\color{poporange}+}\hfill
  \raisebox{0.6ex}{\poppill{Signed course \textbullet\ Premium}{popink}{popcream}}\par
  \vspace{1pt}\popsquiggle{poppurple}\par
  \vspace{8pt}
  \begin{tcolorbox}[enhanced,colback=poporange,colframe=popink,boxrule=2.8pt,arc=13pt,
      drop shadow=popink,left=18pt,right=18pt,top=12pt,bottom=13pt,after skip=10pt]
    \poppill{#2 \textbullet\ #3}{popink}{popcream}\hspace{5pt}\poppill{#4}{white}{popink}\par\vspace{8pt}
    {\popdisplay\fontsize{14}{14}\selectfont\color{popink}CHAPTER #5}\par\vspace{4pt}
    {\raggedright\hyphenpenalty=10000\exhyphenpenalty=10000\popdisplay\fontsize{25}{27}\selectfont\color{popcream}\MakeUppercase{#1}\par}
    \vspace{8pt}
    {\popxbold\fontsize{9.5}{9.5}\selectfont\color{popink}\MakeUppercase{Suggested time: #6}}
  \end{tcolorbox}
  \begin{tcolorbox}[enhanced,colback=white,colframe=popink,boxrule=2.2pt,arc=10pt,
      drop shadow=popink,left=16pt,right=16pt,top=10pt,bottom=10pt,after skip=8pt]
    \poppill{In this chapter}{poppurple}{white}\par\vspace{5pt}
    \popsemi\fontsize{9.3}{12.6}\selectfont\color{popink}#7
  \end{tcolorbox}
  \noindent\begin{minipage}[t]{0.487\linewidth}
    \begin{tcolorbox}[enhanced,colback=popgreen,colframe=popink,boxrule=2.2pt,arc=10pt,
        drop shadow=popink,left=11pt,right=11pt,top=10pt,bottom=10pt,height=3.1cm,valign=center]
      \tcbox[colback=white,colframe=popink,boxrule=1.3pt,arc=5pt,boxsep=0pt,nobeforeafter,
        left=3pt,right=3pt,top=3pt,bottom=3pt,valign=center]{\qrcode[height=1.9cm]{https://play.google.com/store/apps/details?id=com.luvvix.onbuch&hl=en}}%
      \hspace{8pt}\begin{minipage}[c]{0.5\linewidth}
        {\popdisplay\fontsize{11}{12.5}\selectfont\color{white}\MakeUppercase{Download OnBuch}}\par\vspace{4pt}
        {\raggedright\popsemi\fontsize{8.4}{11}\selectfont\color{white}Free \textbullet\ Google Play\\AI tutor, past papers, results\par}
      \end{minipage}
    \end{tcolorbox}
  \end{minipage}\hfill
  \begin{minipage}[t]{0.487\linewidth}
    \begin{tcolorbox}[enhanced,colback=poppurple,colframe=popink,boxrule=2.2pt,arc=10pt,
        drop shadow=popink,left=11pt,right=11pt,top=10pt,bottom=10pt,height=3.1cm,valign=center]
      \tikz[baseline=-0.7ex]\node[circle,fill=white,draw=popink,line width=1.3pt,minimum size=1.6cm,inner sep=0]{\popdisplay\fontsize{24}{24}\selectfont\color{poppurple}+};%
      \hspace{8pt}\begin{minipage}[c]{0.6\linewidth}
        {\popdisplay\fontsize{11}{12.5}\selectfont\color{white}\MakeUppercase{Go OnBuch+}}\par\vspace{4pt}
        {\raggedright\popsemi\fontsize{8.4}{11}\selectfont\color{white}All signed courses, unlimited AI tutor, PDF sheets — OnBuch+ tab in the app\par}
      \end{minipage}
    \end{tcolorbox}
  \end{minipage}
  \vspace{8pt}
  \popcontenu
  \vfill
  \signatureonbuch
  \restoregeometry
  \newpage
}

% Rangée « Ce cours contient » (page de garde)
\newcommand{\poptile}[3]{% #1 couleur #2 gros texte #3 légende
  \begin{tcolorbox}[enhanced,colback=#1,colframe=popink,boxrule=2pt,arc=9pt,shadow={2.5pt}{-2.5pt}{0pt}{popink},
      left=6pt,right=6pt,top=9pt,bottom=9pt,halign=center,valign=center,height=1.75cm,width=\linewidth,nobeforeafter]
    {\popdisplay\fontsize{12.5}{13}\selectfont\color{white}\MakeUppercase{#2}}\par\vspace{3pt}
    {\centering\popsemi\fontsize{7.8}{9.5}\selectfont\color{white}#3\par}
  \end{tcolorbox}}
\newcommand{\popcontenu}{%
  \par\noindent{\popxboldls\fontsize{7.5}{7.5}\selectfont\color{popmuted}THIS SIGNED COURSE CONTAINS}\par\vspace{7pt}
  \noindent\begin{minipage}[t]{0.235\linewidth}\poptile{popblue}{Course}{complete, illustrated, step by step}\end{minipage}\hfill
  \begin{minipage}[t]{0.235\linewidth}\poptile{poporange}{Methods}{and worked examples}\end{minipage}\hfill
  \begin{minipage}[t]{0.235\linewidth}\poptile{poppink}{Exercises}{exam style + solutions}\end{minipage}\hfill
  \begin{minipage}[t]{0.235\linewidth}\poptile{popgreen}{Summary sheet}{the essentials to remember}\end{minipage}\par}

% Sommaire
\newtcolorbox{sommaire}{enhanced,colback=white,colframe=popink,boxrule=2.2pt,arc=10pt,
  drop shadow=popink,left=18pt,right=18pt,top=13pt,bottom=13pt,before skip=6pt,after skip=20pt,
  before upper={\poppill{Contents}{poppurple}{white}\par\vspace{9pt}\popsemi\fontsize{10.6}{14.5}\selectfont\color{popink}}}

% Signature de fin de cours — poussée tout en bas de la dernière page
\newcommand{\findecours}{%
  \par\vfill\nopagebreak
  \begin{center}\popsquiggle{poporange}\end{center}\vspace{4pt}
  \signatureonbuch
}
\AtBeginDocument{\def\llbracket{\mathopen{[\mkern-3.5mu[}}\def\rrbracket{\mathclose{]\mkern-3.5mu]}}} % intervalles d'entiers (glyphe ⟦ absent de la police)
% Glyphes absents de Fira Math : redéfinis à partir de glyphes disponibles
\AtBeginDocument{%
  \def\setminus{\mathbin{\backslash}}\let\smallsetminus\setminus
  \def\vdots{\mathord{\vbox{\baselineskip3.2pt\lineskiplimit0pt\kern4pt\hbox{.}\hbox{.}\hbox{.}}}}%
  \def\ddots{\mathinner{\mkern1mu\raise6.5pt\hbox{.}\mkern2mu\raise3.6pt\hbox{.}\mkern2mu\raise0.7pt\hbox{.}\mkern1mu}}%
  \def\triangle{\mathord{\scalebox{0.9}{$\symup{\Delta}$}}}%
  \def\nabla{\mathord{\raisebox{\depth}{\scalebox{1}[-1]{$\symup{\Delta}$}}}}%
  \def\checkmark{\popcheck{popdark}}%
  \def\star{\mathbin{\ast}}\def\bigstar{\mathord{\ast}}%
  \def\diamond{\mathbin{\scalebox{0.6}{$\Diamond$}}}%
  \def\bigcup{\mathop{\mathchoice{\raisebox{-0.3em}{\scalebox{1.7}{$\cup$}}}{\scalebox{1.25}{$\cup$}}{\cup}{\cup}}\displaylimits}%
  \def\bigcap{\mathop{\mathchoice{\raisebox{-0.3em}{\scalebox{1.7}{$\cap$}}}{\scalebox{1.25}{$\cap$}}{\cap}{\cap}}\displaylimits}%
  \def\lesseqqgtr{\mathrel{\lessgtr}}\def\bigoplus{\mathop{\oplus}\displaylimits}%
}
\providecommand{\ssi}{if and only if} % abréviation fréquente
\AtBeginDocument{\providecommand{\wideparen}{\overparen}} % arc de cercle

% Fonctions usuelles que les modèles utilisent sans que amsmath les fournisse
\providecommand{\arsinh}{\operatorname{arsinh}}\providecommand{\arcosh}{\operatorname{arcosh}}
\providecommand{\artanh}{\operatorname{artanh}}\providecommand{\arsech}{\operatorname{arsech}}
\providecommand{\arcsch}{\operatorname{arcsch}}\providecommand{\arcoth}{\operatorname{arcoth}}
\providecommand{\arcsinh}{\operatorname{arsinh}}\providecommand{\arccosh}{\operatorname{arcosh}}
\providecommand{\arctanh}{\operatorname{artanh}}\providecommand{\sech}{\operatorname{sech}}
\providecommand{\csch}{\operatorname{csch}}\providecommand{\arcsec}{\operatorname{arcsec}}
\providecommand{\arccsc}{\operatorname{arccsc}}\providecommand{\arccot}{\operatorname{arccot}}
\providecommand{\sgn}{\operatorname{sgn}}\providecommand{\lcm}{\operatorname{lcm}}
\providecommand{\Var}{\operatorname{Var}}\providecommand{\Cov}{\operatorname{Cov}}
\providecommand{\permil}{\ensuremath{{}^{0}\!/_{00}}}\providecommand{\textperthousand}{\ensuremath{{}^{0}\!/_{00}}}

%% FIGFIX-BEGIN v5 — correctifs globaux des figures (docs/onbuch-generation/tools/figfix)
\usepackage{adjustbox}\usepackage{etoolbox}\tcbuselibrary{raster}
% toute figure TikZ/pgfplots plus large que la ligne est réduite à la largeur disponible
\BeforeBeginEnvironment{tikzpicture}{\begin{adjustbox}{max width=\linewidth}}
\AfterEndEnvironment{tikzpicture}{\end{adjustbox}}
% légendes pgfplots lisibles par-dessus les courbes
\pgfplotsset{every axis/.append style={legend style={fill=white,fill opacity=0.92,text opacity=1,draw=black!15,inner sep=3pt}}}
% étiquettes d'arêtes (syntaxe edge["..."]) avec fond blanc
\tikzset{every edge quotes/.append style={fill=white,inner sep=1.2pt}}
%% FIGFIX-END
\begin{document}

\pagegarde{Economic growth and Development(continue)}{Economics}{Upper Sixth}{Upper Sixth — Economics}{8}{21 periods}{This chapter completes the study of economic growth and development by examining development strategies, industrialisation, the business cycle and the major macroeconomic problems — unemployment, inflation and deflation — together with the policy instruments (privatisation, nationalisation, regional and population policies) governments use to manage them. It blends theory with Cameroonian case studies to prepare students fully for the GCE Advanced Level Economics examination.
\vspace{4pt}
\begin{multicols}{2}\small
\begin{itemize}
\item By the end of this chapter I can distinguish economic growth from economic development and explain the indicators and characteristics of development
\item By the end of this chapter I can explain the processes of industrialisation, deindustrialisation and reindustrialisation with real-world examples
\item By the end of this chapter I can analyse the development strategies of Newly Industrialised Countries and draw lessons for Cameroon
\item By the end of this chapter I can evaluate strategies for promoting second generation agriculture and agribusiness in Cameroon
\item By the end of this chapter I can describe and illustrate the phases of the trade (business) cycle and explain its causes
\item By the end of this chapter I can state and explain the main government economic objectives and the conflicts between them
\end{itemize}\end{multicols}}

\begin{sommaire}
\begin{multicols}{2}
\begin{enumerate}
\item Section 1: Economic Development — Meaning, Indicators and Characteristics
\item Section 2: Industrialisation, Deindustrialisation and Reindustrialisation
\item Section 3: Development Strategies of Newly Industrialised Countries (NICs) — Further Studies 9
\item Section 4: Promoting Second Generation Agriculture and Agribusiness in Cameroon — Further Studies 10
\item Section 5: The Trade (Business) Cycle and Government Economic Objectives
\item Section 6: Unemployment — Types, Causes, Measurement and the Cameroonian Context
\item Section 7: Inflation and Deflation
\item Section 8: Privatisation, Nationalisation and Regional Policy
\item Section 9: Population — Growing and Ageing Populations, Optimum Population and Population Policy
\item Integration activity
\item Methods and worked examples
\item Exercises
\item Solutions to the exercises
\item Summary sheet
\end{enumerate}
\end{multicols}
\end{sommaire}

\begin{objectifs}
\begin{itemize}
\item By the end of this chapter I can distinguish economic growth from economic development and explain the indicators and characteristics of development
\item By the end of this chapter I can explain the processes of industrialisation, deindustrialisation and reindustrialisation with real-world examples
\item By the end of this chapter I can analyse the development strategies of Newly Industrialised Countries and draw lessons for Cameroon
\item By the end of this chapter I can evaluate strategies for promoting second generation agriculture and agribusiness in Cameroon
\item By the end of this chapter I can describe and illustrate the phases of the trade (business) cycle and explain its causes
\item By the end of this chapter I can state and explain the main government economic objectives and the conflicts between them
\item By the end of this chapter I can define, measure, classify and explain the causes and consequences of unemployment, with reference to Cameroon
\item By the end of this chapter I can explain the causes, types, measurement and effects of inflation and deflation and the policies used to control them
\item By the end of this chapter I can compare privatisation and nationalisation and assess their application in Cameroon
\item By the end of this chapter I can analyse population issues — growing and ageing populations, optimum population and population policy
\end{itemize}
\end{objectifs}

\begin{prerequis}
\begin{itemize}
\item Economic growth: meaning, measurement (GDP/GNP) and sources (Lower Sixth / ECO07)
\item Aggregate demand and aggregate supply analysis
\item The role of government in the economy and fiscal/monetary policy basics
\item Population concepts: birth rate, death rate, natural increase (Lower Sixth)
\item Price mechanism and market structures
\end{itemize}
\end{prerequis}

\newpage

\coursec{Section 1: Economic Development — Meaning, Indicators and Characteristics}

In 2023, a young graduate in Yaoundé sells phone credit by the roadside while inflation erodes the price of a bag of rice, and the government announces yet another agribusiness scheme in the West Region. Why do some countries like South Korea transform from poor to industrial giants in one generation while others stagnate? Why do prices rise, jobs disappear, and governments privatise companies one decade and nationalise them the next? This chapter takes you inside the big questions of economic development, the trade cycle, unemployment, inflation and population that shape the daily life of every Cameroonian. We begin with the most fundamental question of all: what does it really mean for a country to \emph{develop}?

\courssub{The Meaning of Economic Development}

\textbf{Intuition first.} Imagine two Cameroonian villages. In the first, total income has doubled in ten years because a foreign company opened a timber concession — but the school still has no roof, the health centre has no nurse, and most villagers remain subsistence farmers. In the second, income has also doubled, but in addition children are schooled, life expectancy has risen, small workshops process cassava locally, and women participate in decisions. Both villages experienced \cle{growth}; only the second experienced \cle{development}.

\begin{definition}[Economic development]
\cle{Economic development} is a multidimensional process involving major changes in social structures, popular attitudes and national institutions, as well as the acceleration of economic growth, the reduction of inequality and the eradication of absolute poverty. (This definition follows the approach of the development economist Michael Todaro.)
\end{definition}

\begin{definition}[Economic growth]
\cle{Economic growth} is the sustained increase over time in a country's real output of goods and services, usually measured by the percentage change in real GDP or real GDP per capita. It is a purely \emph{quantitative} change.
\end{definition}

The essential distinction is therefore between the \emph{quantitative} and the \emph{qualitative}:

\begin{center}
\begin{tabularx}{\linewidth}{|l|X|X|}
\hline
\rowcolor{popblueL} \textbf{Criterion} & \textbf{Economic growth} & \textbf{Economic development} \\
\hline
Nature of change & Quantitative (more output) & Qualitative and quantitative (better lives) \\
\hline
Main measure & Real GDP growth rate & HDI, poverty rates, literacy, health \\
\hline
Time horizon & Short to medium term & Long term \\
\hline
Distribution & Ignores who receives income & Central concern (equity) \\
\hline
Structures & Economy may stay unchanged & Requires structural transformation \\
\hline
\end{tabularx}
\end{center}

\courssub{The Dimensions of Development}

Development is not one thing but a bundle of simultaneous improvements:

\begin{itemize}
\item \textbf{Rising real incomes per head:} people can consume more goods and services.
\item \textbf{Poverty reduction:} a falling proportion of the population living below the poverty line (absolute poverty) and reduced vulnerability.
\item \textbf{Improved health and education:} longer life expectancy, lower infant mortality, higher literacy and school enrolment — what economists call improved \cle{human capital}.
\item \textbf{Structural transformation:} the shift of labour and output from subsistence agriculture towards industry and modern services, with urbanisation and rising productivity.
\item \textbf{Reduced inequality:} a fairer distribution of income, assets and opportunities, including between men and women and between regions.
\end{itemize}

\begin{savaistu}[Growth without development in Africa]
Several African oil exporters have recorded impressive GDP growth driven by crude petroleum, yet poverty, unemployment and malnutrition barely improved because oil revenues were concentrated in few hands and the rest of the economy stagnated. Economists call this the \cle{resource curse} or ``growth without development''. Cameroon's own oil boom of the 1980s illustrates the risk: high growth coincided with rising inequality and a neglected agricultural sector. We return to this in the exam extra below.
\end{savaistu}

\courssub{Indicators of Development}

\textbf{GDP and GNP per capita.} The oldest and simplest indicator is \cle{GDP per capita} (GDP divided by population). \cle{GNP per capita} adds net income from abroad, which matters for countries with large foreign-owned sectors or large diasporas. For international comparison, values are converted at \emph{purchasing power parity} (PPP) to account for differences in price levels.

\begin{attention}[Common mistake]
Never use GDP per capita \emph{alone} as proof that a country is developed. It says nothing about: (i) the \textbf{distribution} of income (a few oil barons can pull the average up while most citizens stay poor); (ii) \textbf{non-income dimensions} such as health, education and freedom; (iii) the \textbf{informal sector}, which is large in Cameroon and largely unrecorded; (iv) \textbf{environmental damage} caused by production; and (v) whether output consists of weapons or of schools. In an essay, always pair GDP per capita with at least one composite indicator and a discussion of distribution.
\end{attention}

\textbf{The Human Development Index (HDI).} To correct these weaknesses, the United Nations Development Programme (UNDP) created a composite index.

\begin{definition}[Human Development Index]
The \cle{Human Development Index} (HDI) is a composite index, ranging from 0 to 1, combining three dimensions of development:
\begin{enumerate}
\item \textbf{A long and healthy life}, measured by life expectancy at birth;
\item \textbf{Knowledge}, measured by mean and expected years of schooling;
\item \textbf{A decent standard of living}, measured by GNI per capita (PPP, in US dollars).
\end{enumerate}
Countries are grouped as very high, high, medium or low human development.
\end{definition}

\textbf{Other composite indices.} The \cle{Multidimensional Poverty Index} (MPI) counts the joint deprivations households face in health, education and living standards (for example, lack of clean water, electricity or a school-aged child out of school). The \cle{Gender Inequality Index} (GII) measures losses in development due to inequality between women and men in reproductive health, empowerment and the labour market.

\begin{exemplebox}[Reading development indicators]
Country A has GDP per capita of \$4{,}000 and an HDI of 0.82; Country B has GDP per capita of \$6{,}000 and an HDI of 0.61. Although B is ``richer'' on average, A delivers longer lives, better schooling and more widely shared prosperity. A is the more \emph{developed} country. This shows why examiners reward candidates who use several indicators together.
\end{exemplebox}

\begin{popfigure}
\begin{center}
\begin{tikzpicture}
\begin{axis}[
    ybar, bar width=0.55cm,
    width=11cm, height=6.5cm,
    ymin=0, ymax=40,
    ylabel={GDP per capita (thousand USD, PPP, approx. 2022)},
    symbolic x coords={Cameroon, Nigeria, SouthKorea, Norway},
    xtick=data,
    xticklabels={Cameroon, Nigeria, South Korea, Norway},
    nodes near coords, 
    every node near coord/.append style={font=\small, anchor=south, yshift=2pt},
    ymajorgrids, grid style={popline!50},
    legend style={at={(0.5,-0.22)}, anchor=north, legend columns=1},
]
\addplot[fill=popblue] coordinates {
    (Cameroon,4) 
    (Nigeria,6) 
    (SouthKorea,33) 
    (Norway,32)
};
\addlegendentry{GDP per capita (\$ '000{,} PPP)}
\end{axis}
% Add HDI annotations above bars
\node[font=\small, text=popdark] at (1.3, 4.5) {HDI: 0.59};
\node[font=\small, text=popdark] at (3.3, 6.5) {HDI: 0.54};
\node[font=\small, text=popdark] at (5.3, 33.5) {HDI: 0.93};
\node[font=\small, text=popdark] at (7.3, 32.5) {HDI: 0.96};
\end{tikzpicture}
\end{center}
\legende{GDP per capita (PPP, thousands USD, approx. 2022) and HDI values (2022) for Cameroon, Nigeria, South Korea and Norway. South Korea and Norway combine high income with very high human development; Cameroon and Nigeria lag on both, illustrating that resource income (Nigeria's oil) does not automatically translate into high HDI.}
\end{popfigure}

\courssub{Characteristics of Developing Countries}

Although developing countries differ greatly, most share a recognisable cluster of features:

\begin{itemize}
\item \textbf{Low per capita income} and widespread absolute poverty.
\item \textbf{High population growth}, with a young age structure and heavy dependency ratios (many children per working adult).
\item \textbf{Dependence on primary exports:} a few unprocessed commodities — cocoa, coffee, cotton, crude oil, timber for Cameroon — whose world prices fluctuate sharply.
\item \textbf{Dualism:} a small modern, monetised sector (banks, factories, ministries in Yaoundé and Douala) coexisting with a large traditional subsistence sector.
\item \textbf{Informal sector dominance:} a large share of jobs in unregistered, unprotected activities — roadside trade, motorcycle taxis (\emph{bendskins}), small tailoring.
\item \textbf{Capital shortage and low productivity:} little machinery, poor infrastructure, and low output per worker, especially in agriculture.
\end{itemize}

\courssub{Obstacles to Development}

\textbf{The vicious circle of poverty.} The most famous explanation of persistent underdevelopment is the circular trap: poor countries are poor because they are poor.

\begin{popfigure}
\begin{center}
\begin{tikzpicture}[
  node distance=1.1cm,
  box/.style={draw=popink, fill=popblueL, rounded corners, align=center, text width=3.1cm, minimum height=1.1cm, font=\small},
  arr/.style={-{Stealth[length=3mm]}, very thick, poporange}
]
\node[box] (income) {Low income per head};
\node[box, right=2.6cm of income] (saving) {Low saving};
\node[box, below=1.6cm of saving] (invest) {Low investment in capital};
\node[box, left=2.6cm of invest] (prod) {Low productivity};
\draw[arr] (income) -- (saving);
\draw[arr] (saving) -- (invest);
\draw[arr] (invest) -- (prod);
\draw[arr] (prod) -- (income);
\end{tikzpicture}
\end{center}
\legende{The vicious circle of poverty: low income $\rightarrow$ low saving $\rightarrow$ low investment $\rightarrow$ low productivity $\rightarrow$ low income. Breaking the circle requires an injection from outside the loop — foreign aid, foreign direct investment, better governance or export earnings.}
\end{popfigure}

Beyond this circle, economists identify further obstacles:

\begin{itemize}
\item \textbf{The savings gap:} domestic savings are too small to finance the investment needed for a target growth rate.
\item \textbf{The foreign exchange gap:} export earnings are insufficient to pay for the imported machines, fuel and inputs that investment requires (the ``two-gap'' model).
\item \textbf{The debt burden:} servicing external debt absorbs scarce foreign exchange and budget resources that could fund schools and roads — a severe problem for heavily indebted poor countries.
\item \textbf{Poor governance:} corruption, weak institutions, political instability and misdirected public spending waste resources and discourage investment.
\end{itemize}

\begin{methode}[Arguing that ``growth is necessary but not sufficient for development'']
\begin{enumerate}
\item \textbf{Define both terms} precisely (quantitative vs qualitative change) — this earns definition marks immediately.
\item \textbf{Show necessity:} without rising output, there are no resources to finance schools, hospitals and poverty reduction; sustained development has never occurred without growth.
\item \textbf{Show insufficiency:} give counter-examples of growth without development (oil enclaves, rising inequality, jobless growth); explain the transmission channels that can fail — distribution, public services, employment creation.
\item \textbf{Use indicators:} contrast GDP per capita with HDI or MPI to prove the two can diverge.
\item \textbf{Conclude with a judgement:} growth is the engine, but institutions, equity and human capital decide whether the engine moves the whole society forward.
\end{enumerate}
\end{methode}

\courssub{Cameroon's Development Profile}

Cameroon is classified by the UNDP as a country of \emph{medium} human development (HDI $\approx$ 0.587 in the 2022 report), ranking in the middle of the world distribution — below the North African countries and far below the very-high-HDI group. Its economy is often described as ``Africa in miniature'' because of its diversity:

\begin{itemize}
\item \textbf{Structure of output:} services contribute the largest share of GDP (approx. 50\%), followed by industry (including oil refining and food processing, approx. 25\%) and then agriculture (approx. 15\%).
\item \textbf{Structure of employment:} the pattern is reversed — agriculture still employs the largest share of the labour force (approx. 40--45\%), mostly at low productivity, which is the classic signature of incomplete structural transformation.
\item \textbf{Exports} remain concentrated in primary products: crude oil, cocoa, coffee, cotton, timber and bananas, exposing the country to world price shocks.
\end{itemize}

This gap between what Cameroon \emph{produces} and where Cameroonians \emph{work} is precisely why the government promotes agribusiness and industrialisation — themes we develop later in this chapter.

\begin{exoresolu}[Interpreting development data]
\textbf{Statement.} In a given year, Country C records real GDP growth of 6\%, yet its HDI is unchanged, the share of the population in extreme poverty rises slightly, and 80\% of new output comes from a foreign-owned mining enclave. Has Country C developed? Justify your answer.
\tcblower
\textbf{Solution.} Country C has experienced \emph{economic growth} (real output rose by 6\%) but not \emph{economic development}. Three arguments: (1) development is multidimensional — the unchanged HDI shows no progress in health, education or broadly shared living standards; (2) rising extreme poverty shows the benefits were not distributed to the population, consistent with an enclave sector whose revenues accrue to foreign owners and a small elite; (3) there is no structural transformation, since the traditional sector and employment patterns are untouched. Conclusion: growth is necessary but not sufficient; without inclusive distribution, human capital formation and structural change, higher GDP does not mean development.
\end{exoresolu}

\begin{exercice}[Applying the concepts]
\begin{enumerate}
\item Distinguish carefully between economic growth and economic development, giving one Cameroonian example of each.
\item Explain two limitations of GDP per capita as an indicator of development and show how the HDI attempts to overcome them.
\item Using a diagram, explain the vicious circle of poverty and suggest two policies by which Cameroon could break out of it.
\end{enumerate}
\end{exercice}

\begin{corrige}[Exercise 1]
\textbf{1.} Growth is a sustained rise in real output (e.g., an increase in Cameroon's GDP driven by higher cocoa and oil output); development adds qualitative change — better health, education, equity and institutions (e.g., rising school enrolment and falling infant mortality alongside income growth). \textbf{2.} GDP per capita ignores distribution and non-income dimensions (and the informal sector); the HDI adds life expectancy and schooling, capturing health and education directly, though it still averages out inequality. \textbf{3.} The circle runs: low income $\rightarrow$ low saving $\rightarrow$ low investment $\rightarrow$ low productivity $\rightarrow$ low income. Escape routes include foreign aid and FDI to fill the savings and foreign exchange gaps, and improved governance plus investment in education and rural infrastructure to raise productivity directly.
\end{corrige}

\begin{attention}[Exam extra — growth without development]
A favourite GCE essay question asks whether high GDP growth proves development. The strongest answers use the case of \textbf{oil-exporting economies}: GDP per capita can look respectable while the majority remains poor, because oil is capital-intensive (creates few jobs), revenues are easily captured by elites, and the rest of the economy — agriculture and manufacturing — is neglected (the ``Dutch disease''). Always conclude that development requires \emph{inclusive} growth plus deliberate investment in human capital and diversification.
\end{attention}

\begin{aretenir}[Key points of Section 1]
\begin{itemize}
\item Development is \textbf{multidimensional}: income, health, education, equity and structural change; growth is only the quantitative part.
\item GDP per capita is useful but limited; the \textbf{HDI} combines life expectancy, schooling and GNI per capita; the MPI and GII refine the picture.
\item Developing countries share low incomes, high population growth, primary-export dependence, dualism, large informal sectors and capital shortage.
\item Key obstacles: the \textbf{vicious circle of poverty}, the savings and foreign exchange gaps, debt and poor governance.
\item Cameroon is a medium-HDI country whose employment remains dominated by low-productivity agriculture — the mark of incomplete transformation.
\end{itemize}
\end{aretenir}

\coursec{Section 2: Industrialisation, Deindustrialisation and Reindustrialisation}

In Section 1, we established that economic development is a structural transformation of society, not merely an increase in GDP per head. But \emph{what actually drives} this transformation? If you observe the history of today's rich nations — Britain, Germany, Japan, South Korea — one pattern repeats itself: each of them moved labour and capital out of low-productivity agriculture into \cle{manufacturing}, and living standards rose as factories multiplied. No country has become rich by remaining a pure exporter of raw cocoa beans, crude oil or unprocessed timber. This observation leads us naturally to the central theme of this section: \cle{industrialisation}, its possible reversal (\cle{deindustrialisation}), and the modern attempts to rebuild industry (\cle{reindustrialisation}).

\courssub{Definition and Rationale of Industrialisation}

\begin{definition}[Industrialisation]
\cle{Industrialisation} is the process by which an economy transforms itself from a primary-based structure (agriculture, fishing, mining of raw materials) towards one dominated by \textbf{secondary activities} — manufacturing, construction, energy and public works — characterised by the use of machinery, the division of labour, mass production and rising productivity. It is usually measured by the \textbf{share of industry (or manufacturing) in GDP} and in \textbf{total employment}.
\end{definition}

Why do economists call manufacturing an \emph{engine of growth}? Three rigorous arguments explain this.

\textbf{(1) Linkages.} A factory does not stand alone. It creates \cle{backward linkages} (demand for inputs: a brewery demands barley, bottles, labels, transport) and \cle{forward linkages} (its output feeds other industries: flour feeds bakeries; cement feeds construction). One industrial investment therefore pulls a whole chain of activities behind it.

\textbf{(2) Productivity and increasing returns.} Manufacturing benefits from \emph{economies of scale} and rapid technological progress: output per worker rises much faster in a factory than on a traditional farm. As labour moves from low-productivity agriculture to high-productivity industry, national income per head rises mechanically.

\textbf{(3) Exports and foreign exchange.} Manufactured goods generally command more stable and higher prices than raw commodities, whose world prices fluctuate violently (think of cocoa or crude oil). Exporting manufactures earns the foreign exchange needed to import machinery and technology.

\begin{exemplebox}[From cocoa beans to chocolate — illustrative example]
Cameroon exports raw cocoa beans and re-imports chocolate at many times the price. If 1 kg of beans sells for, say, 1 200 FCFA while the equivalent processed chocolate sells abroad for 8 000 FCFA, the \emph{value added} — the difference — is captured by the foreign processor. Industrialisation means capturing that value added at home: grinding cocoa into liquor, butter and powder locally, as some plants in Douala now attempt to do. (Note: the prices above are illustrative orders of magnitude, not current market quotations.)
\end{exemplebox}

\begin{popfigure}
\begin{center}
\begin{tikzpicture}
\begin{axis}[
  width=12cm, height=7.5cm,
  xlabel={Level of development (GDP per head, log scale)},
  ylabel={Share of manufacturing in GDP (\%)},
  xmin=0, xmax=10, ymin=0, ymax=40,
  xtick=\empty, ytick={10,20,30},
  axis lines=left,
  every axis plot/.append style={very thick},
]
\addplot[popblue, smooth, domain=0.3:9.7, samples=100]
  {30*exp(-((x-5.5)^2)/8)};
\node[popdark, align=center, font=\small] at (axis cs:1.2,8) {Pre-industrial\\(agriculture)};
\node[popdark, align=center, font=\small] at (axis cs:5.5,33) {Industrial\\take-off};
\node[popdark, align=center, font=\small] at (axis cs:9,8) {Post-industrial\\(services)};
\end{axis}
\end{tikzpicture}
\end{center}
\legende{The inverted-U pattern: the share of manufacturing in GDP rises during industrialisation, peaks, then falls as services dominate (deindustrialisation phase).}
\end{popfigure}

\begin{attention}[Industrialisation is not a free lunch]
A very common exam mistake is to present industrialisation as \emph{only} beneficial. A balanced essay must also discuss its \textbf{costs}: industrial pollution (air, water — think of factory effluents around Douala's industrial zone), urban congestion and slums as rural migrants flood the cities, the destruction of traditional crafts, and a possible new \textbf{dependency} on foreign capital, technology and imported inputs. In the GCE, evaluation marks are earned precisely by this two-sided discussion.
\end{attention}

\courssub{Strategies of Industrialisation: ISI versus EOI}

Once a country decides to industrialise, it must choose \emph{where the demand for its manufactures will come from}. Two grand strategies have dominated development policy.

\begin{definition}[The two strategies]
\begin{itemize}
\item \cle{Import-Substitution Industrialisation (ISI)}: the country produces at home the manufactured goods it previously imported, protecting its infant industries behind \textbf{tariffs and quotas}. Demand comes from the \emph{domestic} market.
\item \cle{Export-Oriented Industrialisation (EOI)}: the country produces manufactures for the \emph{world} market, competing on price and quality, often with state support (subsidies, cheap credit, competitive exchange rate) but with openness to foreign technology and competition.
\end{itemize}
\end{definition}

\begin{popfigure}
\begin{center}
\begin{tikzpicture}[
  box/.style={draw=popblue, thick, rounded corners, fill=popblueL, text width=4.6cm, align=center, font=\small, minimum height=1.1cm},
  ebox/.style={draw=popgreen, thick, rounded corners, fill=popgreenL, text width=4.6cm, align=center, font=\small, minimum height=1.1cm},
  arr/.style={->, thick, popdark}
]
\node[box, align=center] (isi) at (0,4.5){\textbf{ISI}\\Protect domestic market};
\node[box] (isi1) at (0,2.8) {Tariffs \& quotas on imports};
\node[box] (isi2) at (0,1.1) {Local firms replace imports};
\node[box, fill=popblueL] (isi3) at (0,-0.6) {Risk: inefficiency, small market, rent-seeking};
\node[ebox, align=center] (eoi) at (8,4.5){\textbf{EOI}\\Target world market};
\node[ebox] (eoi1) at (8,2.8) {Competitive exchange rate, subsidies};
\node[ebox] (eoi2) at (8,1.1) {Firms export \& face competition};
\node[ebox] (eoi3) at (8,-0.6) {Risk: dependence on world demand};
\draw[arr] (isi) -- (isi1); \draw[arr] (isi1) -- (isi2); \draw[arr] (isi2) -- (isi3);
\draw[arr] (eoi) -- (eoi1); \draw[arr] (eoi1) -- (eoi2); \draw[arr] (eoi2) -- (eoi3);
\end{tikzpicture}
\end{center}
\legende{The logical chains of the two industrialisation strategies.}
\end{popfigure}

\begin{aretenir}[ISI vs EOI — features, merits and drawbacks]
\begin{center}
\begin{tabularx}{\linewidth}{|l|X|X|}
\hline
\rowcolor{popblueL} \textbf{Criterion} & \textbf{ISI} & \textbf{EOI} \\
\hline
Market targeted & Domestic & Foreign (world) \\
\hline
Trade policy & Protectionist (tariffs, quotas) & Open, export incentives \\
\hline
Merits & Builds infant industries; saves foreign exchange; creates urban jobs quickly & Forces efficiency and quality; earns foreign exchange; market size unlimited; technology transfer \\
\hline
Drawbacks & Firms stay inefficient behind protection; domestic market too small and poor; corruption and rent-seeking; goods too dear for consumers & Vulnerable to world recessions and protectionism abroad; requires strong state capacity and infrastructure; harsh competition \\
\hline
Typical examples & Latin America, most of Africa (incl. Cameroon) after independence & South Korea, Taiwan, Singapore, later China and Vietnam \\
\hline
\end{tabularx}
\end{center}
\end{aretenir}

The historical verdict is striking: the East Asian \emph{Tigers} that chose EOI industrialised spectacularly, while most ISI economies (including Cameroon) ran into crisis by the 1980s. Yet the lesson is nuanced: the Tigers \emph{first} protected infant industries (an ISI phase) before pushing them to export. Sequencing matters.

\courssub{Preconditions for Industrialisation}

Why can a country not simply decree factories into existence? Industrialisation requires a set of preconditions:

\begin{aretenir}[Key preconditions for successful industrialisation]
\begin{itemize}
\item \textbf{Capital}: savings and investment to finance plants and machinery; domestic savings or foreign direct investment (FDI).
\item \textbf{Infrastructure}: reliable roads, ports (Douala, Kribi), railways, water and telecommunications — without them, transport costs kill competitiveness.
\item \textbf{Human capital}: literate, trained and healthy workers, engineers, technicians and managers.
\item \textbf{Energy supply}: cheap and reliable electricity; Cameroon's recurrent power cuts (\emph{délestage}) are a major brake on its industries.
\item \textbf{Stable institutions}: rule of law, secure property rights, low corruption, predictable taxation — investors fear arbitrary environments.
\item \textbf{Entrepreneurship and a supportive state}: credit access, industrial zones, and coherent policy.
\end{itemize}
\end{aretenir}

\begin{exoresolu}[Analysing an industrialisation constraint]
\textbf{Statement.} A Cameroonian SME producing fruit juice pays for electricity generator fuel for 60 days a year because of power cuts. Its unit production cost is 950 FCFA per carton, of which 150 FCFA is generator-related. A Ghanaian competitor produces the same carton at 800 FCFA. (a) Compute the cost reduction if reliable grid electricity eliminated generator costs. (b) Explain the consequence for the competitiveness of Cameroonian industry under AfCFTA free trade. (c) Conclude on one precondition of industrialisation.
\tcblower
\textbf{Solution.}
\begin{itemize}
\item (a) New unit cost $= 950 - 150 = 800$ FCFA, a reduction of $\dfrac{150}{950}\times 100 \approx 15.8\%$.
\item (b) With unreliable power, the Cameroonian firm sells at a cost 150 FCFA (about 18.8\% of the Ghanaian cost) \emph{above} the Ghanaian rival; under AfCFTA tariff-free competition, it would lose market share or close. With reliable electricity, costs equalise at 800 FCFA and the firm can compete on equal footing.
\item (c) Conclusion: \textbf{energy infrastructure} is a binding precondition of industrialisation. Investment in generation and distribution (e.g. Memve'ele, Nachtigal dams) is therefore not a luxury but the foundation of any industrial strategy.
\end{itemize}
\end{exoresolu}

\courssub{Deindustrialisation: Definition, Causes and Consequences}

\begin{definition}[Deindustrialisation]
\cle{Deindustrialisation} is a sustained \textbf{decline in the share of manufacturing} in GDP and/or in total employment. It may be \emph{relative} (manufacturing grows more slowly than services) or \emph{absolute} (factory closures, falling industrial output and jobs).
\end{definition}

Deindustrialisation has several distinct causes:

\begin{itemize}
\item \textbf{Mature (normal) deindustrialisation}: in rich countries, as incomes rise, consumers spend more on services; productivity gains mean fewer workers produce the same goods. This is the downward slope of the inverted-U curve above.
\item \textbf{Global competition}: cheap manufactures from Asia have destroyed textile, shoe and assembly industries in both rich and poor countries (e.g. CICAM's struggles in Cameroon facing imported fabrics).
\item \textbf{Premature deindustrialisation}: developing countries begin to \emph{lose} industry at far lower income levels than the rich countries did — they deindustrialise before ever becoming industrialised.
\item \textbf{Dutch disease}: a booming natural-resource sector indirectly strangles manufacturing.
\end{itemize}

\begin{definition}[Dutch disease]
\cle{Dutch disease} is the harmful effect a natural-resource boom can have on the rest of the economy: export earnings from the resource (oil, gas, minerals) push up the \textbf{real exchange rate} and drain labour and capital towards the booming sector, making manufactured exports \textbf{uncompetitive}. The name refers to the Netherlands' experience after North Sea gas discoveries in the 1960s--70s. In Cameroon and CEMAC, oil revenues have similarly crowded out non-oil industry and agriculture, leaving the economy exposed when oil prices collapse.
\end{definition}

\textbf{Consequences of deindustrialisation} include mass \textbf{job losses} in industrial towns, \textbf{regional decline} (whole cities built around one factory collapse when it closes), loss of technological know-how, rising trade deficits in manufactures, and social unrest as well-paid industrial jobs are replaced by insecure informal work.

\courssub{Reindustrialisation: Meaning and Policies}

\begin{definition}[Reindustrialisation]
\cle{Reindustrialisation} is the deliberate policy of \textbf{rebuilding the industrial base} of an economy that has deindustrialised — raising again the share of manufacturing in output and employment, usually through modern, higher-technology industries rather than a simple return to old smokestack sectors.
\end{definition}

Typical reindustrialisation policies include:
\begin{itemize}
\item \textbf{Industrial policy}: targeted state support to strategic sectors (agro-processing, pharmaceuticals, digital manufacturing), industrial zones and special economic zones.
\item \textbf{Innovation}: support for research and development (R\&D), links between universities and firms, adoption of automation and digital technologies.
\item \textbf{Skills}: vocational and technical training to supply engineers and technicians.
\item \textbf{Reshoring}: incentives for firms to bring production back home (a strong trend in the USA and Europe after supply-chain shocks), and \emph{local content} rules requiring foreign firms to source inputs locally.
\item \textbf{Competitiveness measures}: reliable energy, cheaper credit, reduced import dependence for inputs.
\end{itemize}

\courssub{Industrialisation in Cameroon: A Historical Sketch}

\begin{itemize}
\item \textbf{After independence (1960s--1980s)}: Cameroon pursued classic \textbf{ISI}. The state created or protected import-replacing industries — textiles (CICAM), sugar (SOSUCAM), aluminium smelting (ALUCAM at Edéa), oil refining (SONARA at Limbe), cement (CIMENCAM). Protected from imports, these firms supplied the domestic market.
\item \textbf{The crisis of the 1980s--90s}: the collapse of world prices of oil, cocoa and coffee exposed the fragility of protected, inefficient industries. Under structural adjustment programmes (SAPs), the state cut subsidies, privatised or liquidated many public enterprises, and liberalised trade. Several industries closed; Cameroon experienced a clear episode of \textbf{premature deindustrialisation}, and the informal sector absorbed the displaced workers.
\item \textbf{Current ambitions}: the \textbf{Vision 2035} (``emergence by 2035'') and the \textbf{National Development Strategy 2020--2030 (SND30)} place industrialisation at the centre of policy, through \emph{import-substitution} revived in a modern form (local processing of cocoa, cotton, timber; ``Made in Cameroon'' promotion), agro-industrial poles, the Kribi deep-sea port and industrial zone, and major energy projects. The strategy today blends ISI (substituting imports of rice, fish, pharmaceuticals) with EOI ambitions (processed cocoa, aluminium, regional exports within CEMAC and AfCFTA).
\end{itemize}

\begin{savaistu}[The premature deindustrialisation debate — exam extra]
Economist Dani Rodrik popularised the idea that Africa is \emph{deindustrialising before industrialising}: manufacturing's share of GDP in many African countries peaked at income levels far below those at which Britain or Korea peaked. The causes debated include Chinese competition, Dutch disease from oil and minerals, weak infrastructure and premature trade liberalisation under SAPs. The stakes are high: if manufacturing can no longer absorb masses of young workers, can services (telecoms, digital) play that role instead? In an essay, citing this debate — and nuancing it with Cameroon's agro-industrial ambitions — is exactly the kind of \emph{evaluation} that earns top marks.
\end{savaistu}

\begin{methode}[Writing a balanced essay on industrialisation strategies]
\begin{enumerate}
\item \textbf{Define} industrialisation and link it to structural transformation (Section 1).
\item \textbf{Explain} the engine-of-growth arguments: linkages, productivity, export earnings.
\item \textbf{Contrast} ISI and EOI with a clear table (market, trade policy, merits, drawbacks, examples).
\item \textbf{Apply} to Cameroon: post-independence ISI, SAP-era deindustrialisation, current SND30 blend.
\item \textbf{Evaluate} critically: costs (pollution, congestion, dependency), missing preconditions (energy, skills, institutions), Dutch disease from oil, premature deindustrialisation.
\item \textbf{Conclude} with a nuanced judgement: industrialisation is necessary but must be sequenced, green, inclusive, and supported by the preconditions listed above.
\end{enumerate}
\end{methode}

\begin{exercice}[Essay practice]
``Industrialisation remains the only reliable road to economic emergence for Cameroon.'' Discuss, with reference to ISI and EOI strategies, the preconditions of industrialisation, and the risks of Dutch disease and premature deindustrialisation. (20 marks, GCE style.)
\end{exercice}

\begin{corrige}[Essay practice — guidance]
A strong answer (i) defines industrialisation and links it to development (Section 1); (ii) explains the engine-of-growth arguments (linkages, productivity, exports); (iii) contrasts ISI and EOI with Cameroon's post-independence ISI and the Asian Tigers' EOI; (iv) evaluates: costs (pollution, congestion, dependency), preconditions still lacking (energy, skills, institutions), Dutch disease from oil, premature deindustrialisation after the SAPs; (v) concludes with a nuanced judgement — industrialisation is necessary but must be sequenced, green and inclusive, as SND30 attempts. Balance and Cameroonian examples are what distinguish an A-grade script.
\end{corrige}

\coursec{Section 3: Development Strategies of Newly Industrialised Countries (NICs) — Further Studies 9}

In the previous section we studied industrialisation, deindustrialisation and reindustrialisation, and we saw that industrialisation is widely regarded as the royal road to development. But this raises a burning question: if industrialisation is the key, why have some countries succeeded spectacularly at it while others — including most of Africa — have struggled? The answer lies in the experience of a remarkable group of economies that moved from poverty to industrial power within a single generation: the \cle{Newly Industrialised Countries} (NICs). In 1960, South Korea had a real income per head comparable to that of Cameroon; today it is a high-income economy exporting cars, ships and smartphones to the whole world. Understanding \emph{how} this was done — and what of it can be copied — is the object of this section.

\courssub{Definition and Identification of NICs}

\begin{definition}[Newly Industrialised Country (NIC)]
A \cle{Newly Industrialised Country} is a country whose level of economic development and industrialisation places it between the developing countries and the advanced industrialised countries: it has experienced rapid, sustained industrial growth, a rising share of manufacturing in output and exports, and a fast increase in real income per head, without yet having reached the living standards of the richest nations.
\end{definition}

Economists usually distinguish two generations of NICs:

\begin{itemize}
\item The \textbf{first generation}, the famous \cle{Asian Tigers} (or ``Four Dragons''): \textbf{South Korea, Taiwan, Singapore and Hong Kong}. From the 1960s onwards they recorded growth rates of real GDP often exceeding $7\%$ per year for three decades — a performance without historical precedent.
\item The \textbf{second generation} of NICs, which industrialised later, from the 1970s--1980s: \textbf{Malaysia, Thailand, Indonesia}, and on a much larger scale \textbf{China}, together with Latin American cases such as \textbf{Brazil} and \textbf{Mexico}.
\end{itemize}

\begin{savaistu}[Did you know?]
At independence in 1960, South Korea and Cameroon had broadly similar levels of income per head, and Korea was emerging from a devastating war (1950--1953). Many observers at the time considered that African countries, richly endowed with natural resources, had \emph{better} development prospects than resource-poor Korea. Sixty years later, Korea's income per head is roughly ten to twenty times that of Cameroon. Development is not a matter of natural endowment alone — it is a matter of \textbf{strategy}.
\end{savaistu}

\begin{popfigure}
\begin{center}
\begin{tikzpicture}
\begin{axis}[
width=13cm, height=7.5cm,
xlabel={Year}, ylabel={Real GDP per capita (international dollars, index)},
xmin=1960, xmax=2020, ymin=0, ymax=30,
xtick={1960,1970,1980,1990,2000,2010,2020},
ytick={0,5,10,15,20,25,30},
legend pos=north west,
grid=major, grid style={popline!40},
axis line style={popdark},
]
\addplot[color=popblue, very thick, smooth] coordinates {
(1960,1.0) (1965,1.2) (1970,1.8) (1975,2.6) (1980,3.6)
(1985,5.4) (1990,8.2) (1995,12.0) (2000,15.0) (2005,19.0)
(2010,23.5) (2015,26.5) (2020,28.5)};
\addlegendentry{South Korea (illustrative index{,} 1960 = 1)}
\addplot[color=poporange, very thick, smooth] coordinates {
(1960,1.0) (1965,1.1) (1970,1.3) (1975,1.6) (1980,2.0)
(1985,2.1) (1990,1.9) (1995,1.7) (2000,1.8) (2005,1.9)
(2010,2.1) (2015,2.3) (2020,2.2)};
\addlegendentry{Cameroon (illustrative index{,} 1960 = 1)}
\end{axis}
\end{tikzpicture}
\end{center}
\legende{The great divergence: illustrative index of real GDP per capita, South Korea versus Cameroon, 1960--2020 (1960 = 1). Starting from similar levels, the two economies followed radically different paths.}
\end{popfigure}

The figure above illustrates what economists call the \cle{great divergence}: two economies starting from a comparable point, one taking off through industrialisation, the other stagnating around a low-income equilibrium. The rest of this section explains the ingredients of the take-off.

\courssub{The Common Features of NIC Success}

Although each country had its own history, the NICs shared a striking number of common features. Examiners expect you to know at least four of them.

\textbf{1. Export-led growth.} Rather than producing only for their small domestic markets (the \emph{import-substitution} strategy tried in much of Africa and Latin America), the Tigers oriented their industries towards \textbf{world markets}. Exports of textiles, then steel, ships, electronics and cars provided the demand that allowed factories to grow, to exploit \emph{economies of scale} and to earn foreign currency. Export targets also forced firms to meet international standards of quality and price — a powerful discipline.

\textbf{2. High savings and investment rates.} Industrialisation requires massive investment in machines, infrastructure and skills. The NICs financed it largely from \textbf{domestic savings}: gross savings and investment rates frequently reached $30\%$ to $40\%$ of GDP, against often less than $20\%$ in many African economies. High savings were encouraged by stable banks, positive real interest rates and, in some cases, compulsory saving schemes (Singapore's Central Provident Fund).

\textbf{3. Investment in education.} From the start, the Tigers devoted a large share of public spending to primary, secondary and technical education. A literate, disciplined and trainable workforce allowed firms to absorb foreign technology quickly and to climb the ladder of sophistication.

\textbf{4. Macroeconomic stability.} The NICs generally kept inflation moderate, public finances under control and exchange rates competitive (often deliberately undervalued to support exporters). Stability lowered uncertainty and encouraged long-term investment.

\begin{popfigure}
\begin{center}
\begin{center}\tcbox[colback=popblue,colframe=popblue,coltext=white,arc=9pt,boxsep=4pt,left=6pt,right=6pt]{\bfseries\small Pillars of NIC success}\end{center}
\vspace{-2pt}
\begin{tcbraster}[raster columns=3,raster equal height=rows,raster column skip=6pt,raster row skip=6pt,enhanced,blankest]
\begin{tcolorbox}[enhanced,colback=poporange,colframe=poporange,coltext=white,boxrule=0.9pt,arc=3pt,left=4pt,right=4pt,top=3pt,bottom=3pt,boxsep=1pt,halign=left]\bfseries\scriptsize Export-led growth\end{tcolorbox}
\begin{tcolorbox}[enhanced,colback=popgreen,colframe=popgreen,coltext=white,boxrule=0.9pt,arc=3pt,left=4pt,right=4pt,top=3pt,bottom=3pt,boxsep=1pt,halign=left]\bfseries\scriptsize High savings \& investment\end{tcolorbox}
\begin{tcolorbox}[enhanced,colback=poppurple,colframe=poppurple,coltext=white,boxrule=0.9pt,arc=3pt,left=4pt,right=4pt,top=3pt,bottom=3pt,boxsep=1pt,halign=left]\bfseries\scriptsize Education \& human capital\end{tcolorbox}
\begin{tcolorbox}[enhanced,colback=poppink,colframe=poppink,coltext=white,boxrule=0.9pt,arc=3pt,left=4pt,right=4pt,top=3pt,bottom=3pt,boxsep=1pt,halign=left]\bfseries\scriptsize Active state industrial policy\end{tcolorbox}
\begin{tcolorbox}[enhanced,colback=popgold,colframe=popgold,coltext=white,boxrule=0.9pt,arc=3pt,left=4pt,right=4pt,top=3pt,bottom=3pt,boxsep=1pt,halign=left]\bfseries\scriptsize Macroeconomic stability \& FDI\end{tcolorbox}
\end{tcbraster}

\end{center}
\legende{Mind-map of the five pillars of NIC success. None of them alone is sufficient: it is their combination that produced the ``miracle''.}
\end{popfigure}

\courssub{The Role of the State: Industrial Policy}

\begin{attention}[Common mistake — ``the NICs succeeded thanks to free markets alone'']
A frequent error in essays is to attribute the East Asian miracle purely to laissez-faire. In reality, the \textbf{state played a central role}. Markets were used, but they were also \emph{guided}. Writing that the Tigers simply ``let markets work'' will cost you marks for evaluation.
\end{attention}

In South Korea and Taiwan, governments practised a deliberate \cle{industrial policy}:

\begin{itemize}
\item \textbf{Targeted credit}: state-controlled banks channelled cheap loans to selected industries and firms that met export targets.
\item \textbf{Export promotion}: exporters received tax rebates, subsidised inputs, information about foreign markets and diplomatic support.
\item \textbf{Protection of infant industries}: young industries (steel, shipbuilding, electronics) were temporarily sheltered from foreign competition by tariffs and quotas — but protection was \emph{conditional} on performance and was gradually removed. Firms that failed to become competitive lost their privileges.
\item \textbf{Discipline}: crucially, state support was tied to \emph{results}, especially export performance. This ``carrot and stick'' approach limited the rent-seeking and inefficiency that protectionism often creates elsewhere.
\end{itemize}

\begin{propriete}[The ``flying geese'' pattern of industrial upgrading]
East Asian industrialisation followed a sequence known as the \cle{flying geese} model: a country begins by exporting simple labour-intensive goods (textiles, toys), then, as wages rise and skills accumulate, it upgrades to more capital- and technology-intensive products (steel, ships, electronics, cars, semiconductors), while the simpler industries migrate to the next ``goose'' — poorer neighbouring countries (Malaysia, Thailand, then China and Vietnam). The lead economy (Japan, then the Tigers) pulls the whole flock along. \emph{This property is descriptive; its formal proof is not examinable, but you must be able to explain the mechanism.}
\end{propriete}

\courssub{Human Capital, Technology and Foreign Direct Investment}

\textbf{Human capital.} Education was not a luxury but a productive investment. By the time Korean industry needed engineers and technicians, the education system was already producing them in large numbers. This is a powerful lesson: \emph{physical capital without human capital yields little}.

\textbf{Technology transfer and absorption.} The NICs did not invent most of their technology; they \textbf{imported, imitated and improved} it — through licensing, reverse engineering, joint ventures and sending students and engineers abroad. The key was the domestic capacity to \emph{absorb} foreign knowledge, which again depended on education.

\textbf{Foreign direct investment (FDI).} Multinational firms brought capital, technology and access to world distribution networks. Singapore and Hong Kong relied heavily on FDI; Korea relied more on domestic firms (the \emph{chaebols}) borrowing abroad. In all cases, integration into \textbf{world markets} — not isolation — was the engine.

\courssub{Limits and Criticisms of the NIC Model}

A rigorous candidate must evaluate the model, not just admire it.

\begin{itemize}
\item \textbf{Authoritarianism}: several Tigers industrialised under authoritarian regimes that restricted trade unions and political freedoms. Whether repression was necessary for growth is contested — but the social cost was real.
\item \textbf{Environmental costs}: rapid industrialisation brought severe pollution of air and water, only addressed later, at high cost.
\item \textbf{Vulnerability to external shocks}: openness cuts both ways. The \textbf{1997 Asian financial crisis} showed that economies dependent on foreign capital inflows and world demand can be hit brutally when confidence reverses; currencies collapsed and several countries needed international rescue packages.
\item \textbf{Working conditions}: the early phases relied on long hours and low wages in labour-intensive factories.
\end{itemize}

\begin{exoresolu}[Applying the NIC model to a comparison]
\textbf{Statement.} In 1960, country A (a Tiger) and country B (an African economy) both had GDP per capita of about 100 (index). Over 60 years, A grew at an average of $5.5\%$ per year and B at $1.5\%$ per year. Compute the ratio of their incomes per head after 60 years and comment.
\tcblower
\textbf{Solution.} Using compound growth, income after $n$ years is $Y_n = Y_0(1+g)^n$.
\begin{align*}
Y_A &= 100 \times (1.055)^{60} \approx 100 \times 24.8 = 2480,\\
Y_B &= 100 \times (1.015)^{60} \approx 100 \times 2.44 = 244.
\end{align*}
The ratio is $\dfrac{2480}{244} \approx 10$: a gap of about two percentage points in annual growth, sustained for 60 years, produces an income ratio of roughly \textbf{ten to one}. \emph{Comment:} this is the arithmetic of the great divergence — small differences in growth rates, compounded over a generation, explain why the NICs became rich while others stagnated. It also shows why the growth \emph{rate} (and hence strategy) matters more than the starting point.
\end{exoresolu}

\courssub{Lessons for Cameroon and Africa: What Can and Cannot Be Copied}

\textbf{What can be copied:}
\begin{itemize}
\item The priority given to \textbf{education and technical skills} — the foundation of everything else.
\item The \textbf{export orientation} of agriculture and industry (processed cocoa, timber products, cotton textiles) rather than dependence on a few raw commodities; the African Continental Free Trade Area (AfCFTA) offers a large market for this.
\item \textbf{Macroeconomic discipline} — credible public finances and a competitive environment for investors.
\item A \textbf{strategic, performance-based state}: support to infant industries conditional on results, not permanent protection.
\item High \textbf{domestic savings mobilisation} to finance investment without excessive foreign debt.
\end{itemize}

\textbf{What cannot simply be copied:}
\begin{itemize}
\item The world economy of the 1960s--1980s (cheap access to Western markets, Cold War aid to Korea and Taiwan) no longer exists; today's Africa faces competition from already-established giants like China.
\item Authoritarian methods are neither desirable nor necessary; Botswana and Mauritius show that growth is compatible with democracy.
\item Cameroon must build on its \emph{own} endowments — agriculture, energy, a central position in CEMAC — rather than imitate Korean shipbuilding.
\end{itemize}

\begin{methode}[Structuring an essay: ``To what extent can Cameroon replicate the NIC model?'']
\begin{enumerate}
\item \textbf{Introduction}: define NICs; state the issue; announce a balanced thesis (partial replication).
\item \textbf{Part 1 — What Cameroon can borrow}: education and human capital; export-led and agro-industrial strategy; macroeconomic stability; strategic state support tied to performance. Illustrate with the divergence graph and the compound-growth example.
\item \textbf{Part 2 — The limits}: different historical and geopolitical context; weak institutions and governance problems; infrastructure deficits; competition from established NICs; risk of copying protection without discipline.
\item \textbf{Conclusion}: answer ``to what extent'' explicitly — Cameroon can adopt the \emph{principles} (skills, exports, stability, disciplined state) but must adapt them to its own resources and democratic context.
\end{enumerate}
\end{methode}

\begin{aretenir}[Key points]
\begin{itemize}
\item NICs are economies between developing and advanced status: Asian Tigers (Korea, Taiwan, Singapore, Hong Kong), then Malaysia, Thailand, China, Brazil.
\item Success rested on five pillars: export-led growth, high savings and investment, education, an active but disciplining state, and macroeconomic stability with openness to FDI.
\item The ``flying geese'' pattern describes upgrading from labour-intensive to high-technology industries.
\item The model has limits: authoritarianism, pollution, and vulnerability to shocks (1997 crisis).
\item For Cameroon: copy the principles, not the details — skills, exports, stability and a performance-based state.
\end{itemize}
\end{aretenir}

\begin{exercice}[Practice]
\begin{enumerate}
\item Define a Newly Industrialised Country and give two examples from each generation.
\item Explain why high domestic savings were crucial to the NICs' industrialisation.
\item ``The Asian Tigers prove that the state should stay out of the economy.'' Discuss.
\item Using the compound-growth formula, show how a $4\%$ growth gap sustained over 40 years affects relative incomes. Comment on the lesson for Cameroon.
\end{enumerate}
\end{exercice}

\begin{corrige}[Exercise 3]
\textbf{Guidance.} This statement is a caricature. \emph{Against it}: targeted credit, export promotion, infant-industry protection and massive public education show the state was central in Korea and Taiwan. \emph{Nuance}: the NIC state used market signals (export competition) to discipline firms — it \emph{guided} markets rather than replacing them, and macroeconomic policy was prudent. \emph{Conclusion}: the lesson is not ``less state'' but a \emph{smarter, performance-oriented state} combined with competitive markets.
\end{corrige}

\coursec{Section 4: Promoting Second Generation Agriculture and Agribusiness in Cameroon — Further Studies 10}

In Section 3 we saw how the Newly Industrialised Countries (NICs) of East Asia built their take-off on export-oriented manufacturing. But we also noted a crucial lesson: countries such as South Korea and Taiwan did \emph{not} neglect agriculture. On the contrary, they first raised agricultural productivity, which released labour, food and savings for industry. Cameroon, with its fertile soils, varied climate and large rural population, cannot industrialise by ignoring its farms. This section studies how Cameroon can move from a traditional, subsistence agriculture to a \cle{second generation agriculture} driven by \cle{agribusiness}, as a pillar of its emergence strategy.

\courssub{The Role of Agriculture in Economic Development}

\textbf{Intuition.} Imagine a village near Bafoussam where every family grows just enough maize and cassava to feed itself. No surplus is sold, nobody can leave the farm to work in a workshop, and there is no money to buy manufactured goods. Now imagine the same village with improved seeds, a feeder road and a maize mill: yields double, surpluses are sold in Douala, some young people open a poultry unit, and the mill buys spare parts from local mechanics. Agriculture has financed and fed the beginnings of industry. This is exactly the mechanism analysed by \cle{Johnston and Mellor} (1961), who identified five contributions of agriculture to economic development.

\begin{propriete}[The Five Contributions of Agriculture to Development (Johnston \& Mellor)]
Agriculture contributes to economic development through five channels:
\begin{enumerate}
\item \textbf{Food supply:} rising farm output provides cheap food for the growing urban and industrial labour force, keeping wages and inflation low.
\item \textbf{Labour release:} higher farm productivity frees workers who migrate to industry and services (the Lewis mechanism).
\item \textbf{Foreign exchange earnings:} export crops (cocoa, coffee, cotton, banana) earn the foreign currency needed to import machines and technology.
\item \textbf{Raw materials and capital:} agriculture supplies inputs to agro-industry (cotton for textiles, cocoa for grinding) and, through taxation and rural savings, capital for investment.
\item \textbf{Market for industry:} rising rural incomes create demand for manufactured goods (tools, soap, cement, phones), widening the domestic market.
\end{enumerate}
This property is directly examinable; learn to state and illustrate each channel with a Cameroonian example.
\end{propriete}

\begin{exemplebox}[Agriculture in the Cameroonian Economy]
Agriculture remains the backbone of the Cameroonian economy: it employs the largest share of the labour force (roughly half of workers, mostly in small family farms) and contributes a substantial share of GDP (around 15--17\%) and of export earnings through cocoa, coffee, cotton, bananas, rubber and timber-related products. Cocoa from the Centre and South-West regions, cotton from the North (via SODECOTON) and bananas from the Littoral illustrate the foreign-exchange contribution, while food crops (maize, cassava, plantain, yams) feed the cities of Douala and Yaoundé.
\end{exemplebox}

\begin{popfigure}
\begin{center}
\begin{tikzpicture}
\begin{axis}[
    ybar, bar width=22pt,
    width=0.85\linewidth, height=6.2cm,
    ymin=0, ymax=70,
    ylabel={Share (\%)},
    symbolic x coords={GDP share, Employment share},
    xtick=data,
    nodes near coords, nodes near coords align={vertical},
    every node near coord/.append style={font=\small},
    axis lines=left,
    ymajorgrids=true, grid style={popline!40},
]
\addplot[fill=popblue, draw=popdark] coordinates {(GDP share,16) (Employment share,50)};
\end{axis}
\end{tikzpicture}
\end{center}
\legende{Indicative shares of agriculture in Cameroon's GDP and employment: agriculture employs about half of the labour force but produces only about one-sixth of GDP — a sign of low productivity and the huge potential of modernisation.}
\end{popfigure}

\begin{attention}[Reading the Chart Correctly]
The gap between the employment share (about 50\%) and the GDP share (about 16\%) means that average labour productivity in agriculture is far below that of the rest of the economy. Do \textbf{not} conclude that agriculture is ``useless'': conclude that raising farm productivity is the single biggest lever for raising rural incomes and releasing labour for industry — precisely the Johnston--Mellor argument.
\end{attention}

\courssub{From First to Second Generation Agriculture}

\textbf{Intuition.} A farmer who clears a new plot every year, uses a hoe, plants local seeds and eats most of the harvest is practising what Cameroonian policy documents call \cle{first generation agriculture}. It feeds families but transforms little and sells little. \emph{Second generation agriculture} is different in nature, not just in degree: the farmer is an \emph{entrepreneur} who produces for a market, uses improved inputs, and thinks in terms of profit, processing and contracts.

\begin{definition}[First and Second Generation Agriculture]
\cle{First generation agriculture} designates traditional, largely subsistence and extensive farming: low-input, low-yield, hoe-based cultivation on small plots, with surpluses sold occasionally on local markets. \par
\cle{Second generation agriculture} designates a modern, \textbf{market-oriented} and \textbf{value-adding} agriculture: it uses improved seeds, fertilisers, mechanisation and irrigation; it is organised in \textbf{value chains} linked to processing, storage and marketing; and it is run as a business aiming at profit, quality standards and access to national, regional and export markets.
\end{definition}

\begin{center}
\begin{tabularx}{\linewidth}{|l|X|X|}
\hline
\rowcolor{popblueL} \textbf{Criterion} & \textbf{First generation} & \textbf{Second generation} \\
\hline
Orientation & Subsistence, own consumption & Market, profit, contracts \\
\hline
Techniques & Hoe, local seeds, extensive & Improved inputs, mechanisation, irrigation \\
\hline
Land use & Shifting cultivation, new plots & Intensification on existing plots \\
\hline
Output & Raw produce, little processing & Graded, processed, packaged products \\
\hline
Organisation & Isolated family farms & Cooperatives, value chains, agropoles \\
\hline
Financing & Self-financing, tontines & Bank credit, microfinance, subsidies \\
\hline
\end{tabularx}
\end{center}

\courssub{Agribusiness: Meaning and Scope}

\begin{definition}[Agribusiness]
\cle{Agribusiness} is the set of all business activities connected to agriculture, covering the whole \textbf{value chain}: (i) \textbf{upstream} — supply of inputs (seeds, fertilisers, equipment, veterinary services, finance); (ii) \textbf{production} — farming, livestock and fisheries; (iii) \textbf{downstream} — processing, packaging, storage, transport, marketing, distribution and export. Agribusiness therefore includes not only farmers but also input dealers, processors, transporters, wholesalers, exporters and retailers.
\end{definition}

The key idea is \cle{value addition}: a tonne of raw cocoa beans exported earns far less than the same beans transformed into cocoa paste, butter or chocolate. Each stage of the chain adds value, creates jobs and keeps income inside the country.

\begin{popfigure}
\begin{center}
\begin{tikzpicture}[
  box/.style={draw=popdark, fill=popblueL, rounded corners=2pt, minimum height=1.1cm, text width=2.6cm, align=center, font=\small},
  arr/.style={-{Stealth[length=2.5mm]}, thick, poporange},
  node distance=0.55cm
]
\node[box, align=center] (input){\textbf{Input supply}\\ seeds, fertiliser, credit, equipment};
\node[box, right=of input, align=center] (prod){\textbf{Production}\\ farms, cooperatives};
\node[box, right=of prod, align=center] (proc){\textbf{Processing}\\ drying, milling, grading};
\node[box, right=of proc, align=center] (stor){\textbf{Storage \& transport}\\ warehouses, cold chain};
\node[box, right=of stor, align=center] (mark){\textbf{Marketing \& export}\\ Douala port, CEMAC, AfCFTA};
\draw[arr] (input) -- (prod);
\draw[arr] (prod) -- (proc);
\draw[arr] (proc) -- (stor);
\draw[arr] (stor) -- (mark);
\node[below=0.35cm of proc, font=\small\itshape, text=popmuted, align=center, text width=5.5cm] {Value added rises at each stage};
\end{tikzpicture}
\end{center}
\legende{The cocoa value chain in Cameroon, from input supply to export through the port of Douala. Second generation agriculture means mastering the whole chain, not only the production stage.}
\end{popfigure}

\courssub{Constraints on Cameroonian Agriculture}

Why has Cameroonian agriculture remained largely first generation despite its potential? The main constraints are:

\begin{itemize}
\item \textbf{Land tenure insecurity:} much rural land is held under customary rights without formal titles, discouraging long-term investment and making land unusable as collateral for loans.
\item \textbf{Low mechanisation and input use:} most farms depend on the hoe and manual labour; improved seeds, fertiliser and irrigation are under-used, keeping yields low.
\item \textbf{Poor rural infrastructure:} bad feeder roads raise transport costs, cut off production basins during the rainy season and cause produce to rot before reaching markets; electricity and storage facilities are scarce.
\item \textbf{Limited access to credit:} banks consider smallholders risky; interest rates are high and collateral demands exclude most farmers.
\item \textbf{Climate risk:} dependence on rainfall exposes farmers to droughts and floods, aggravated by climate change, with almost no agricultural insurance.
\item \textbf{Post-harvest losses:} a large share of perishable output (tomatoes, plantain, fish, fruits) is lost between farm and market because of poor storage, processing and transport.
\item \textbf{Ageing farmers and weak extension:} young people desert villages, and agricultural extension services (advice, training) are understaffed.
\end{itemize}

\courssub{Government Strategies for Second Generation Agriculture}

Cameroon's long-term vision, the \cle{SND30} (National Development Strategy 2020--2030), places the modernisation of agriculture and the development of agro-industry at the heart of its structural transformation objective. The main policy tools are:

\begin{itemize}
\item \textbf{Value chain development:} public support targets complete chains (cassava, maize, rice, cocoa, palm oil, poultry, fish) rather than isolated production, combining input distribution, processing units and marketing.
\item \textbf{Youth agripreneurship:} programmes encourage young graduates and school-leavers to create agricultural enterprises, with training, starter kits and subsidised credit, so as to rejuvenate the farm population and create jobs.
\item \textbf{Subsidies and input support:} fertiliser and improved-seed subsidies lower production costs, though they must be well targeted to avoid waste and dependency.
\item \textbf{Extension services and research:} revitalising agricultural advisory services and research institutes (e.g.\ IRAD) spreads improved techniques and varieties.
\item \textbf{Agropoles and incubators:} integrated zones that concentrate infrastructure, processing units and services around production basins (see case study below).
\item \textbf{Rural infrastructure:} rehabilitation of rural roads, rural electrification and storage facilities to cut post-harvest losses and connect farmers to markets.
\end{itemize}

\begin{experience}[Case Study — Agropoles and Agribusiness Incubators in Cameroon (exam-useful)]
An \textbf{agropole} is a geographically concentrated zone where production, processing, storage and marketing of agricultural products are brought together with shared infrastructure (roads, power, water, laboratories). Cameroon's agropole programme, supported by international partners, aims to develop competitive value chains (maize, cassava, livestock, aquaculture) in several regions. \textbf{Agribusiness incubators} complement agropoles: they host young entrepreneurs for one to three years, providing training in farm management, access to land or greenhouses, mentoring and help in obtaining finance. \par
\textbf{Analysis for the exam:} agropoles attack several constraints \emph{simultaneously} — infrastructure, processing, market access and youth employment — which is why they are more promising than isolated measures (e.g.\ fertiliser subsidies alone). Their limits are high cost, dependence on public funding, and the risk of poor maintenance of shared infrastructure.
\end{experience}

\courssub{Opportunities: Value Chains and Regional Markets}

Cameroon enjoys exceptional agro-ecological diversity — often called ``Africa in miniature'' — which opens opportunities in several value chains:

\begin{itemize}
\item \textbf{Cassava:} high demand for gari, flour and starch; processing into industrial starch and high-quality flour can substitute imports.
\item \textbf{Maize:} strategic for food security and for poultry feed; better storage could stabilise prices between harvests.
\item \textbf{Cocoa:} Cameroon is a significant world producer; local grinding and chocolate manufacturing would capture more value than bean exports.
\item \textbf{Palm oil:} strong domestic and regional demand; modern mills raise extraction rates and smallholder incomes.
\item \textbf{Poultry:} rapid urban demand growth; a full chain (feed, hatcheries, veterinary services, cold chain) creates many urban and rural jobs.
\end{itemize}

Beyond the domestic market, \cle{regional integration} widens outlets: the \textbf{CEMAC} market (Cameroon, Chad, Central African Republic, Congo, Equatorial Guinea, Gabon) absorbs Cameroonian food products, and the \textbf{African Continental Free Trade Area (AfCFTA)} offers, in the long run, access to a continental market of over a billion consumers — provided Cameroonian products meet quality and phytosanitary standards.

\begin{savaistu}[Did You Know? — Cameroon's Agro-Ecological Zones]
Cameroon is often described as ``Africa in miniature'' because it contains almost every African climate zone: the humid forest zone (South, Centre, East) for cocoa, coffee, oil palm; the highland zone (West, North-West) for arabica coffee, tea, temperate vegetables; the Sudano-Sahelian zone (North, Far North) for cotton, millet, sorghum, livestock; and the coastal zone (Littoral, South-West) for bananas, rubber, plantains. This diversity allows year-round production of a wide range of crops and is a unique asset for agribusiness development.
\end{savaistu}

\begin{methode}[How to Evaluate an Agricultural Policy]
To assess any agricultural policy (subsidies, an agropole, an extension programme) in an essay, apply five criteria systematically:
\begin{enumerate}
\item \textbf{Output and productivity:} does production and yield per hectare rise?
\item \textbf{Farm incomes:} do producers' real incomes increase, and is the gain fairly shared?
\item \textbf{Employment:} how many jobs are created on farms and along the chain (processing, transport, marketing)?
\item \textbf{Food security:} do consumers get more abundant, cheaper and safer food?
\item \textbf{Sustainability:} are soils, water and forests preserved, and is the policy financially sustainable for the State?
\end{enumerate}
A good exam answer applies each criterion, then gives a balanced judgement with advantages and limits.
\end{methode}

\begin{exoresolu}[Worked Exercise — Evaluating a Fertiliser Subsidy]
\textbf{Statement.} The government of Cameroon introduces a subsidy covering 30\% of the price of fertiliser for maize producers in the West Region. Using the five evaluation criteria, discuss the likely effects of this policy and its limits. \tcblower
\textbf{Solution.}
\begin{itemize}
\item \emph{Output:} cheaper fertiliser raises input use and yields; maize output should increase, ceteris paribus.
\item \emph{Incomes:} farmers' costs fall, so margins rise — but mainly for farmers who can already afford the remaining 70\%; the poorest may be excluded.
\item \emph{Employment:} higher output creates some farm and transport jobs, but fertiliser alone creates few processing jobs.
\item \emph{Food security:} larger maize supply lowers prices in cities, benefiting consumers.
\item \emph{Sustainability:} the subsidy burdens the State budget every year; excessive fertiliser use can degrade soils and pollute water.
\end{itemize}
\textbf{Judgement:} the subsidy is useful as a short-term stimulus, but it is not a second-generation strategy by itself. It must be combined with extension services, rural roads, storage and processing (value addition) and market access — otherwise higher yields simply rot in the villages.
\end{exoresolu}

\begin{attention}[Common Mistake — Modernisation Is Not Mechanisation Only]
Many candidates write that modernising agriculture simply means ``giving tractors to farmers''. This is a serious error. A tractor without feeder roads, credit, spare parts, storage and buyers changes nothing. \textbf{Second generation agriculture = improved inputs \emph{and} mechanisation \emph{and} value addition (processing) \emph{and} market access \emph{and} business organisation.} In your essays, always present modernisation as a \emph{value chain} transformation, not as a single machine.
\end{attention}

\begin{aretenir}[Key Points to Remember]
\begin{itemize}
\item Agriculture contributes to development through five channels (Johnston \& Mellor): food, labour, foreign exchange, raw materials/capital, and a market for industry.
\item First generation agriculture is subsistence and extensive; second generation agriculture is market-oriented, value-adding and organised in chains.
\item Agribusiness covers the entire value chain: inputs, production, processing, storage, marketing and export.
\item Cameroon's constraints: land tenure, low mechanisation, poor infrastructure, limited credit, climate risk, post-harvest losses.
\item SND30 strategies: value chains, youth agripreneurship, subsidies, extension, agropoles and incubators.
\item Evaluate any policy with five criteria: output, incomes, employment, food security, sustainability.
\end{itemize}
\end{aretenir}

\begin{exercice}[Practice Questions]
\begin{enumerate}
\item Define second generation agriculture and distinguish it from first generation agriculture using four criteria.
\item Explain, with Cameroonian examples, any three of Johnston and Mellor's five contributions of agriculture to development.
\item Choose one Cameroonian value chain (cocoa, cassava, maize, palm oil or poultry) and show how value addition at each stage can raise national income and employment.
\item To what extent can agropoles solve the constraints facing Cameroonian agriculture? (Use the five evaluation criteria.)
\end{enumerate}
\end{exercice}

\begin{corrige}[Exercise 1]
Second generation agriculture is a modern, market-oriented and value-adding agriculture organised in value chains and run as a business, whereas first generation agriculture is traditional, subsistence and extensive. Distinctions: (i) orientation — profit and contracts versus own consumption; (ii) techniques — improved inputs, mechanisation and irrigation versus the hoe and local seeds; (iii) output — processed, graded and packaged products versus raw produce; (iv) organisation — cooperatives, financing and market linkages versus isolated self-financed family farms.
\end{corrige}

\begin{corrige}[Exercise 2]
Any three of: (i) \emph{Food supply} — maize, cassava and plantain from the West and Centre regions feed Douala and Yaoundé, moderating urban food prices; (ii) \emph{Labour release} — as farm productivity rises, young workers can move to construction, industry and services in the cities; (iii) \emph{Foreign exchange} — cocoa, coffee, cotton and banana exports earn the currency used to import machinery; (iv) \emph{Raw materials} — cotton supplies SODECOTON's gins and the textile industry, cocoa supplies local grinding plants; (v) \emph{Market for industry} — richer rural households buy cement, soap, phones and tools, stimulating domestic manufacturing.
\end{corrige}

\begin{corrige}[Exercise 3]
Take cocoa: at the \emph{input} stage, improved varieties and fertiliser raise yields; at \emph{production}, cooperatives improve quality and bargaining power; at \emph{processing}, local grinding into paste and butter captures more value than exporting raw beans; at \emph{storage and transport}, better roads and warehouses cut losses; at \emph{marketing and export}, branding and quality certification open premium markets. Each stage adds value, raises GDP, and creates jobs (factory workers, transporters, traders), while higher export earnings improve the trade balance. The same logic applies to cassava (flour, starch), maize (feed), palm oil (refining) and poultry (cold chain).
\end{corrige}

\begin{corrige}[Exercise 4]
\emph{Output:} agropoles concentrate inputs, infrastructure and processing, so yields and marketed volumes should rise. \emph{Incomes:} farmers gain from better prices and lower losses, though benefits depend on fair contracts with processors. \emph{Employment:} jobs are created in farming, processing plants, transport and services, especially for youths. \emph{Food security:} larger, steadier supplies stabilise urban prices. \emph{Sustainability:} high set-up costs and dependence on public funds are risks; maintenance and good governance are essential. \emph{Conclusion:} agropoles address several constraints simultaneously and are therefore a powerful instrument of second generation agriculture, but they must be complemented by land reform, credit access and extension services to be fully effective.
\end{corrige}

\coursec{Section 5: The Trade (Business) Cycle and Government Economic Objectives}

In Section 4 we saw how Cameroon is trying to transform its agriculture into a modern, job-creating agribusiness sector. But even the best development strategy does not produce smooth, uninterrupted progress. Anyone who has followed the Cameroonian economy knows this from experience: the years of oil-fuelled euphoria in the early 1980s were followed by the deep crisis of 1986--1994, then by recovery after the 1994 devaluation of the CFA franc. Economies do not grow in a straight line; they \emph{fluctuate} around their long-run trend. Understanding these fluctuations --- the \cle{trade cycle} --- and the \cle{objectives} that governments pursue in trying to manage them is the purpose of this section.

\courssub{The Trade (Business) Cycle: Meaning and Phases}

\begin{definition}[The trade (business) cycle]
The \cle{trade cycle} (or \cle{business cycle}) is the recurring pattern of expansion and contraction in the level of economic activity --- measured by real national output (real GDP) --- around the long-run trend rate of growth. A complete cycle runs from one peak (or one trough) to the next.
\end{definition}

Two points in this definition deserve emphasis. First, the cycle concerns \emph{real} output, not merely prices or money. Second, fluctuations occur \emph{around a trend}: a growing economy like Cameroon can experience a recession (falling output) even though its long-run trend is upward.

A complete cycle is conventionally divided into \textbf{four phases}.

\begin{definition}[Phases of the trade cycle]
\begin{itemize}
\item \textbf{Boom (peak / prosperity):} the phase in which economic activity reaches its maximum. Output is at or above full capacity, employment is very high, investment and consumer spending are strong, business confidence is high, and demand-pull inflationary pressure builds up.
\item \textbf{Recession (downturn / contraction):} the phase in which real output falls. Firms cut production and investment, unemployment rises, incomes and spending decline, confidence collapses, and inflationary pressure eases. A fall in real GDP for two consecutive quarters is the conventional technical definition of a recession.
\item \textbf{Trough (slump / depression):} the lowest point of the cycle. Output and employment are at their minimum, many firms operate far below capacity or close down, investment is very low, and pessimism is generalised. A prolonged and severe trough is called a \emph{depression}.
\item \textbf{Recovery (upturn / expansion):} the phase in which activity picks up again. Investment resumes, employment and incomes rise, confidence returns, and output climbs back towards and then beyond the trend line.
\end{itemize}
\end{definition}

\begin{popfigure}
\begin{center}
\begin{tikzpicture}
\begin{axis}[
width=13cm, height=7cm,
axis lines=middle,
xlabel={Time}, ylabel={Real GDP},
xtick=\empty, ytick=\empty,
xmin=0, xmax=12.8, ymin=-1.8, ymax=3.4,
every axis x label/.style={at={(ticklabel* cs:1.02)}, anchor=west},
every axis y label/.style={at={(ticklabel* cs:1.02)}, anchor=south},
]
\addplot[domain=0:12.6, samples=200, thick, popblue] {0.18*x + 1.1*sin(deg(1.05*x))};
\addplot[domain=0:12.6, samples=2, dashed, popmuted] {0.18*x};
\node[popmuted, align=center] at (axis cs:11.6,1.35) {\small Trend\\ \small growth};
\node[popdark, align=center] at (axis cs:1.55,2.35) {\small \textbf{Boom}};
\node[popdark, align=center] at (axis cs:4.6,0.15) {\small \textbf{Recession}};
\node[popdark, align=center] at (axis cs:4.75,-1.35) {\small \textbf{Trough}};
\node[popdark, align=center] at (axis cs:8.0,1.9) {\small \textbf{Recovery}};
\node[popdark, align=center] at (axis cs:7.6,3.05) {\small \textbf{Boom}};
\node[popdark, align=center] at (axis cs:10.7,0.55) {\small \textbf{Recession}};
\end{axis}
\end{tikzpicture}
\legende{The trade cycle: fluctuations of real GDP around the long-run trend, with the four phases.}
\end{center}
\end{popfigure}

\begin{attention}[Cycles are not regular]
A very common examination mistake is to describe the trade cycle as a regular, predictable, clock-like movement with phases of fixed length. This is wrong. Real cycles vary enormously in \textbf{duration} (from a few years to a decade or more) and in \textbf{amplitude} (some recessions are mild, others are deep depressions). The curve above is a stylised representation, not a timetable. Always state this explicitly in your essays --- examiners reward it.
\end{attention}

\courssub{Causes of the Trade Cycle}

Why does economic activity fluctuate? Several explanations exist; they are complements rather than rivals.

\textbf{1. The Keynesian explanation: fluctuations in aggregate demand.} For Keynesians, the cycle is driven by swings in \cle{aggregate demand}, especially in \emph{investment}, which is volatile because it depends on business expectations (Keynes's ``animal spirits''). A fall in investment reduces aggregate demand; through the \cle{multiplier}, national income falls by a larger amount; falling income then reduces consumption and induces further cuts in investment. The interaction of the multiplier with the \emph{accelerator} (studied below) can turn a small initial shock into a full cyclical movement.

\textbf{2. Monetary theories.} For monetarists, cycles result from erratic changes in the \cle{money supply} and credit. Excessive credit expansion fuels a boom; when the central bank tightens money to fight inflation, spending contracts and a recession follows. In the CEMAC zone, monetary policy is conducted by the BEAC (Banque des États de l'Afrique Centrale), so the money supply channel operates at the regional level.

\textbf{3. Real business cycle and external shocks.} This school argues that fluctuations originate in \emph{real} (supply-side) shocks: changes in technology, harvests, or the prices of key inputs such as oil. For open, commodity-exporting economies like those of the CEMAC zone, this explanation is particularly relevant.

\begin{savaistu}[The 1986--1994 crisis in Cameroon]
Cameroon offers a textbook illustration of an externally triggered cycle. In the early 1980s, oil revenues and high world prices for cocoa and coffee sustained a boom. From 1986, world prices of oil, cocoa and coffee collapsed. Because the CFA franc was pegged to the French franc, Cameroon could not devalue to restore competitiveness. The result was a long and painful recession: falling output, unpaid salaries, bank failures and rising unemployment, lasting until the devaluation of the CFA franc in January 1994, after which recovery gradually began. A single external shock, amplified by rigid domestic structures, produced nearly a decade of contraction.
\end{savaistu}

\courssub{The Accelerator Principle and the Multiplier--Accelerator Interaction}

Keynesian economists explain the \emph{amplification} and \emph{self-sustaining} character of cycles by combining two mechanisms. The \cle{multiplier} (studied in Lower Sixth) says that a change in investment causes a larger change in national income: $k = \dfrac{\Delta Y}{\Delta I} = \dfrac{1}{1 - MPC}$. The \cle{accelerator} works in the opposite direction.

\begin{definition}[The accelerator principle]
The \cle{accelerator principle} states that the level of \emph{net investment} depends on the \emph{rate of change} of national income (or output), not on its level:
\[ I_t = v\,\Delta Y = v\,(Y_t - Y_{t-1}) \]
where $v$ is the \cle{capital--output ratio} (accelerator coefficient), i.e. the value of capital needed to produce one extra unit of output per period.
\end{definition}

The intuition is simple. If a firm needs \SI{3}{} FCFA of machines to produce \SI{1}{} FCFA of extra output per year ($v = 3$), then when demand rises by 100 million FCFA, firms must invest 300 million FCFA in new equipment. But notice the crucial implication: investment depends on the \emph{change} in output. If output keeps rising but \emph{more slowly}, $\Delta Y$ falls and investment collapses --- even though demand is still growing! This explains why booms can turn into recessions without any fall in consumption.

\begin{exoresolu}[Multiplier--accelerator interaction: a numerical illustration]
An economy has a marginal propensity to consume $MPC = 0.5$ and a capital--output ratio $v = 2$. In period 1, government spending rises permanently by 100 billion FCFA. Trace the effects on income and induced investment in the first three periods, assuming consumption in each period depends on the previous period's income: $C_t = 0.5\,Y_{t-1}$, and induced investment is $I_t = 2\,(Y_{t-1} - Y_{t-2})$. Take $Y_0 = 1000$.
\tcblower
\textbf{Solution.} We track the \emph{change} in income relative to the initial equilibrium ($\Delta Y_t = Y_t - 1000$). Each period, the change in income equals the sum of induced consumption, induced investment, and the autonomous government injection: $\Delta Y_t = \Delta C_t + \Delta I_t + \Delta G$, with $\Delta G = 100$.

\textbf{Period 1:} No lagged effects yet. $\Delta C_1 = 0$, $\Delta I_1 = 0$. $\Delta Y_1 = 100$, so $Y_1 = 1100$.

\textbf{Period 2:} Consumption responds to period-1 income change: $\Delta C_2 = 0.5 \times \Delta Y_1 = 0.5 \times 100 = 50$. Induced investment: $\Delta I_2 = v \times \Delta Y_1 = 2 \times 100 = 200$.
\[ \Delta Y_2 = 50 + 200 + 100 = 350 \quad\Rightarrow\quad Y_2 = 1350. \]

\textbf{Period 3:} $\Delta C_3 = 0.5 \times \Delta Y_2 = 0.5 \times 350 = 175$; $\Delta I_3 = 2 \times \Delta Y_2 = 2 \times 350 = 700$.
\[ \Delta Y_3 = 175 + 700 + 100 = 975 \quad\Rightarrow\quad Y_3 = 1975. \]

Income is rising fast (boom). But observe what will happen later: as the growth of income eventually slows, $\Delta Y$ shrinks, induced investment $I_t = v\Delta Y$ falls sharply, and through the multiplier income itself contracts --- the downturn begins. The multiplier and accelerator thus feed each other, generating a self-reinforcing upswing followed by a self-reinforcing downswing: the cycle.
\end{exoresolu}

\begin{popfigure}
\begin{center}
\begin{tikzpicture}[
node distance=1.1cm,
box/.style={draw=popblue, thick, fill=popblueL, rounded corners, text width=3.4cm, align=center, minimum height=1cm, font=\small},
arr/.style={->, thick, popdark}]
\node[box] (inv) {Change in autonomous investment or government spending};
\node[box, right=2.2cm of inv] (inc) {National income rises by a multiplied amount ($k = \frac{1}{1-MPC}$)};
\node[box, below=1.1cm of inc] (acc) {Rising income induces new investment: $I = v\,\Delta Y$};
\node[box, left=2.2cm of acc] (feed) {Extra investment raises demand again \ldots{} until growth slows, then $I$ collapses};
\draw[arr] (inv) -- node[above, font=\small]{multiplier} (inc);
\draw[arr] (inc) -- node[right, font=\small]{accelerator} (acc);
\draw[arr] (acc) -- (feed);
\draw[arr] (feed) -- node[left, font=\small]{multiplier again} (inv.south -| feed.west);
\draw[arr] (feed.north) -- (inv.south);
\end{tikzpicture}
\legende{The multiplier--accelerator interaction: each mechanism amplifies the other, generating cyclical fluctuations.}
\end{center}
\end{popfigure}

\courssub{Government Economic Objectives}

Because cycles impose heavy costs --- unemployment in slumps, inflation in booms --- governments intervene to stabilise and steer the economy. Their action is organised around a set of macroeconomic objectives.

\begin{aretenir}[The main macroeconomic objectives and their indicators]
\begin{itemize}
\item \textbf{Economic growth:} a sustained rise in real national output. \emph{Indicator:} rate of growth of real GDP (or real GDP per head).
\item \textbf{Full employment / low unemployment:} the highest possible use of the labour force. \emph{Indicator:} the unemployment rate (percentage of the labour force without work).
\item \textbf{Price stability:} a low, stable and predictable rate of inflation. \emph{Indicator:} the annual percentage change in the consumer price index (CPI).
\item \textbf{Balance of payments equilibrium:} avoiding persistent large deficits (or surpluses) on the external accounts. \emph{Indicator:} the current account balance, and the level of foreign exchange reserves.
\item \textbf{Equitable distribution of income and wealth:} reducing excessive inequalities and poverty. \emph{Indicator:} the Gini coefficient, poverty headcount ratios.
\item \textbf{Environmental sustainability:} growth that does not destroy the natural capital of future generations. \emph{Indicator:} measures of pollution, deforestation, and adjusted ``green'' national income.
\end{itemize}
The first four are the traditional ``magic square'' of macroeconomic policy; the last two have become increasingly prominent, including in Cameroon's long-term development vision documents (e.g. Vision 2035, SND30).
\end{aretenir}

\courssub{Conflicts and Trade-Offs Between Objectives}

The difficulty of economic policy is that these objectives cannot all be maximised at once: pursuing one often undermines another.

\textbf{1. Inflation versus unemployment.} Expanding aggregate demand to reduce unemployment (through higher government spending or easier credit) also bids up prices when the economy approaches full capacity. Conversely, deflating demand to reduce inflation raises unemployment. This is the logic of the \emph{Phillips curve}, studied formally in the next sections on unemployment and inflation.

\begin{methode}[Explaining a policy conflict with Phillips curve logic]
To analyse any demand-management dilemma in an essay:
\begin{enumerate}
\item Identify the policy and its effect on \emph{aggregate demand} (e.g. increased public investment shifts AD to the right).
\item Trace the \emph{beneficial} effect on the target objective: higher real output, more jobs, lower unemployment.
\item Trace the \emph{cost} on the competing objective: as spare capacity shrinks and labour markets tighten, wages and prices rise --- demand-pull inflation.
\item Conclude with the \emph{trade-off}: lower unemployment is bought at the price of higher inflation; the government must choose its preferred combination (its position on the Phillips curve).
\item Evaluate: in the long run, the expectations-augmented Phillips curve is vertical at the NAIRU (Non-Accelerating Inflation Rate of Unemployment). Supply-side policies (training, infrastructure, agribusiness development as in Section 4) can shift the NAIRU leftwards, easing the trade-off by increasing capacity.
\end{enumerate}
\end{methode}

\textbf{2. Growth versus balance of payments equilibrium.} Rapid growth raises household incomes, and part of that income is spent on imports. If exports do not grow as fast, the trade balance deteriorates. Cameroon experiences this regularly: investment booms suck in imported machinery, fuel and consumer goods, putting pressure on the foreign reserves pooled within the BEAC. Growth may also attract short-term capital inflows whose reversal can destabilise the external accounts.

\textbf{3. Growth versus equity.} Growth led by capital-intensive sectors (oil extraction, large agro-industrial plantations) may enrich a small urban elite while rural poverty persists. A government seeking equity through heavy redistribution and high taxation may, in turn, discourage investment and slow growth --- although well-designed redistribution (education, rural roads) can also support long-run growth.

\textbf{4. Growth versus the environment.} Expanding logging, mining or industrial output raises GDP but may degrade forests and pollute water, compromising sustainability.

\begin{definition}[The policy mix]
The \cle{policy mix} is the combination of policy instruments --- fiscal policy (taxation and public spending), monetary policy (money supply and interest rates, conducted for CEMAC by the BEAC), exchange rate policy, and supply-side policies --- that a government uses simultaneously to pursue its objectives while managing the conflicts between them.
\end{definition}

Because no single instrument can achieve several conflicting objectives at once, economists invoke the Tinbergen principle: a policymaker generally needs \emph{at least as many independent instruments as targets}. This is why modern policy combines demand management (to smooth the cycle) with supply-side measures (to raise trend growth and ease trade-offs).

\begin{exercice}[Applying the analysis]
\begin{enumerate}
\item Define the trade cycle and name its four phases, stating one characteristic of each.
\item Using $MPC = 0.75$ and $v = 2.5$, calculate the multiplier and the induced investment when income rises from 800 to 1000 billion FCFA.
\item Explain, using the multiplier--accelerator interaction, how a collapse in world cocoa prices could trigger a recession in Cameroon.
\item ``A government cannot simultaneously reduce inflation and reduce unemployment.'' Discuss, using Phillips curve logic.
\end{enumerate}
\end{exercice}

\begin{corrige}[Exercise 5]
\textbf{1.} The trade cycle is the recurring fluctuation of real output around its long-run trend. Boom: output and employment at maximum, inflationary pressure. Recession: falling output and rising unemployment. Trough: minimum activity, low investment, pessimism. Recovery: rising output, employment and confidence.

\textbf{2.} Multiplier: $k = \dfrac{1}{1 - 0.75} = 4$. Induced investment: $I = v\,\Delta Y = 2.5 \times (1000 - 800) = 500$ billion FCFA.

\textbf{3.} A fall in cocoa prices reduces export earnings and farm incomes; consumption falls; through the multiplier (in reverse) national income falls by a larger amount. As income growth turns negative, $\Delta Y < 0$, so induced investment $I = v\Delta Y$ collapses or becomes negative (firms scrap capacity). Falling investment further reduces demand via the multiplier, deepening the downturn --- exactly the mechanism observed in Cameroon after 1986.

\textbf{4.} In the short run, the statement is broadly correct: expanding aggregate demand to cut unemployment generates demand-pull inflation, and restrictive demand policies that cut inflation raise unemployment --- a movement along the Phillips curve. However, the trade-off can be eased: supply-side policies (training, infrastructure, agricultural modernisation) increase productive capacity, allowing lower unemployment \emph{and} stable prices. Moreover, credible, stable policies anchor expectations and reduce the inflationary cost of expansion. In the long run, the Phillips curve is vertical at the NAIRU, so there is no permanent trade-off. The conflict is real in the short run but not absolute.
\end{corrige}

\begin{attention}[Examination warnings]
\begin{itemize}
\item Do not confuse the \textbf{trend} (long-run growth) with the \textbf{cycle} (short-run fluctuations around it). A country can be in recession while its trend growth remains positive.
\item In the accelerator formula, investment depends on the \emph{change} in output, $\Delta Y$, not on the \emph{level} of output $Y$. Confusing the two is the most frequent error.
\item When discussing government objectives, always mention the \emph{indicator} used to measure each one, and always discuss at least one \emph{conflict} between objectives --- a mere list of objectives earns few marks at Advanced Level.
\item Remember that for CEMAC countries, monetary policy and exchange rate policy are conducted at the regional level (BEAC), which constrains the national policy mix.
\end{itemize}
\end{attention}

Having seen that one of the central objectives of government is to keep unemployment low --- and that recessions are precisely the periods when joblessness surges --- we now turn, in the next section, to a detailed study of \emph{unemployment}: its types, its causes, its measurement, and its particular features in Cameroon.

\coursec{Section 6: Unemployment — Types, Causes, Measurement and the Cameroonian Context}

In Section 5 we saw that government economic objectives include \cle{full employment}, and that during the downswing of the trade cycle, firms lay off workers as aggregate demand falls. Unemployment is therefore both a symptom of macroeconomic instability and a permanent structural challenge, especially in developing economies like Cameroon. It is one of the most heavily examined topics at the GCE Advanced Level, both in essays and in data-response questions. This section defines and measures unemployment, classifies its types and causes, evaluates its consequences, and analyses the policies available to reduce it, before applying all of this to the Cameroonian labour market.

\courssub{Defining and Measuring Unemployment}

\textbf{The labour force.} Not everyone who is without a job is unemployed. A student in Upper Sixth, a retired pensioner in Bafoussam and a full-time housewife are not counted as unemployed because they are not seeking work. Economists therefore begin with the concept of the \cle{labour force} (or working population): all persons of working age who are either \emph{employed} or \emph{actively seeking employment}.

\begin{definition}[Unemployment, underemployment and the unemployment rate]
According to the \cle{International Labour Organisation (ILO)}, a person is \textbf{unemployed} if he or she is of working age, is \emph{without work}, is \emph{available for work}, and has \emph{actively sought work} in a recent reference period (usually the past four weeks).
\par\smallskip
\textbf{Underemployment} exists when workers are employed but not fully utilised: either they work fewer hours than they wish (visible underemployment) or they work in jobs far below their qualifications and productivity (invisible underemployment). \textbf{Disguised unemployment} is a form of underemployment common in traditional agriculture: five family members cultivate a plot that three could farm, so the marginal product of the extra workers is approximately zero.
\par\smallskip
The \textbf{unemployment rate} is:
\[
\text{Unemployment rate} = \frac{\text{Number of unemployed}}{\text{Labour force}} \times 100
\]
\end{definition}

\begin{exemplebox}[Calculating the unemployment rate]
Suppose a country has a working-age population of 12 million. Of these, 8 million are in the labour force, of whom 7.2 million are employed. The number unemployed is $8 - 7.2 = 0.8$ million, so:
\[
\text{Unemployment rate} = \frac{0.8}{8} \times 100 = 10\%
\]
Note that the 4 million people outside the labour force (students, retirees, discouraged workers) do not enter the calculation.
\end{exemplebox}

\textbf{Measuring unemployment in practice.} Two broad methods exist:
\begin{itemize}
\item The \textbf{claimant count}: the number of people registered as unemployed and claiming unemployment-related benefits. It is cheap and quick, but it \emph{understates} true unemployment where benefit systems are weak — in Cameroon very few unemployed persons receive any benefit, so registration data capture only a tiny fraction of joblessness.
\item \textbf{Labour force surveys} (such as Cameroon's \emph{Enquêtes sur l'Emploi et le Secteur Informel} conducted by the National Institute of Statistics): households are sampled and questioned using ILO criteria. This is more accurate and also reveals underemployment and informal-sector work, but it is costly and infrequent.
\end{itemize}

\begin{attention}[Common mistake]
Do not quote the official unemployment rate of a developing country as the whole story. In Cameroon, the measured \emph{unemployment} rate may look modest, but \emph{underemployment} — especially in the informal sector — is very high. A graduate selling phone credit at a roadside kiosk in Douala is ``employed'' in the statistics, yet clearly underemployed. Always mention both dimensions.
\end{attention}

\begin{popfigure}
\begin{tikzpicture}
\begin{axis}[
    width=0.85\linewidth, height=6.5cm,
    ybar, bar width=14pt,
    ymin=0, ymax=45,
    ylabel={Rate (\%)},
    symbolic x coords={15--24, 25--34, 35--49, 50+},
    xtick=data,
    legend style={at={(0.5,1.02)}, anchor=south, legend columns=2, font=\small},
    nodes near coords, nodes near coords style={font=\scriptsize},
    ymajorgrids, grid style={popline!40},
    axis line style={popdark},
    tick label style={font=\small},
]
\addplot[fill=popblue, draw=popdark] coordinates {(15--24,12) (25--34,7) (35--49,4) (50+,3)};
\addplot[fill=poporange, draw=popdark] coordinates {(15--24,38) (25--34,30) (35--49,22) (50+,18)};
\legend{Unemployment, Underemployment}
\end{axis}
\end{tikzpicture}
\legende{Illustrative pattern of unemployment and underemployment by age group in Cameroon: both fall with age, and underemployment exceeds unemployment in every group — the hallmark of a large informal sector (illustrative values in the spirit of INS employment surveys).}
\end{popfigure}

\courssub{Types and Causes of Unemployment}

The single most important skill in this topic is to \emph{diagnose the type} of unemployment before prescribing a cure, because each type has a different cause and therefore a different remedy.

\begin{definition}[Types of unemployment]
\begin{itemize}
\item \textbf{Frictional unemployment}: temporary unemployment while workers move between jobs or search for their first job. It exists even in a healthy economy because information and matching take time.
\item \textbf{Structural unemployment}: unemployment caused by a long-term mismatch between the skills workers have and the skills employers demand, or between the location of workers and the location of jobs — typically following the decline of an industry or changes in technology and consumer tastes.
\item \textbf{Cyclical (demand-deficient) unemployment}: unemployment caused by a general deficiency of aggregate demand during the downswing or recession phase of the trade cycle.
\item \textbf{Seasonal unemployment}: unemployment linked to seasons — farm labourers between harvests, tourist-sector workers in the low season.
\item \textbf{Technological unemployment}: a form of structural unemployment in which machines and automation replace labour (e.g.\ automated teller machines reducing bank cashiers).
\item \textbf{Classical (real-wage) unemployment}: unemployment caused by real wages being held \emph{above} the market-clearing level — by minimum wage laws, strong trade unions or efficiency wages — so that the supply of labour exceeds demand.
\item \textbf{Voluntary unemployment}: persons who choose not to work at the going wage, for instance because reservation wages or family support make work unattractive.
\end{itemize}
\end{definition}

\begin{center}
\begin{tabularx}{\linewidth}{|l|X|X|X|}
\hline
\rowcolor{popblueL} \textbf{Type} & \textbf{Cause} & \textbf{Example} & \textbf{Appropriate policy} \\
\hline
Frictional & Job search takes time & Graduate seeking first job in Yaoundé & Job centres, better labour-market information \\
\hline
Structural & Skills/location mismatch & Retrenched textile worker after factory closure & Retraining, education reform, regional policy \\
\hline
Cyclical & Deficient aggregate demand & Lay-offs during a recession & Expansionary fiscal and monetary policy \\
\hline
Seasonal & Seasonal demand & Cocoa farm labourers off-season & Diversification, off-season activities \\
\hline
Technological & Automation replaces labour & Bank cashiers replaced by ATMs & Re-skilling, support for new sectors \\
\hline
Classical & Real wage above equilibrium & Minimum wage above market-clearing level & Wage flexibility (controversial), raising labour demand \\
\hline
Voluntary & Choice not to work at going wage & Person living on family support & Work incentives, benefit reform \\
\hline
\end{tabularx}
\end{center}

\textbf{Cyclical unemployment and the deflationary gap.} When aggregate demand falls short of the level needed for full employment, the economy settles at an equilibrium output below potential output. The difference is the \cle{deflationary gap}, and the associated joblessness is demand-deficient unemployment.

\begin{popfigure}
\begin{tikzpicture}
\begin{axis}[
    width=0.8\linewidth, height=7cm,
    axis lines=left,
    xlabel={Real national output (income)},
    ylabel={Price level},
    xmin=0, xmax=10, ymin=0, ymax=10,
    xtick=\empty, ytick=\empty,
    axis line style={popdark},
    clip=false,
]
\addplot[domain=1:9, samples=50, thick, popblue] {8 - 0.75*x};
\node[popblue, font=\small] at (axis cs:8.6,1.9) {$AD_1$};
\addplot[domain=1:9, samples=50, thick, popblue!60] {6.2 - 0.6*x};
\node[popblue!70, font=\small] at (axis cs:8.9,0.9) {$AD_2$};
\addplot[domain=0.8:8.5, samples=50, thick, poporange] {0.55*x + 0.5};
\node[poporange, font=\small] at (axis cs:8.2,5.6) {$AS$};
\addplot[thick, popgreen, dashed] coordinates {(7.5,0) (7.5,9)};
\node[popgreen, font=\small] at (axis cs:7.5,9.4) {$Y_f$};
\draw[dashed, popmuted] (axis cs:5.3,0) -- (axis cs:5.3,3.4);
\node[font=\small] at (axis cs:5.3,-0.5) {$Y_1$};
\node[font=\small] at (axis cs:7.5,-0.5) {$Y_f$};
\draw[<->, thick, poppink] (axis cs:5.3,0.35) -- (axis cs:7.5,0.35);
\node[poppink, font=\small, align=center] at (axis cs:6.4,0.85) {deflationary gap};
\end{axis}
\end{tikzpicture}
\legende{Demand-deficient unemployment: a fall in aggregate demand from $AD_1$ to $AD_2$ pushes equilibrium output below the full-employment level $Y_f$, creating a deflationary gap and cyclical unemployment.}
\end{popfigure}

\begin{definition}[The natural rate of unemployment]
The \cle{natural rate of unemployment} is the rate of unemployment that persists when the labour market is in equilibrium and aggregate demand is at its full-employment level. It consists of frictional and structural unemployment — the unemployment that cannot be removed by boosting demand. The related concept of the \cle{NAIRU} (Non-Accelerating Inflation Rate of Unemployment) is the rate of unemployment at which inflation is stable; attempting to push unemployment below it through demand expansion only accelerates inflation. \emph{(Exam flag: define it, state its components, and explain why demand policies cannot reduce it.)}
\end{definition}

\courssub{Consequences of Unemployment}

Unemployment imposes costs at three levels:
\begin{itemize}
\item \textbf{Economic costs}: lost output (the economy operates inside its production possibility frontier); lost tax revenue and higher public spending on social support, worsening the budget balance; a waste of human capital as skills deteriorate during long jobless spells.
\item \textbf{Social costs}: rising poverty and inequality; links to crime, urban insecurity and social unrest; \cle{brain drain}, as frustrated graduates emigrate, depriving Cameroon of skills it paid to produce.
\item \textbf{Individual costs}: loss of income and falling living standards; loss of self-esteem, stress and health problems; deskilling, which makes re-employment progressively harder (the ``scarring effect'').
\end{itemize}

\courssub{Policies to Reduce Unemployment}

\begin{methode}[Matching the policy to the type of unemployment]
In data-response and essay questions, proceed in three steps:
\begin{enumerate}
\item \textbf{Diagnose}: identify from the data which type(s) of unemployment dominate (e.g.\ economy-wide lay-offs in a recession $\Rightarrow$ cyclical; declining industries with unfilled vacancies elsewhere $\Rightarrow$ structural).
\item \textbf{Match}: choose the policy aimed at that cause — demand management for cyclical unemployment; supply-side measures (training, mobility, flexibility) for structural and frictional unemployment.
\item \textbf{Evaluate}: discuss time lags, cost, side effects (inflation risk from demand expansion) and whether the policy suits a developing economy.
\end{enumerate}
\end{methode}

The main policy instruments are:
\begin{itemize}
\item \textbf{Demand management}: expansionary fiscal policy (higher government spending, tax cuts) and monetary policy (lower interest rates) to close a deflationary gap. Effective only against \emph{cyclical} unemployment.
\item \textbf{Education and vocational training}: aligning curricula with labour-market needs (technical education, ICT, agribusiness skills) to cure structural and frictional unemployment.
\item \textbf{Labour market flexibility}: reducing rigidities, improving job information and placement services, encouraging mobility between regions and occupations.
\item \textbf{Regional and industrial policy}: attracting investment to depressed regions and promoting new industries to absorb workers from declining ones.
\item \textbf{Support to SMEs and self-employment}: microfinance, tax incentives and business training, since small enterprises are the main creators of jobs in Cameroon.
\end{itemize}

\begin{attention}[Common mistake]
Do not treat \emph{all} unemployment as curable by raising aggregate demand. If unemployment is structural — workers with the wrong skills in the wrong places — boosting demand mainly creates \emph{inflation} and unfilled vacancies, not jobs. Structural unemployment requires \emph{supply-side} measures: training, mobility and industrial restructuring. This is precisely the message of the natural rate / NAIRU analysis.
\end{attention}

\courssub{Guided Work 7 — Problems of Unemployment in Cameroon}

\begin{experience}[Case study: unemployment and underemployment in Cameroon]
\textbf{Facts and features.} Cameroon's labour market is marked by (i) high \emph{youth and graduate unemployment}: thousands of university graduates each year compete for very few formal jobs, and many remain jobless for years after graduation; (ii) a dominant \emph{informal sector} — the great majority of workers earn their living in small trade, transport (e.g.\ motorcycle taxis, ``benskin'' riders), artisanal activities and subsistence agriculture, with low productivity, no contracts and no social protection; (iii) severe \emph{underemployment}, which affects young people far more than open unemployment.
\par\smallskip
\textbf{Causes.}
\begin{itemize}
\item \emph{Skills mismatch}: general academic education produces graduates whose profiles do not match the technical and entrepreneurial skills employers demand.
\item \emph{Slow and job-poor growth}: economic growth has not been rapid or labour-intensive enough to absorb the large annual flow of new entrants into the labour market.
\item \emph{Public sector saturation}: the civil service, traditionally the main destination of graduates, can no longer absorb them, while the formal private sector remains small.
\item \emph{Rapid population growth and rural exodus}, which swell urban labour supply faster than urban jobs are created.
\end{itemize}
\par\smallskip
\textbf{Government responses.} Among the measures deployed are the \cle{PIFMAS} (\emph{Projet Intégré de Formation et de Financement des Micro-entreprises Agricoles et de Services}), which trains and finances young entrepreneurs in agriculture and services; youth employment and insertion programmes (such as the national youth employment plans and support funds); the expansion of \emph{vocational and technical training} (including agricultural and crafts training centres); and incentives for SMEs and self-employment through microfinance and business-creation support.
\par\smallskip
\textbf{Evaluation.} These programmes attack structural and frictional unemployment at its roots, but their scale remains small relative to the number of job seekers, funding and follow-up are often inadequate, and durable progress requires faster, more labour-intensive growth — notably through second generation agriculture and agribusiness (Section 4) and industrialisation (Section 2).
\end{experience}

\begin{exoresolu}[Worked exercise: diagnosing and prescribing]
A country experiences simultaneously: (a) a recession in which 200\,000 workers across all sectors lose their jobs; (b) the permanent closure of its shipbuilding industry, leaving 40\,000 specialised workers jobless; (c) 15\,000 school leavers taking an average of four months to find their first job. Identify each type of unemployment and propose one appropriate policy for each.
\tcblower
\textbf{Solution.}
\begin{itemize}
\item (a) The economy-wide, recession-linked job losses are \textbf{cyclical (demand-deficient) unemployment}. Appropriate policy: \emph{expansionary fiscal and/or monetary policy} to raise aggregate demand and close the deflationary gap.
\item (b) The shipbuilding workers suffer \textbf{structural unemployment}: their specific skills no longer match the demands of a changed industrial structure. Appropriate policy: \emph{retraining and re-skilling programmes}, plus regional/industrial policy to attract new employers to the affected area. Raising aggregate demand would \emph{not} restore their jobs.
\item (c) The school leavers experience \textbf{frictional unemployment}. Appropriate policy: \emph{improved job information and placement services} (job centres, school-to-work counselling) to shorten the search period.
\end{itemize}
Key lesson: one symptom (joblessness), three diagnoses, three different cures.
\end{exoresolu}

\begin{exercice}[Practice]
\begin{enumerate}
\item Define the labour force and the unemployment rate. A country has a labour force of 10 million and 9.1 million employed persons; calculate its unemployment rate.
\item Explain, with the aid of an AD/AS diagram, how a fall in aggregate demand causes cyclical unemployment.
\item Distinguish structural from frictional unemployment and explain why the policies appropriate to each differ.
\item ``In Cameroon, underemployment is a more serious problem than open unemployment.'' Discuss with reference to the informal sector.
\item Outline two government measures to reduce youth unemployment in Cameroon and evaluate their likely effectiveness.
\end{enumerate}
\end{exercice}

\begin{corrige}[Exercise 6]
\begin{enumerate}
\item Labour force: all persons of working age who are employed or actively seeking work. Unemployment rate $= \dfrac{10 - 9.1}{10}\times 100 = 9\%$.
\item A fall in AD shifts the AD curve leftward; equilibrium output falls below the full-employment level $Y_f$, opening a deflationary gap; firms reduce output and lay off workers — demand-deficient unemployment (see the figure above).
\item Structural unemployment arises from a lasting mismatch of skills or location and requires supply-side cures (retraining, mobility, industrial policy); frictional unemployment is short-term search unemployment cured by better information and placement. Demand expansion helps neither directly.
\item Strong answers note that measured unemployment understates job market distress: most workers cannot afford to be openly unemployed, so they accept low-productivity informal work; hence underemployment (visible and invisible) is widespread, especially among the youth, while open unemployment is concentrated among graduates and urban youth.
\item For example: vocational and technical training aligned to market needs (effective against skills mismatch but slow and costly); PIFMAS-type support to youth agribusiness and SMEs (creates self-employment, but limited by funding, follow-up and market access). Evaluation should weigh scale against the number of annual labour-market entrants.
\end{enumerate}
\end{corrige}

\begin{aretenir}[Key points of Section 6]
\begin{itemize}
\item Unemployment (ILO): without work, available for work, actively seeking work; the unemployment rate is the unemployed divided by the labour force, times 100.
\item In developing economies, \emph{underemployment} and the informal sector matter as much as open unemployment.
\item Know the seven types — frictional, structural, cyclical, seasonal, technological, classical, voluntary — each with its cause and its cure.
\item Cyclical unemployment stems from a deflationary gap; the natural rate (NAIRU) is the frictional-plus-structural core that demand policy cannot remove.
\item Unemployment costs output, public finances, social cohesion and individual well-being (including brain drain).
\item In Cameroon: youth and graduate unemployment, skills mismatch, slow job-rich growth, public sector saturation; responses include PIFMAS, vocational training and SME support.
\end{itemize}
\end{aretenir}

\coursec{Section 7: Inflation and Deflation}

In Section 6 we studied unemployment, one of the most painful macroeconomic problems a government can face. We now turn to its traditional rival: \cle{inflation}. The two are often linked --- so closely that economists spent decades debating whether a government must ``choose'' between them (the famous Phillips curve, which we shall examine below). For a Cameroonian household, inflation is not an abstract number: it is the price of a bag of rice in Mokolo market, the cost of transport from Douala to Yaound\'e, the school fees that rise faster than salaries. Understanding inflation --- its measurement, its causes, its effects and its cures --- is therefore central to evaluating government economic policy, and it is a guaranteed topic at the GCE Advanced Level.

\courssub{Definition and Measurement of Inflation}

\begin{definition}[Inflation, disinflation and deflation]
\cle{Inflation} is a \textbf{sustained and generalised rise in the general price level} over time. Equivalently, it is a fall in the purchasing power of money: each franc buys fewer goods than before.
\begin{itemize}
\item \cle{Disinflation} is a \textbf{fall in the rate of inflation}: prices are still rising, but more slowly (e.g. inflation falls from 6\% to 3\%).
\item \cle{Deflation} is a \textbf{sustained fall in the general price level} (a negative inflation rate).
\end{itemize}
\end{definition}

Three points in this definition deserve emphasis. First, inflation concerns the \emph{general} price level: a rise in the price of petrol alone is not inflation. Second, it must be \emph{sustained}: a one-off jump in prices after a tax change is a price shock, not ongoing inflation. Third, inflation means money loses value --- if prices double, the real value of 10\,000 FCFA is halved.

\textbf{Measuring inflation: the Consumer Price Index (CPI).} Statisticians (in Cameroon, the National Institute of Statistics) track the prices of a representative ``basket'' of goods and services consumed by a typical household --- food, housing, transport, clothing, health, education. Each item is weighted according to its share in household spending. The \cle{Consumer Price Index} compares the cost of this basket in a given year with its cost in a base year (where the index is set to 100):
\[ \mathrm{CPI}_t = \frac{\text{cost of the basket in year } t}{\text{cost of the basket in the base year}} \times 100. \]

\begin{propriete}[The inflation rate formula]
The rate of inflation between year 1 and year 2 is the percentage change in the CPI:
\[ \pi = \frac{\mathrm{CPI}_2 - \mathrm{CPI}_1}{\mathrm{CPI}_1} \times 100. \]
This formula is examinable: candidates must be able to compute an inflation rate from index numbers.
\end{propriete}

\begin{exoresolu}[Calculating a price index and an inflation rate]
A simplified household basket costs 80\,000 FCFA in the base year (Year 0). The same basket costs 88\,000 FCFA in Year 1 and 92\,400 FCFA in Year 2.
\begin{enumerate}
\item Calculate the CPI for Years 1 and 2.
\item Calculate the inflation rate between Year 1 and Year 2.
\end{enumerate}
\tcblower
\textbf{Solution.}
\begin{enumerate}
\item $\mathrm{CPI}_1 = \dfrac{88\,000}{80\,000}\times 100 = 110$; \quad $\mathrm{CPI}_2 = \dfrac{92\,400}{80\,000}\times 100 = 115.5$.
\item $\pi = \dfrac{115.5 - 110}{110}\times 100 = 5\%$. Prices rose by 5\% during Year 2.
\end{enumerate}
\textbf{Interpretation:} a worker whose nominal salary rose by only 3\% over the same year suffered a fall of about 2\% in real purchasing power.
\end{exoresolu}

\courssub{Types and Causes of Inflation}

Economists classify inflation by its origin. The two great families are \cle{demand-pull} and \cle{cost-push} inflation, to which we add monetary and structural explanations.

\textbf{1. Demand-pull inflation.} When aggregate demand (consumption, investment, government spending, exports) grows faster than the economy's capacity to produce, buyers compete for limited goods and bid prices up: ``too much money chasing too few goods.'' On an AD/AS diagram, the AD curve shifts right and the price level rises.

\begin{popfigure}
\begin{center}
\begin{tikzpicture}[scale=1.0]
\draw[->, thick] (0,0) -- (6.2,0) node[below left]{\small Real output $Y$};
\draw[->, thick] (0,0) -- (0,4.6) node[below left]{\small Price level $P$};
\draw[very thick, popblue] (0.6,4.2) -- (5.4,0.6) node[pos=0.05, above right, popblue]{\small $AD_1$};
\draw[very thick, poporange] (1.6,4.2) -- (6.0,0.9) node[pos=0.02, above right, poporange]{\small $AD_2$};
\draw[very thick, popgreen] (0.8,0.5) -- (5.6,4.3) node[pos=0.9, below left, popgreen]{\small $SRAS$};
\fill[popink] (2.72,2.32) circle (1.5pt); \fill[popink] (3.55,2.75) circle (1.5pt);
\draw[dashed, gray] (2.72,2.32) -- (0,2.32) node[left]{\small $P_1$};
\draw[dashed, gray] (3.55,2.75) -- (0,2.75) node[left]{\small $P_2$};
\draw[dashed, gray] (2.72,2.32) -- (2.72,0) node[below]{\small $Y_1$};
\draw[dashed, gray] (3.55,2.75) -- (3.55,0) node[below]{\small $Y_2$};
\end{tikzpicture}
\end{center}
\legende{Demand-pull inflation: a rightward shift of aggregate demand raises the price level from $P_1$ to $P_2$.}
\end{popfigure}

\textbf{2. Cost-push inflation.} Here prices rise because firms' costs of production rise, and firms pass these costs on to consumers. The short-run aggregate supply curve shifts left. Three sources are distinguished:
\begin{itemize}
\item \emph{Wage-push}: trade unions win wage increases above productivity growth, raising unit labour costs.
\item \emph{Imported inflation}: a rise in the price of imported inputs (fuel, wheat, machinery) or a depreciation of the currency. This is crucial for Cameroon, which imports fuel, food and manufactured goods; the 1994 devaluation of the CFA franc doubled the CFA price of imports overnight.
\item \emph{Profit-push}: firms with market power widen their profit margins.
\end{itemize}

\begin{popfigure}
\begin{center}
\begin{tikzpicture}[scale=1.0]
\draw[->, thick] (0,0) -- (6.2,0) node[below left]{\small Real output $Y$};
\draw[->, thick] (0,0) -- (0,4.6) node[below left]{\small Price level $P$};
\draw[very thick, popblue] (0.6,4.2) -- (5.6,0.5) node[pos=0.05, above right, popblue]{\small $AD$};
\draw[very thick, popgreen] (1.6,0.5) -- (5.8,4.3) node[pos=0.9, below left, popgreen]{\small $SRAS_1$};
\draw[very thick, poporange] (0.5,0.6) -- (4.6,4.3) node[pos=0.9, below left, poporange]{\small $SRAS_2$};
\fill[popink] (3.9,2.42) circle (1.5pt); \fill[popink] (3.1,2.93) circle (1.5pt);
\draw[dashed, gray] (3.9,2.42) -- (0,2.42) node[left]{\small $P_1$};
\draw[dashed, gray] (3.1,2.93) -- (0,2.93) node[left]{\small $P_2$};
\draw[dashed, gray] (3.9,2.42) -- (3.9,0) node[below]{\small $Y_1$};
\draw[dashed, gray] (3.1,2.93) -- (3.1,0) node[below]{\small $Y_2$};
\end{tikzpicture}
\end{center}
\legende{Cost-push inflation: a leftward shift of SRAS raises the price level but \emph{reduces} output ($Y_2 < Y_1$) --- the recipe for stagflation.}
\end{popfigure}

\begin{methode}[Distinguishing demand-pull from cost-push inflation on an AD/AS diagram]
\begin{enumerate}
\item Identify which curve shifted: $AD$ shifting right $\Rightarrow$ demand-pull; $SRAS$ shifting left $\Rightarrow$ cost-push.
\item Check the direction of output: demand-pull raises both $P$ and $Y$ in the short run; cost-push raises $P$ but \textbf{lowers} $Y$.
\item Conclude on the policy dilemma: cost-push inflation is harder to fight, because reducing demand to lower prices worsens unemployment (stagflation).
\end{enumerate}
\end{methode}

\textbf{3. Monetary inflation.} Monetarists, following Milton Friedman, argue that ``inflation is always and everywhere a monetary phenomenon'': if the money supply grows persistently faster than real output, prices must rise.

\begin{propriete}[The quantity theory of money]
The \textbf{Fisher equation of exchange} states:
\[ MV = PY, \]
where $M$ is the money supply, $V$ the velocity of circulation, $P$ the price level and $Y$ real output. Assuming $V$ is stable and $Y$ is determined by real factors in the long run, an increase in $M$ translates proportionally into an increase in $P$: money growth causes inflation. The statement itself is examinable; a full proof is not required.
\end{propriete}

\begin{exemplebox}[Applying $MV = PY$]
Suppose $M = 2\,000$ billion FCFA, $V = 3$ and $Y = 12\,000$ billion FCFA of real output. Then $P = \dfrac{MV}{Y} = \dfrac{2\,000 \times 3}{12\,000} = 0.5$. If the central bank doubles $M$ to 4\,000 billion while $V$ and $Y$ are unchanged, $P$ doubles to $1.0$: the price level rises by 100\%. This illustrates why uncontrolled money creation (for example, printing money to finance a budget deficit) is inflationary.
\end{exemplebox}

\textbf{4. Structural (bottleneck) inflation.} In developing economies, inflation often arises from supply rigidities: poor roads raise transport costs, electricity shortages constrain production, and agricultural output cannot respond quickly to rising urban demand. These ``bottlenecks'' mean prices rise even when aggregate demand is not excessive --- a point highly relevant to Cameroon.

\courssub{The Effects of Inflation}

Inflation is not neutral: it redistributes income and distorts decisions.

\begin{itemize}
\item \textbf{Purchasing power:} money incomes buy less; workers on fixed nominal wages and pensioners lose most.
\item \textbf{Savers and lenders lose, borrowers gain:} if inflation is 6\% while a savings account pays 3\%, the real return is negative. Debtors repay in ``cheaper'' money --- the real burden of debt falls.
\item \textbf{Income distribution:} inflation arbitrarily redistributes wealth from creditors to debtors and from those on fixed incomes to those who can index their earnings.
\item \textbf{International competitiveness:} if Cameroon's inflation exceeds that of its trading partners, Cameroonian exports become dearer and imports cheaper, worsening the trade balance (though within CEMAC the common currency limits this divergence).
\item \textbf{Menu costs and shoe-leather costs:} firms must constantly reprint price lists (menu costs), and households waste time and effort minimising cash holdings (shoe-leather costs).
\item \textbf{Uncertainty:} high and variable inflation discourages investment and long-term planning.
\end{itemize}

A brief nuance: \emph{moderate} inflation (around 2--3\%) can be beneficial. It ``greases the wheels'' of the labour market (real wages can fall without nominal cuts), reduces the real burden of debt, and keeps the economy safely away from deflation. This is why the CEMAC convergence criterion sets an inflation ceiling of 3\% rather than zero.

\courssub{The Phillips Curve: the Inflation--Unemployment Trade-off}

In 1958, A. W. Phillips observed a stable inverse relationship between wage inflation and unemployment in British data. Generalised to price inflation, the \cle{Phillips curve} suggests that governments face a trade-off: lower unemployment can be ``bought'' at the price of higher inflation, and vice versa.

\begin{propriete}[The short-run Phillips trade-off]
In the \textbf{short run}, there is an inverse relationship between the inflation rate and the unemployment rate: expanding aggregate demand reduces unemployment but raises inflation. In the \textbf{long run}, once workers and firms adjust their expectations, the economy returns to the \textbf{natural rate of unemployment} and the long-run Phillips curve is \textbf{vertical} at that rate: there is no permanent trade-off. (This expectations-augmented extension is a frequent exam question.)
\end{propriete}

\begin{popfigure}
\begin{center}
\begin{tikzpicture}
\begin{axis}[
  width=10cm, height=7cm,
  axis lines=left,
  xlabel={Unemploying rate (\%)}, ylabel={Inflation rate (\%)},
  xmin=0, xmax=10, ymin=0, ymax=10,
  xtick={5}, xticklabels={$U^{*}$},
  ytick=\empty,
  legend style={at={(0.97,0.97)}, anchor=north east, font=\small},
  every axis plot/.append style={very thick}
]
\addplot[popblue, domain=1.2:9, samples=100] {12/x};
\addlegendentry{Short-run Phillips curve}
\addplot[poporange, domain=0:10, samples=2, dashed] coordinates {(5,0) (5,9.5)};
\addlegendentry{Long-run curve (natural rate $U^{*}$)}
\end{axis}
\end{tikzpicture}
\end{center}
\legende{The Phillips curve: a short-run trade-off between inflation and unemployment, but a vertical long-run curve at the natural rate $U^{*}$.}
\end{popfigure}

The intuition for the vertical long-run curve is expectations. If the government accepts 8\% inflation to cut unemployment, workers eventually expect 8\% inflation and demand matching wage rises; the stimulus evaporates and unemployment returns to its natural rate, but now with higher inflation. Attempts to exploit the trade-off then produce \emph{stagflation} --- high inflation \emph{and} high unemployment --- as many industrialised countries experienced in the 1970s.

\courssub{Anti-Inflationary Policies}

\begin{itemize}
\item \textbf{Monetary policy:} raising interest rates and restricting money-supply growth reduce spending and credit. In the CEMAC zone, this is the role of the \cle{BEAC} (Bank of Central African States), which sets key policy rates, while \cle{COBAC} supervises commercial banks. Cameroon, having pooled its monetary sovereignty, cannot conduct an independent national monetary policy.
\item \textbf{Fiscal policy:} cutting government spending or raising taxes reduces aggregate demand (deflationary fiscal policy in the old sense of ``deflating demand'').
\item \textbf{Supply-side policies:} investment in roads, energy, ports, education and competition policy raises productive capacity, attacking structural and cost-push inflation at its root --- slow but durable.
\item \textbf{Prices and incomes policy:} direct controls on prices and wage increases. These can calm inflation temporarily but create shortages and black markets if maintained.
\item \textbf{Exchange rate policy:} a strong or appreciated currency cheapens imports and fights imported inflation. Within the CFA franc zone, the fixed peg to the euro anchors import prices, and devaluation (as in 1994) does the opposite.
\end{itemize}

\courssub{Deflation: Definition, Causes and Consequences}

\begin{attention}[Deflation is not disinflation]
A classic exam trap. \textbf{Disinflation} = the inflation rate \emph{falls} but remains positive (from 6\% to 2\%): prices still rise. \textbf{Deflation} = the inflation rate is \emph{negative}: the general price level actually falls. Confusing the two costs marks every year.
\end{attention}

\textbf{Causes of deflation.} Deflation may come from a collapse of aggregate demand (recession, falling investment, credit contraction) --- the harmful kind --- or from technological progress and productivity gains that lower production costs (the ``benign'' kind observed in prices of electronics).

\textbf{Consequences.} Falling prices sound pleasant, but demand-side deflation is dangerous:
\begin{itemize}
\item \emph{Deferred spending:} if consumers expect prices to be lower tomorrow, they postpone purchases, deepening the fall in demand.
\item \emph{Rising real debt:} deflation increases the real value of debts, strangling indebted firms and households (the mirror image of inflation).
\item \emph{The deflationary spiral:} falling demand $\rightarrow$ falling prices $\rightarrow$ falling profits and wages $\rightarrow$ lay-offs $\rightarrow$ still lower demand. Japan's ``lost decades'' illustrate this trap.
\end{itemize}

\textbf{Policies against deflation.} Governments and central banks practise \cle{reflation}: cutting interest rates, expanding the money supply and increasing public spending to revive aggregate demand. When interest rates approach zero, central banks may resort to \emph{quantitative easing} (large-scale purchase of financial assets to inject liquidity) or even \emph{negative interest rates} (charging banks for holding reserves, to push them to lend). These tools, used by the European Central Bank and the Bank of Japan, remain largely theoretical for the BEAC, whose room for manoeuvre is constrained by the fixed exchange rate and reserve-cover rules.

\courssub{The Cameroonian Context}

\begin{savaistu}[Inflation episodes in Cameroon]
Cameroon's modern inflation history is marked by two salient episodes. First, the \textbf{January 1994 devaluation} of the CFA franc (from 50 to 100 FCFA per French franc) doubled the domestic-currency price of imports, producing a sharp burst of imported inflation --- with annual inflation reaching roughly 30\% in 1994 before subsiding. Second, \textbf{from 2021 onwards}, the combined effects of the COVID-19 supply disruptions, the war in Ukraine and higher fuel and food import prices pushed inflation above the CEMAC ceiling, squeezing household purchasing power even though the BEAC tightened monetary policy. As a member of CEMAC, Cameroon is bound by the \textbf{convergence criterion of a maximum 3\% average annual inflation rate}, designed to protect the stability of the common currency.
\end{savaistu}

\begin{popfigure}
\begin{center}
\begin{tikzpicture}
\begin{axis}[
  width=11.5cm, height=6.5cm,
  axis lines=left,
  xlabel={Year}, ylabel={Inflation rate (\%)},
  xmin=1990, xmax=2023, ymin=-4, ymax=32,
  xtick={1990,1995,2000,2005,2010,2015,2020},
  ytick={0,10,20,30},
  legend style={at={(0.03,0.97)}, anchor=north west, font=\small}
]
\addplot[popblue, very thick, mark=*, mark size=1pt] coordinates {
(1990,1.5) (1991,0.5) (1992,1.9) (1993,-3.0) (1994,12.7) (1995,25.8)
(1996,6.4) (1997,4.5) (1998,2.8) (1999,-1.2) (2000,0.8) (2001,2.8)
(2002,6.3) (2003,0.6) (2004,0.3) (2005,2.0) (2006,5.1) (2007,1.1)
(2008,5.3) (2009,3.0) (2010,1.3) (2011,2.9) (2012,2.4) (2013,2.1)
(2014,1.9) (2015,2.7) (2016,0.9) (2017,0.6) (2018,1.1) (2019,2.5)
(2020,2.5) (2021,2.3) (2022,6.3) (2023,7.4)
};
\addlegendentry{Cameroon: indicative inflation path}
\addplot[poporange, dashed, thick, domain=1990:2023] {3};
\addlegendentry{CEMAC 3\% criterion}
\end{axis}
\end{tikzpicture}
\end{center}
\legende{Cameroon's inflation over recent decades (indicative orders of magnitude): the 1994--95 devaluation spike, two decades of relative stability, and the post-2021 imported inflation above the CEMAC 3\% ceiling.}
\end{popfigure}

\begin{exoresolu}[Essay-style question]
``Inflation in Cameroon is mainly imported; therefore domestic policies are powerless.'' Discuss.
\tcblower
\textbf{Plan of a good answer.} \emph{Step 1 --- Define:} inflation, imported inflation (via import prices and exchange rate). \emph{Step 2 --- Support the claim:} Cameroon imports fuel, food and manufactures; the 1994 devaluation and post-2021 commodity shocks show the power of external causes; monetary sovereignty is pooled in the BEAC, limiting national action. \emph{Step 3 --- Nuance:} structural bottlenecks (roads, energy, storage) amplify imported shocks; domestic demand pressures and fiscal deficits also matter; supply-side policies (second-generation agriculture, local transformation of cocoa and timber) reduce import dependence over time; the BEAC's interest-rate policy affects domestic credit. \emph{Step 4 --- Conclude:} imported inflation is real, but a combination of regional monetary policy, prudent fiscal policy and supply-side reforms gives Cameroon genuine, if partial, instruments.
\end{exoresolu}

\begin{exercice}[Practice]
\begin{enumerate}
\item The CPI rises from 125 to 132.5. Calculate the inflation rate.
\item Using an AD/AS diagram, explain why a rise in world oil prices causes cost-push inflation in Cameroon, and state the effect on output.
\item Distinguish carefully, with a numerical example, between disinflation and deflation.
\item Explain why the long-run Phillips curve is vertical. What does this imply for a government tempted to reduce unemployment below its natural rate?
\end{enumerate}
\end{exercice}

\begin{corrige}[Exercise]
\begin{enumerate}
\item $\pi = \dfrac{132.5 - 125}{125}\times 100 = 6\%$.
\item Oil is an imported input; higher oil prices raise production and transport costs, shifting SRAS left. The price level rises while output falls (risk of stagflation).
\item Disinflation: inflation falls from 6\% to 3\% --- prices still rise. Deflation: inflation is $-2\%$ --- the price level falls.
\item Once expectations adjust, wage demands catch up with inflation, so unemployment returns to the natural rate at any inflation rate. Persistently pushing unemployment below $U^{*}$ only produces accelerating inflation, not lasting jobs.
\end{enumerate}
\end{corrige}

\begin{aretenir}[Key points of Section 7]
\begin{itemize}
\item Inflation is a sustained rise in the general price level, measured by the CPI: $\pi = \frac{\mathrm{CPI}_2-\mathrm{CPI}_1}{\mathrm{CPI}_1}\times 100$.
\item Causes: demand-pull (AD shifts right), cost-push (SRAS shifts left --- wages, imports, profits), monetary ($MV=PY$) and structural bottlenecks.
\item Inflation erodes purchasing power, penalises savers and lenders, redistributes income and harms competitiveness; moderate inflation can be useful.
\item The Phillips curve shows a short-run inflation--unemployment trade-off, but the long-run curve is vertical at the natural rate.
\item Deflation (falling price level) differs from disinflation (falling inflation rate); deflation feeds a dangerous spiral of deferred spending and rising real debt.
\item In Cameroon, inflation is often imported (1994 devaluation, post-2021 shocks) and constrained by the CEMAC 3\% convergence criterion, with monetary policy conducted by the BEAC.
\end{itemize}
\end{aretenir}

\coursec{Section 8: Privatisation, Nationalisation and Regional Policy}

The analysis of inflation and deflation in Section 7 showed how macro‑economic stability influences the environment in which firms operate.  When the State decides to change the ownership or the spatial distribution of productive assets, it uses two powerful — and often controversial — tools: \cle{privatisation} and \cle{nationalisation}.  A third tool, \cle{regional policy}, tries to correct the geographical imbalances that both ownership changes and market forces can create.  This section examines each tool, its theoretical justification, its practical forms, and the Cameroonian experience, covering TU 73, TU 74, TU 75, TU 76 and TU 77 of the official syllabus.

\courssub{Privatisation: Meaning, Forms and Arguments}

\begin{definition}[Privatisation]
\cle{Privatisation} is the transfer of ownership, control or management of public‑sector assets to the private sector.  It can be total (outright sale) or partial (sale of a minority stake), and it may involve a change in legal status without a change in ownership (e.g. corporatisation).
\end{definition}

\begin{aretenir}[Main Forms of Privatisation]
\begin{itemize}
    \item \textbf{Asset sale / divestiture:} Government sells shares or whole enterprises (e.g. SONEL concessioned to AES-SONEL, later ENEO).
    \item \textbf{Contracting out / outsourcing:} Public services are delivered by private firms under contract (e.g. waste collection in Yaoundé and Douala).
    \item \textbf{Deregulation / liberalisation:} Legal barriers to entry are removed, allowing private firms to compete with former monopolies (e.g. telecommunications after the 1998 law).
    \item \textbf{Public–Private Partnerships (PPP):} Long‑term contracts where the private partner finances, builds and operates infrastructure (e.g. Kribi deep‑sea port concession to KCT).
    \item \textbf{Corporatisation:} Transforming a state department into a commercial company under public law, as a first step toward privatisation (e.g. CAMTEL before mobile licences).
\end{itemize}
\end{aretenir}

\begin{propriete}[Arguments For Privatisation]
\begin{itemize}
    \item \textbf{Allocative and productive efficiency:} Private firms face hard budget constraints and profit incentives, reducing X‑inefficiency (organisational slack).
    \item \textbf{Reduced fiscal burden:} The State no longer covers recurring losses; sale proceeds can reduce public debt or finance infrastructure.
    \item \textbf{Competition and innovation:} Entry of rivals forces better products, lower prices and technological adoption.
    \item \textbf{Access to capital markets:} Private firms can raise equity and debt more easily than state‑owned enterprises (SOEs) constrained by sovereign ceilings.
    \item \textbf{Depoliticisation of management:} Decisions based on commercial criteria rather than political patronage.
\end{itemize}
\end{propriete}

\begin{propriete}[Arguments Against Privatisation]
\begin{itemize}
    \item \textbf{Natural monopoly problem:} In industries with huge fixed costs (water, electricity transmission), a single private firm may exploit monopoly power.
    \item \textbf{Job losses:} Restructuring often leads to retrenchment; social costs can be high, especially where safety nets are weak.
    \item \textbf{Higher prices for essential services:} Without effective regulation, tariffs may rise above socially optimal levels, hurting the poor.
    \item \textbf{Inequality:} Assets are often bought by wealthy insiders or foreign investors, widening wealth gaps.
    \item \textbf{Loss of strategic control:} The State may lose leverage over sectors vital for national security (energy, defence, food security).
    \item \textbf{Regulatory capture risk:} Weak regulators may be dominated by the private concessionaire.
\end{itemize}
\end{propriete}

\begin{popfigure}
\begin{tikzpicture}[scale=0.9, every node/.style={font=\small}]
    % Pivot
    \draw[thick, popdark] (0,0) -- (0,2.5) node[above] {Policy Decision};
    % Beam
    \draw[thick, popdark] (-4,2.5) -- (4,2.5);
    % Left pan (For)
    \draw[thick, popblue] (-4,2.5) -- (-4,1) -- (-3,0.5) -- (-4,1);
    \node[popblue, align=center, text width=3.2cm] at (-4,3.3) {\textbf{Arguments FOR}\\\textit{Efficiency, Fiscal relief,\\Competition, Investment,\\Depoliticisation}};
    % Right pan (Against)
    \draw[thick, poporange] (4,2.5) -- (4,1) -- (3,0.5) -- (4,1);
    \node[poporange, align=center, text width=3.2cm] at (4,3.3) {\textbf{Arguments AGAINST}\\\textit{Monopoly power, Job losses,\\Higher prices, Inequality,\\Strategic loss}};
    % Weights on left
    \foreach \y/\txt in {1.3/Efficiency, 0.9/Fiscal relief, 0.5/Competition, 0.1/Investment, -0.3/Depoliticisation}
        \node[popblue, anchor=east] at (-4.4,\y) {\txt};
    % Weights on right
    \foreach \y/\txt in {1.3/Monopoly, 0.9/Job losses, 0.5/Higher prices, 0.1/Inequality, -0.3/Strategic loss}
        \node[poporange, anchor=west] at (4.4,\y) {\txt};
    % Scale indicator
    \draw[popgreen, thick, ->] (0,2.5) -- (0,3.3) node[above] {Net welfare effect?};
\end{tikzpicture}
\legende{Balance‑scale diagram: weighing the main arguments for and against privatisation.  The net welfare effect depends on the regulatory framework, market contestability and complementary social policies.}
\end{popfigure}

\begin{methode}[Evaluating a Privatisation Case — Stakeholder Analysis]
\begin{enumerate}
    \item \textbf{Identify stakeholders:} Government (fiscal, political), Workers (jobs, wages, conditions), Consumers (price, quality, access), Investors (return, risk), Regulators (independence, capacity), Local communities (environment, spillovers).
    \item \textbf{List expected changes for each:} Ownership structure, pricing rule, investment plan, employment level, service coverage, universal service obligations.
    \item \textbf{Assess net benefit/cost:} Use qualitative scores ($+$, $0$, $-$) or quantitative estimates (NPV of fiscal savings, consumer surplus change, employment effects).
    \item \textbf{Check for complementary reforms:} Is there an independent regulator with price‑cap or revenue‑cap powers? Is there a competition law? Are there social safety nets for displaced workers? Are universal service obligations legally enforced?
    \item \textbf{Conclude:} Does the privatisation increase total welfare?  If not, what design changes (e.g. price caps, employee share ownership, phased liberalisation) could improve it?
\end{enumerate}
\end{methode}

\begin{exoresolu}[Case Study: Privatisation of SONEL (Cameroon)]
\textbf{Statement:} In 2001, the Cameroon electricity utility SONEL was concessioned to a consortium (Actis, EDF) creating AES-SONEL (later ENEO).  Evaluate the outcome using the stakeholder method.

\textbf{Solution:}
\begin{itemize}
    \item \textbf{Government:} Received concession fees; reduced direct subsidies.  However, contingent liabilities (e.g. guaranteed tariffs, power purchase agreements) remained.  \textit{Score: $+$ (fiscal relief) / $-$ (contingent risk)}.
    \item \textbf{Workers:} Significant job cuts in the first years (estimates vary); remaining staff saw wage alignment with private sector norms.  \textit{Score: $-$ (short term) / $0$ (long term)}.
    \item \textbf{Consumers:} Tariffs moved toward cost‑reflective levels; supply reliability improved in major urban areas but rural electrification lagged behind targets.  \textit{Score: $-$ (price) / $+$ (quality urban) / $-$ (rural access)}.
    \item \textbf{Investors:} Earned returns via regulated asset base; invested in new generation capacity (e.g. Kribi gas-fired plant, Memve'ele hydro).  \textit{Score: $+$}.
    \item \textbf{Regulator (ARSEL, created 1998):} Legal framework exists but capacity constraints limited effective price‑cap enforcement and monitoring of investment obligations.  \textit{Score: $0$}.
    \item \textbf{Local communities/Environment:} New generation projects brought local employment but also environmental and social impact concerns (resettlement, water use).  \textit{Score: $0$ / $-$}.
\end{itemize}
\textbf{Overall:} Net welfare gain in urban centres; persistent rural access gap and affordability concerns for low‑income households.  Lesson: privatisation without strong regulation, universal service obligations and social mitigation can leave the poor behind.
\tcblower
\textbf{Detailed solution:} The stakeholder table above summarises the evaluation.  The key insight for GCE essays is that \emph{ownership change alone does not guarantee competition}; the electricity transmission and distribution network remains a natural monopoly, so a private concessionaire needs a robust price‑cap or revenue‑cap regime, clear quality standards and enforceable universal service obligations.  Cameroon's experience illustrates the \cle{warning} below.
\end{exoresolu}

\begin{attention}[Common Mistake]
\textbf{Assuming privatisation automatically creates competition.}  Changing ownership from public to private does \emph{not} alter the market structure.  If the industry is a natural monopoly (water distribution, electricity transmission, railway infrastructure), a private monopolist can be even more rent‑seeking than a public one unless an independent regulator imposes price caps, quality standards and universal service obligations.  Always ask: \emph{Is there effective competition or a credible regulatory contract?}
\end{attention}

\courssub{Nationalisation: Meaning, Reasons and Criticisms}

\begin{definition}[Nationalisation]
\cle{Nationalisation} is the acquisition by the State of privately owned assets, usually with compensation, bringing them under public ownership and control.  It can be total (expropriation) or partial (golden share, strategic stake).
\end{definition}

\begin{aretenir}[Main Reasons for Nationalisation]
\begin{itemize}
    \item \textbf{Strategic control:} Sectors deemed vital for sovereignty (oil refining — SONARA, defence, telecommunications backbone, food security).
    \item \textbf{Natural monopoly:} To avoid private monopoly pricing, the State may keep or reclaim the network infrastructure (e.g. CAMTEL's fixed‑line infrastructure).
    \item \textbf{Equity and universal access:} Ensuring affordable services for low‑income groups (water, electricity, transport).
    \item \textbf{Employment protection:} Preventing mass lay‑offs in politically sensitive firms or regions (e.g. debates around CDC, SOCAPALM).
    \item \textbf{Rescue of failing firms:} Temporary nationalisation to restructure and later reprivatise (bank bailouts during financial crises).
    \item \textbf{Resource nationalism:} Capturing a larger share of rents from natural resources (oil, minerals).
\end{itemize}
\end{aretenir}

\begin{propriete}[Arguments Against Nationalisation]
\begin{itemize}
    \item \textbf{X‑inefficiency:} Absence of profit motive and soft budget constraints lead to over‑staffing, under‑investment and lack of innovation.
    \item \textbf{Political interference:} Appointments based on loyalty rather than merit; pricing and investment decisions driven by electoral cycles.
    \item \textbf{Fiscal burden:} Losses are absorbed by the treasury, crowding out other public spending (health, education, infrastructure).
    \item \textbf{Corruption and rent‑seeking:} Procurement and hiring become patronage channels; lack of transparency.
    \item \textbf{Crowding out private investment:} Uncertainty about property rights deters domestic and foreign private capital.
\end{itemize}
\end{propriete}

\begin{aretenir}[Privatisation vs Nationalisation — Comparative Table]
\begin{center}
\begin{tabularx}{\linewidth}{|l|X|X|}
\hline
\rowcolor{popblueL}
\textbf{Dimension} & \textbf{Privatisation} & \textbf{Nationalisation} \\ \hline
\textbf{Primary aim} & Efficiency, fiscal relief, competition & Strategic control, equity, universal access \\ \hline
\textbf{Typical sectors} & Competitive manufacturing, telecoms, hotels, airlines & Utilities, natural resources, defence, strategic infrastructure \\ \hline
\textbf{Advantages} & Hard budget constraint, access to capital, innovation, depoliticisation & Price control for essential goods, job preservation, national sovereignty, resource rent capture \\ \hline
\textbf{Disadvantages} & Possible private monopoly, job losses, inequality, regulatory capture & X‑inefficiency, political capture, fiscal drain, corruption, crowding out \\ \hline
\textbf{Cameroon examples} & SONEL $\rightarrow$ AES-SONEL/ENEO, CAMTEL mobile licences (MTN, Orange), CAMAIR $\rightarrow$ Camair‑Co, hotel privatisations & SONARA (state‑owned refinery), CDC/SOCAPALM (state majority), CAMTEL fixed assets (re‑nationalised), Camair‑Co (state-owned airline) \\ \hline
\end{tabularx}
\end{center}
\end{aretenir}

\begin{savaistu}[Nationalisation in the CEMAC Context]
The CEMAC Treaty (1994) encourages free movement of capital but allows Member States to retain strategic enterprises.  In practice, Cameroon has oscillated: the 1990s Structural Adjustment Programme (SAP) pushed privatisation (SONEL, CAMTEL mobile, CAMAIR, hotels), while the 2010s saw a return to state‑led investment (Kribi port, Lom Pangar dam, Memve'ele hydro, creation of \emph{Société Nationale de Raffinage} — SONARA remains 100\% state).  The debate mirrors the global tension between \cle{market‑friendly} reforms and \cle{developmental state} ambitions.  Recent CEMAC directives emphasise better governance of SOEs rather than blanket privatisation.
\end{savaistu}

\courssub{Regional Policy: Objectives, Instruments and Application in Cameroon}

\begin{definition}[Regional Policy]
\cle{Regional policy} comprises government interventions designed to reduce economic disparities between regions and to promote balanced national development.  It targets \cle{spatial inequalities} in income, employment, infrastructure and access to public services.
\end{definition}

\begin{definition}[Natural Monopoly (recall)]
A \cle{natural monopoly} exists when a single firm can supply the entire market at lower average cost than multiple firms, due to massive economies of scale (e.g. electricity transmission, water pipes, railway tracks).  Regional policy often addresses the location of such infrastructure.
\end{definition}

\begin{aretenir}[Objectives of Regional Policy]
\begin{itemize}
    \item Reduce \cle{inter‑regional disparities} (GDP per capita, poverty rates, human development indicators).
    \item Prevent excessive \cle{rural‑urban migration} and congestion in primate cities (Douala, Yaoundé).
    \item Exploit comparative advantages of each region (agro‑ecological zones, mineral resources, ports, tourism potential).
    \item Strengthen national cohesion and political stability.
    \item Promote \cle{territorial equity} in access to public services (health, education, administration).
\end{itemize}
\end{aretenir}

\begin{propriete}[Main Instruments of Regional Policy]
\begin{itemize}
    \item \textbf{Infrastructure investment:} Roads, railways, energy, digital connectivity (e.g. Yaoundé‑Nsimalen highway, Kribi‑Lolabé railway, fibre optic backbone).
    \item \textbf{Fiscal and financial incentives:} Tax holidays, reduced customs duties, subsidised credit for firms locating in lagging regions (e.g. Law on Investment Incentives, special regimes for industrial free zones).
    \item \textbf{Industrial location policies:} Creation of \cle{development poles} (agropoles, industrial free zones, special economic zones) with ready‑to‑use plots and utilities.
    \item \textbf{Decentralisation and local governance:} Transfer of competences and resources to Regional Councils (Law No. 2019/024) to design bottom‑up Regional Development Plans.
    \item \textbf{Human capital investment:} Regional universities, vocational training centres, health facilities (e.g. University of Maroua, University of Bertoua, University of Ngaoundéré).
    \item \textbf{Public service relocation:} Moving administrative headquarters or agencies to secondary cities.
\end{itemize}
\end{propriete}

\begin{popfigure}
\begin{tikzpicture}[scale=0.85, every node/.style={font=\footnotesize}]
    % Simplified Cameroon outline (rectangle approx)
    \draw[popdark, thick] (0,0) rectangle (12,9);
    % Regions as labelled boxes with text width
    \node[draw, popblue, fill=popblueL, rounded corners, align=center, text width=2.2cm] at (2,7.5) {Far North\\Maroua};
    \node[draw, popblue, fill=popblueL, rounded corners, align=center, text width=2.2cm] at (5.5,7.5) {North\\Garoua};
    \node[draw, popblue, fill=popblueL, rounded corners, align=center, text width=2.2cm] at (9,7.5) {Adamawa\\Ngaoundéré};
    \node[draw, popblue, fill=popblueL, rounded corners, align=center, text width=2.2cm] at (1.5,5) {East\\Bertoua};
    \node[draw, popblue, fill=popblueL, rounded corners, align=center, text width=2.2cm] at (5,5) {Centre\\Yaoundé};
    \node[draw, popblue, fill=popblueL, rounded corners, align=center, text width=2.2cm] at (8.5,5) {West\\Bafoussam};
    \node[draw, popblue, fill=popblueL, rounded corners, align=center, text width=2.2cm] at (2,2.5) {South\\Ebolowa};
    \node[draw, popblue, fill=popblueL, rounded corners, align=center, text width=2.2cm] at (5.5,2.5) {Littoral\\Douala};
    \node[draw, popblue, fill=popblueL, rounded corners, align=center, text width=2.2cm] at (9,2.5) {South‑West\\Buea};
    % Development poles (circles)
    \foreach \pos/\name in {
        (5.5,2.5)/Douala,
        (5,5)/Yaoundé,
        (9,7.5)/Ngaoundéré,
        (2,7.5)/Maroua,
        (1.5,5)/Bertoua,
        (8.5,5)/Bafoussam,
        (2,2.5)/Ebolowa,
        (9,2.5)/Buea
    } {
        \node[circle, draw=poporange, fill=poporange!20, thick, minimum size=6mm] at \pos {};
        \node[poporange, anchor=north, text width=1.5cm, align=center] at \pos {\name};
    }
    % Infrastructure arrows (sample)
    \draw[popgreen, thick, ->] (5.5,2.5) -- (5.5,5) node[midway, right, text width=2cm, align=center] {Road/Rail\\corridor};
    \draw[popgreen, thick, ->] (5,5) -- (9,7.5) node[midway, above, text width=2cm, align=center] {Rail to\\Ngaoundéré};
    \draw[popgreen, thick, ->] (5.5,2.5) -- (2,2.5) node[midway, below, text width=2cm, align=center] {Coastal road};
    \draw[popgreen, thick, ->] (5.5,2.5) -- (9,2.5) node[midway, below, text width=2cm, align=center] {Kribi port link};
    % Agropoles labels
    \node[poppurple, align=left, text width=3cm] at (0.5,8.5) {Agropoles / Development Poles:};
    \node[poppurple, align=left, text width=3cm] at (0.5,8.0) {Maroua: cotton, cereals};
    \node[poppurple, align=left, text width=3cm] at (0.5,7.5) {Ngaoundéré: livestock, maize};
    \node[poppurple, align=left, text width=3cm] at (0.5,7.0) {Bertoua: timber, mining};
    \node[poppurple, align=left, text width=3cm] at (0.5,6.5) {Bafoussam: coffee, cocoa};
    \node[poppurple, align=left, text width=3cm] at (0.5,6.0) {Kribi: deep‑sea port, industry};
    \node[poppurple, align=left, text width=3cm] at (0.5,5.5) {Limbe: oil refinery (SONARA)};
\end{tikzpicture}
\legende{Simplified schematic of Cameroon's ten regions, major urban poles (orange circles) and key infrastructure corridors (green arrows).  Purple labels indicate agropoles / development poles promoted by regional policy.  Not to scale.}
\end{popfigure}

\begin{experience}[Guided Work 8 — Privatisation in Cameroon (Recap)]
Using the stakeholder method, compare the privatisation of \textbf{CAMTEL} (mobile licences to MTN/Orange, fixed‑line retained) with the \textbf{re‑nationalisation of CAMTEL's fixed assets} (circa 2016).  Discuss: (a) fiscal impact, (b) service quality and coverage, (c) competition in mobile vs fixed markets, (d) lessons for future PPPs and universal service provision.
\end{experience}

\begin{exoresolu}[Worked Exercise: Evaluating the Kribi Deep‑Sea Port as a Regional Policy Instrument]
\textbf{Statement:} The Kribi port (Phase 1 operational 2014, Phase 2 under construction) is a PPP concession to \emph{Kribi Conteneur Terminal (KCT)}.  Assess its contribution to regional policy objectives for the South Region and the national economy.

\textbf{Solution:}
\begin{itemize}
    \item \textbf{Infrastructure spillover:} The port triggered the Kribi‑Lolabé railway (mineral evacuation) and the Yaoundé‑Kribi highway, linking the South Region to the centre and reducing transport costs.
    \item \textbf{Employment:} Direct jobs estimated at several thousand; indirect (logistics, hospitality, construction) significantly higher.  Local hiring quotas included in concession agreement.
    \item \textbf{Fiscal revenue:} Concession fees + customs revenue from increased throughput (target 1.4M TEU by 2035).  Potential for transit fees from landlocked neighbours (Chad, CAR).
    \item \textbf{Regional equity:} South Region historically lagged (poverty rate above national average).  Port investment raises its GDP share; however, benefits concentrate in Kribi town — rural hinterland needs complementary agropoles (palm oil, rubber, cocoa) and skills training.
    \item \textbf{Risks:} Concession contract may include revenue guarantees to KCT (contingent liability for the State).  Environmental concerns (mangroves, biodiversity, coastal erosion).  Social tensions over land compensation.
    \item \textbf{Strategic positioning:} Positions Cameroon as a gateway for Central Africa, supporting AfCFTA objectives.
\end{itemize}
\textbf{Conclusion:} Kribi port is a \cle{growth pole} that aligns with regional policy if accompanied by \emph{inclusive} local development plans (skills training, SME support, environmental safeguards, revenue sharing with local communities).
\tcblower
\textbf{Detailed solution:} The answer above follows the evaluation framework: identify objectives, list instruments, trace transmission channels (infrastructure $\rightarrow$ market access $\rightarrow$ firm location $\rightarrow$ jobs), and flag risks (fiscal contingencies, spatial concentration, environmental externalities).  In GCE essays, always link the \emph{instrument} to the \emph{objective} and provide \emph{evidence} (even if qualitative).
\end{exoresolu}

\begin{aretenir}[Key Takeaways for Section 8]
\begin{itemize}
    \item Privatisation transfers ownership to the private sector; forms include asset sales, contracting out, deregulation, PPPs, corporatisation.
    \item Theoretical case for privatisation rests on efficiency, fiscal relief, competition, investment access and depoliticisation; case against stresses natural monopoly, job losses, inequality, strategic control loss and regulatory capture.
    \item \textbf{Warning:} Ownership change $\neq$ competition.  Effective regulation (independent, price‑cap/quality‑cap, universal service obligations) is essential, especially for natural monopolies.
    \item Nationalisation brings assets under state control for strategic, equity, universal access, employment, rescue or resource nationalism reasons; critics highlight X‑inefficiency, political capture, fiscal cost, corruption and crowding out.
    \item Regional policy aims to reduce spatial disparities using infrastructure, fiscal incentives, development poles (agropoles, SEZs), decentralisation, human capital investment and public service relocation.
    \item Cameroon's experience illustrates both tools: 1990s privatisation wave (SONEL, CAMTEL mobile, CAMAIR, hotels), recent state‑led investments (Kribi, Lom Pangar, Memve'ele, SONARA), and the 2019 decentralisation law empowering Regional Councils.
    \item Evaluation of any policy requires a \emph{stakeholder analysis} (government, workers, consumers, investors, regulators, communities) and a check on complementary reforms (regulation, competition law, social safety nets, universal service obligations, environmental safeguards).
\end{itemize}
\end{aretenir}

\begin{exercice}[Practice Questions]
\begin{enumerate}
    \item Distinguish between privatisation and deregulation.  Give a Cameroonian example of each.
    \item Explain why privatising a natural monopoly without an independent regulator may worsen consumer welfare.
    \item List three instruments of regional policy and illustrate how they have been applied in Cameroon since 2010.
    \item "Nationalisation is always inefficient."  Discuss with reference to the Cameroonian context.
    \item Using the stakeholder method, outline the likely effects of a hypothetical full privatisation of SONARA.
\end{enumerate}
\end{exercice}

\begin{corrige}[Exercise 1]
Privatisation is a change of ownership (public $\rightarrow$ private); deregulation is a removal of legal barriers to entry without necessarily changing ownership.  Example: SONEL $\rightarrow$ AES-SONEL/ENEO (privatisation via concession); 1998 telecommunications law allowing MTN/Orange to compete with CAMTEL (deregulation/liberalisation).
\end{corrige}

\begin{corrige}[Exercise 2]
A natural monopoly has declining average costs over the relevant output range due to massive economies of scale.  A private monopolist maximises profit where $MR=MC$, setting $P>MC$ and restricting output below the socially optimal level.  Without a regulator imposing $P=AC$ (average cost pricing) or a price cap ($RPI-X$), the private firm charges a higher price than a welfare‑maximising public firm would, reducing consumer surplus and creating deadweight loss.  The public firm, in theory, could price at $MC$ (with a subsidy) or $AC$ (break-even), achieving allocative or productive efficiency respectively.
\end{corrige}

\begin{corrige}[Exercise 3]
(1) Infrastructure: Kribi‑Lolabé railway, Yaoundé‑Kribi highway, fibre optic backbone extension to northern regions. (2) Fiscal incentives: Tax holidays and customs duty exemptions for firms in the Kribi Industrial Zone and other special economic zones under the 2013 Investment Incentives Law. (3) Development poles: Agropoles in Maroua (cotton, cereals), Ngaoundéré (livestock, maize), Bertoua (timber, mining), Bafoussam (coffee, cocoa). (4) Decentralisation: 2019 Regional Councils (Law No. 2019/024) receive own‑source revenues and design Regional Development Plans.
\end{corrige}

\begin{corrige}[Exercise 4]
Not always.  Nationalisation can correct market failures (natural monopoly, positive externalities, coordination failures) and protect strategic interests (energy security, food sovereignty).  In Cameroon, SONARA ensures domestic fuel supply security and captures refining margins; CDC/SOCAPALM maintain employment and social stability in plantation zones.  However, persistent losses, political appointments, under‑investment and governance challenges illustrate the efficiency risks.  The net outcome depends on governance quality, transparency, and whether the SOE is run on commercial principles with hard budget constraints.
\end{corrige}

\begin{corrige}[Exercise 5]
Stakeholders: Government (fiscal windfall from sale vs loss of strategic control over fuel supply and fiscal revenue from refinery margins), Workers (possible retrenchment, wage harmonisation, loss of public-sector benefits), Consumers (price liberalisation $\rightarrow$ higher pump prices reflecting import parity, but maybe better supply reliability and product quality), Investors (access to a monopoly asset with guaranteed demand, need stable fiscal terms and regulatory certainty), Regulator (capacity to monitor pricing, environmental standards, competition if market opened), Local communities (environmental impact, CSR obligations).  Net effect hinges on whether a credible regulatory contract (price cap, quality standards, universal supply obligation, environmental compliance) replaces state ownership as the discipline device, and whether competition can be introduced in downstream distribution.
\end{corrige}

\coursec{Section 9: Population — Growing and Ageing Populations, Optimum Population and Population Policy}

In Section 8 we examined how the State shapes the economy through privatisation, nationalisation and regional policy. We now turn to the most fundamental economic resource of all: \cle{people}. Every worker, every consumer, every taxpayer is a member of the population. A country like Cameroon, where roughly two-thirds of the population is under 25, faces very different economic challenges from a country like Japan, where more than a quarter of the population is over 65. Understanding how populations grow, age and interact with resources is therefore essential for judging whether population change is a blessing or a burden — and for designing the policies that make it a blessing.

\courssub{Population Growth: Determinants and Measurement}

\begin{definition}[Population growth]
\cle{Population growth} is the change in the size of a country's population over a given period, usually expressed as a percentage per year. It depends on three determinants:
\begin{itemize}
\item the \cle{birth rate} (number of live births per 1\,000 inhabitants per year);
\item the \cle{death rate} (number of deaths per 1\,000 inhabitants per year);
\item \cle{net migration} (immigrants minus emigrants per 1\,000 inhabitants).
\end{itemize}
\end{definition}

\begin{propriete}[Natural growth rate and total growth rate]
The \cle{natural growth rate} of a population is:
\[ \text{Natural growth rate} = \text{Birth rate} - \text{Death rate} \quad (\text{per 1\,000}) \]
The \cle{total population growth rate} adds net migration:
\[ \text{Total growth rate} = (\text{Birth rate} - \text{Death rate}) + \text{Net migration} \]
To convert a rate per 1\,000 into a percentage, divide by 10. (The formula itself is examinable in data-response questions.)
\end{propriete}

\begin{exemplebox}[Calculating Cameroon's natural growth]
Suppose a country records a birth rate of 36 per 1\,000 and a death rate of 10 per 1\,000, with negligible net migration. The natural growth rate is $36 - 10 = 26$ per 1\,000, that is $2.6\%$ per year. At this pace, the population doubles in roughly $\frac{70}{2.6} \approx 27$ years (the "rule of 70"). This is the situation of many Sub-Saharan African countries, including Cameroon.
\end{exemplebox}

\courssub{The Demographic Transition Model}

How do birth and death rates evolve as a country develops? The \cle{demographic transition model} (DTM) describes the historical experience of today's rich countries and provides a framework for analysing countries like Cameroon.

\begin{propriete}[Stages of the demographic transition model]
\begin{enumerate}
\item \textbf{Stage 1 — High stationary:} high birth rate, high and fluctuating death rate; population is small and grows slowly (pre-industrial societies).
\item \textbf{Stage 2 — Early expanding:} death rate falls sharply (better hygiene, food, medicine) while the birth rate stays high; population grows rapidly.
\item \textbf{Stage 3 — Late expanding:} birth rate falls (urbanisation, education of girls, family planning); population still grows but more slowly.
\item \textbf{Stage 4 — Low stationary:} low birth rate, low death rate; population is large but stable.
\item \textbf{Stage 5 (sometimes added):} birth rate falls below the death rate; population declines and ages (Japan, parts of Europe).
\end{enumerate}
\end{propriete}

\begin{popfigure}
\begin{center}
\begin{tikzpicture}
\begin{axis}[
width=12cm, height=7cm,
xlabel={Time / stage of development},
ylabel={Rate per 1\,000},
xmin=0, xmax=4, ymin=0, ymax=45,
xtick={0.5,1.5,2.5,3.5},
xticklabels={Stage 1,Stage 2,Stage 3,Stage 4},
ytick=\empty,
legend style={at={(0.02,0.97)}, anchor=north west, font=\small},
axis lines=left, thick]
\addplot[very thick, popblue, smooth] coordinates {(0,38) (0.8,37) (1.4,34) (2,25) (2.6,15) (3.2,11) (4,10)};
\addlegendentry{Birth rate}
\addplot[very thick, poporange, smooth] coordinates {(0,36) (0.7,35) (1.2,25) (1.7,15) (2.3,11) (3,10) (4,10)};
\addlegendentry{Death rate}
\addplot[very thick, popgreen, smooth] coordinates {(0,3) (1,4) (1.8,9) (2.6,16) (3.2,20) (4,22)};
\addlegendentry{Total population (scaled)}
\end{axis}
\end{tikzpicture}
\end{center}
\legende{The demographic transition model: the gap between birth and death rates drives population growth.}
\end{popfigure}

Cameroon is best described as being in \textbf{late Stage 2 / early Stage 3}: death rates have fallen considerably thanks to vaccination and better health care, while birth rates remain high, producing rapid population growth and a very youthful age structure.

\courssub{Economic Effects of a Growing Population}

\textbf{Positive effects.}
\begin{itemize}
\item \textbf{Larger labour supply:} more workers can raise productive capacity, provided jobs are created.
\item \textbf{Larger domestic market:} a bigger population means more consumers, which can attract investment and allow economies of scale.
\item \textbf{Demographic dividend:} when fertility falls, the share of the working-age population rises relative to dependants, creating a window of accelerated growth — if education and employment follow.
\end{itemize}

\textbf{Negative effects.}
\begin{itemize}
\item \textbf{Pressure on resources:} land, water and food must be shared among more people.
\item \textbf{Dependency burden:} a youthful population means many children per worker, raising consumption needs relative to output.
\item \textbf{Unemployment and underemployment:} if job creation lags behind labour-force growth, urban unemployment and the informal sector swell — a familiar challenge in Douala and Yaoundé.
\item \textbf{Strain on services:} schools, hospitals and housing struggle to keep pace.
\end{itemize}

\begin{definition}[Dependency ratio and demographic dividend]
The \cle{dependency ratio} measures the economic burden carried by workers:
\[ \text{Dependency ratio} = \frac{\text{Population aged 0--14} + \text{Population aged 65+}}{\text{Population aged 15--64}} \times 100 \]
The \cle{demographic dividend} is the potential boost to economic growth that occurs when the working-age share of the population rises and the dependency ratio falls, provided workers are healthy, educated and employed.
\end{definition}

\courssub{The Ageing Population}

An \cle{ageing population} is one in which the proportion of elderly people rises over time. Its causes are \textbf{falling fertility} (fewer children per woman) and \textbf{rising life expectancy} (better medicine and living standards). It is measured by the \cle{old-age dependency ratio}: population aged 65+ divided by population aged 15--64.

\textbf{Economic consequences:}
\begin{itemize}
\item \textbf{Pension and health-care costs:} fewer workers must finance more pensions and medical care, straining public budgets.
\item \textbf{Labour shortages:} shrinking workforces push up wages and may slow growth; some countries respond with automation or immigration.
\item \textbf{Changing demand patterns:} demand shifts towards health services, leisure for retirees and adapted housing, and away from schools and baby products.
\item \textbf{Savings and investment effects:} older societies may save less as retirees draw down their assets.
\end{itemize}

\begin{popfigure}
\begin{center}
\begin{tikzpicture}[scale=0.62]
% Youthful pyramid (Cameroon)
\begin{scope}
\foreach \y/\w in {0/2.4, 1/2.2, 2/2.0, 3/1.8, 4/1.5, 5/1.2, 6/0.9, 7/0.6, 8/0.3}{
\fill[popblue] (-\w,\y) rectangle (0,\y+0.85);
\fill[poppink] (0,\y) rectangle (\w,\y+0.85);}
\node[font=\small\bfseries] at (0,9.3) {Cameroon (youthful)};
\node[font=\scriptsize] at (-1.4,-0.6) {Male};
\node[font=\scriptsize] at (1.4,-0.6) {Female};
\end{scope}
% Ageing pyramid (Japan)
\begin{scope}[xshift=11cm]
\foreach \y/\w in {0/1.1, 1/1.2, 2/1.3, 3/1.5, 4/1.7, 5/1.8, 6/1.7, 7/1.4, 8/0.9}{
\fill[popblue] (-\w,\y) rectangle (0,\y+0.85);
\fill[poppink] (0,\y) rectangle (\w,\y+0.85);}
\node[font=\small\bfseries] at (0,9.3) {Japan (ageing)};
\node[font=\scriptsize] at (-1.4,-0.6) {Male};
\node[font=\scriptsize] at (1.4,-0.6) {Female};
\end{scope}
\end{tikzpicture}
\end{center}
\legende{Population pyramids: a broad base (Cameroon) signals youth and rapid growth; a narrow base with a heavy top (Japan) signals ageing.}
\end{popfigure}

\begin{methode}[Analysing a population pyramid in data-response questions]
\begin{enumerate}
\item \textbf{Look at the base:} a wide base means a high birth rate and a youthful population; a narrow base means falling fertility.
\item \textbf{Look at the top:} a wide top (especially among women, who live longer) indicates an ageing population.
\item \textbf{Look at the shape overall:} a triangle = rapid growth; a beehive/urn shape = slow growth or decline.
\item \textbf{Identify the working-age share} (roughly ages 15--64) and comment on the dependency ratio.
\item \textbf{Draw the economic conclusion:} youthful $\Rightarrow$ need for schools, jobs, family planning; ageing $\Rightarrow$ need for pensions, health care, labour supply policies.
\end{enumerate}
\end{methode}

\courssub{Optimum Population}

Is there an ideal population size? Economists answer with the concept of \cle{optimum population}.

\begin{definition}[Optimum, over- and under-population]
The \cle{optimum population} is the population size which, given a country's resources, technology and capital, maximises \textbf{output (income) per head}.
\begin{itemize}
\item \cle{Under-population}: population is too small to exploit resources fully; output per head is below its maximum (e.g. vast, sparsely populated Canada).
\item \cle{Over-population}: population is too large relative to resources; diminishing returns push output per head below its maximum.
\end{itemize}
\end{definition}

\begin{popfigure}
\begin{center}
\begin{tikzpicture}
\begin{axis}[
width=11cm, height=7cm,
xlabel={Population size}, ylabel={Output per head},
xmin=0, xmax=10, ymin=0, ymax=6,
xtick=\empty, ytick=\empty,
axis lines=left, thick]
\addplot[very thick, poppurple, smooth, domain=0.4:9.6] {5.2*exp(-((x-4.6)^2)/9)};
\addplot[dashed, popmuted] coordinates {(4.6,0) (4.6,5.2)};
\addplot[dashed, popmuted] coordinates {(0,5.2) (4.6,5.2)};
\node[font=\small, popdark] at (axis cs:4.6,-0.35) {$P^{*}$ optimum};
\node[font=\small, popdark] at (axis cs:-0.4,5.2) {$Y^{*}$};
\node[font=\small, popblue] at (axis cs:1.6,2.2) {Under-population};
\node[font=\small, poporange] at (axis cs:8,2.2) {Over-population};
\end{axis}
\end{tikzpicture}
\end{center}
\legende{The optimum population: output per head is maximised at $P^{*}$; below it resources are underused, above it diminishing returns dominate.}
\end{popfigure}

\begin{attention}[Optimum population is not a fixed number]
A frequent exam mistake is to treat the optimum population as a permanent figure. It is \textbf{dynamic}: improvements in technology, discovery of new resources (e.g. oil, better crop varieties) or higher capital stock shift the curve \emph{upwards}, raising both the optimum population and maximum output per head. Conversely, resource depletion lowers it. Always qualify your answer with "given current technology and resources".
\end{attention}

\courssub{Malthus versus the Optimists (Exam-Flagged)}

\textbf{Malthusian theory} (Thomas Malthus, late 18th century) holds that population grows \emph{geometrically} (2, 4, 8, 16...) while food supply grows only \emph{arithmetically} (1, 2, 3, 4...). Unchecked population therefore outruns food, leading to "positive checks" (famine, disease, war) unless "preventive checks" (moral restraint, later marriage) limit births. \textbf{Evaluation:} Malthus underestimated technological progress — agricultural yields have risen faster than population in most of the world — yet his warning still resonates where rapid population growth meets fragile environments.

The \textbf{optimist view} (associated with Ester Boserup) argues that population pressure \emph{stimulates} innovation: necessity is the mother of invention. More people means more ideas, larger markets and incentives to intensify agriculture. In essays, the strongest answers present both views and conclude that the outcome depends on institutions, education and technology.

\courssub{Further Studies 11 — Population Policy}

\begin{definition}[Population policy]
A \cle{population policy} is the set of measures a government uses to influence the size, growth rate, structure and distribution of its population in line with economic and social objectives.
\end{definition}

\textbf{Anti-natalist policies} (to reduce fertility, used where population grows too fast):
\begin{itemize}
\item family planning programmes and access to contraception;
\item education of girls and women (strongly correlated with lower fertility);
\item incentives such as tax advantages for smaller families, or disincentives (China's former one-child policy is the extreme case);
\item raising the legal age of marriage.
\end{itemize}

\textbf{Pro-natalist policies} (to raise fertility, used in ageing societies such as France or Japan):
\begin{itemize}
\item family allowances and child benefits;
\item paid parental leave and subsidised childcare;
\item tax relief for families with children.
\end{itemize}

\textbf{Migration policy} complements both: ageing countries attract immigrant workers to fill labour shortages, while emigration can relieve (or drain, through the "brain drain") youthful countries.

\textbf{Cameroon's population policy.} Cameroon has a \textbf{high fertility rate} and a very \textbf{youthful population}, with well over 40\% of Cameroonians under 15. Successive governments have supported family planning through the Ministry of Public Health and partners, promoted girls' education, and adopted national strategies (notably within the long-term development vision) aimed at \textbf{harnessing the demographic dividend}: investing in education, health and job creation so that the coming bulge of working-age Cameroonians becomes an engine of growth rather than a source of unemployment. The contrast with ageing economies such as Japan or much of Europe is striking: they worry about pensions and labour shortages, while Cameroon worries about schools, jobs and dependency — two opposite faces of the same demographic transition.

\begin{exoresolu}[Worked exercise — dependency ratio and policy]
A country has 8 million people aged 0--14, 12 million aged 15--64 and 2 million aged 65 and over. (a) Calculate the dependency ratio. (b) State whether the country is more likely youthful or ageing, and suggest one appropriate population policy.
\tcblower
\textbf{(a)} 
\[ \text{Dependency ratio} = \frac{8 + 2}{12} \times 100 = 83.3\% \]
There are about 83 dependants for every 100 working-age people.
\textbf{(b)} Since children (8 million) greatly outnumber the elderly (2 million), this is a \textbf{youthful} population (like Cameroon). An appropriate policy is investment in \textbf{education and family planning}, combined with job creation, to prepare for and capture the demographic dividend. (An ageing country would instead prioritise pensions and pro-natalist incentives.)
\end{exoresolu}

\begin{exercice}[Practice]
\begin{enumerate}
\item A country has a birth rate of 40 per 1\,000, a death rate of 14 per 1\,000 and net emigration of 2 per 1\,000. Calculate its total population growth rate in percentage terms.
\item Using a diagram, explain why the optimum population of a country can increase over time.
\item "Rapid population growth is always an obstacle to development." Discuss, with reference to Cameroon and one ageing economy.
\end{enumerate}
\end{exercice}

\begin{corrige}[Exercise 1]
\begin{enumerate}
\item Natural growth $= 40 - 14 = 26$ per 1\,000. Total growth $= 26 - 2 = 24$ per 1\,000 $= 2.4\%$ per year.
\item The optimum population diagram plots output per head against population. Technical progress or new resources shift the output-per-head curve upwards and to the right, so the maximum occurs at a larger population: $P^{*}$ rises. Hence the optimum is dynamic, not fixed.
\item \textbf{Against the statement:} population growth enlarges the labour force and market, and can yield a demographic dividend (Boserup's optimist view; East Asian NICs benefited). \textbf{For the statement:} where job creation, education and food production lag, rapid growth raises the dependency ratio, strains services and fuels unemployment — a real risk for Cameroon if fertility stays high without investment in human capital. \textbf{Ageing economies} (Japan, Europe) show the opposite problem: shrinking workforces and pension burdens. \textbf{Conclusion:} growth is an obstacle only when complementary investments (education, health, jobs) are missing; with the right policies, it becomes an asset.
\end{enumerate}
\end{corrige}

\begin{aretenir}[Key points of Section 9]
\begin{itemize}
\item Population growth = (birth rate $-$ death rate) + net migration; the natural growth rate alone excludes migration.
\item The demographic transition model has four (or five) stages; Cameroon is in the expanding stages, Japan in the final stage.
\item The dependency ratio $= \frac{(0\text{--}14) + (65+)}{15\text{--}64} \times 100$; a falling ratio creates the chance of a demographic dividend.
\item Optimum population maximises output per head and \textbf{shifts with technology and resources} — it is never a fixed number.
\item Malthus feared population would outrun food; optimists (Boserup) reply that population pressure drives innovation.
\item Population policies may be anti-natalist (family planning, girls' education) or pro-natalist (allowances, parental leave); Cameroon seeks to harness its youthful population, while ageing economies fight labour shortages.
\end{itemize}
\end{aretenir}

\newpage

\coursec{Integration activity}

\coursec{Integration Activity: Unlocking Structural Transformation — Policy Brief Simulation}

\courssub{Aims of the Activity}
This integration activity is designed to consolidate your mastery of the entire Upper Sixth Economics syllabus (Module ECO08) by placing you in a realistic policy-making scenario. By the end of this activity, you should be able to:
\begin{itemize}
    \item \textbf{Synthesise} macroeconomic data to diagnose the phase of the \cle{trade (business) cycle} and identify binding constraints on \cle{economic development}.
    \item \textbf{Apply} concepts of \cle{structural unemployment}, \cle{imported inflation}, \cle{demographic dividend}, and \cle{deindustrialisation} to the Cameroonian context.
    \item \textbf{Evaluate} the relevance of \cle{NIC development strategies} (export-oriented industrialisation, human capital investment, state guidance) and \cle{second-generation agriculture} for Cameroon's structural transformation.
    \item \textbf{Design} a coherent policy package that balances \cle{privatisation}, \cle{regional policy}, and \cle{population policy} while explicitly recognising \cle{policy conflicts} (e.g., inflation vs. growth, equity vs. efficiency).
    \item \textbf{Communicate} economic reasoning concisely in a \cle{policy brief} format and defend it orally — a key skill for Paper 2 (Essay) and Paper 3 (Data Response) of the GCE Advanced Level examination.
\end{itemize}

\courssub{Situation and Mock Data}
\textbf{Context:} It is January 2034. The \emph{Ministry of the Economy, Planning and Regional Development (MINEPAT)} has convened a \emph{National Economic Foresight Workshop} to recalibrate the \emph{Cameroon Vision 2035} implementation strategy. The Minister has invited teams of young economists (you) to analyse the latest macroeconomic dossier and submit a \emph{Policy Brief} titled ``\emph{Unlocking Structural Transformation: Priorities for the 2034--2035 Fiscal Year}''.

\textbf{Mock Data File — Cameroon Economic Dashboard (Estimates for FY 2023/2024):}
\begin{center}
\begin{tabularx}{\linewidth}{|l|X|}\hline
\rowcolor{popblueL}
\textbf{Indicator} & \textbf{Value / Description} \\ \hline
\textbf{Real GDP Growth Rate} & 3.8\% (down from 4.2\% in 2022; 10-year average: 4.0\%) \\ \hline
\textbf{GDP per capita (PPP, current int.\$)} & 4,150 (Lower Middle Income threshold: ~\$4,500) \\ \hline
\textbf{HDI (2022)} & 0.576 (Rank 151/193) — \emph{Medium Human Development} \\ \hline
\textbf{Inflation Rate (CPI, annual avg.)} & 7.4\% (BEAC target: $\le$ 3\%; driven by food +12.1\%, transport +8.3\%) \\ \hline
\textbf{Current Account Balance} & -3.2\% of GDP (chronic deficit; oil \& gas exports declining, manufactured imports rising) \\ \hline
\textbf{Fiscal Deficit (incl. grants)} & -2.9\% of GDP (CEMAC convergence criterion: $\ge$ -1.5\%) \\ \hline
\textbf{Public Debt / GDP} & 45.8\% (IMF moderate risk; rising debt service / revenue ratio) \\ \hline
\textbf{Youth Unemployment (15--34 yrs, ILO broad)} & 13.2\% (Underemployment: 68\% of employed youth in informal sector) \\ \hline
\textbf{Structural Unemployment Indicator} & Skills mismatch: 62\% of university graduates in fields with low labour demand (humanities, law) vs. STEM/TVET shortage \\ \hline
\textbf{Population (2024 est.)} & 29.3 million; Growth rate: 2.6\% p.a. \\ \hline
\textbf{Age Structure} & 0--14: 42\%; 15--64: 55\%; 65+: 3\% — \emph{High dependency ratio (81.8\%)} \\ \hline
\textbf{Urbanisation Rate} & 58\% (Douala/Yaoundé = 45\% of urban pop.); Rural poverty rate: 56\% vs Urban: 21\% \\ \hline
\textbf{Industrial Share of GDP} & 14.5\% (Manufacturing: 9.2\%; stagnant since 2000) — \emph{Signs of premature deindustrialisation} \\ \hline
\textbf{Agriculture Share of GDP} & 16.8\% (employs 43\% of labour force; low productivity, rain-fed, ageing farmers) \\ \hline
\textbf{Key State-Owned Enterprises (SOEs)} & \textbf{SONARA} (refinery, loss-making, needs \$1.2bn rehabilitation), \textbf{CAMTEL} (telecom, monopoly eroded), \textbf{CAMRAIL} (rail concession), \textbf{CDC} (agro-industry, privatised 2020, restructuring), \textbf{SONEL} (electricity, concessioned to ENEO), \textbf{CAMWATER}, \textbf{CRTV} \\ \hline
\textbf{Business Climate (Doing Business legacy)} & Rank $\approx$ 160/190; \emph{Weak contract enforcement, electricity access, cross-border trade} \\ \hline
\textbf{Regional Disparities} & \emph{Far North, North, Adamawa, East, North-West, South-West} lag significantly in infrastructure, HDI, poverty incidence ($>$50\%) vs. \emph{Littoral, Centre, West, South} \\ \hline
\end{tabularx}
\end{center}

\textbf{Additional Qualitative Intelligence:}
\begin{itemize}
    \item \textbf{CEMAC Constraints:} Fixed exchange rate (CFA Franc pegged to Euro), common external tariff (CET), regional inflation target $\le$ 3\%, fiscal convergence criteria. Monetary policy set by BEAC (not national).
    \item \textbf{AfCFTA:} Cameroon ratified; customs duties elimination schedule ongoing. Opportunity for agro-processed exports (cocoa powder, cotton yarn, plantain chips) but competitive pressure from Nigeria, Egypt, South Africa.
    \item \textbf{Climate Shocks:} Recurrent floods in Far North/Logone, droughts in North, insecurity in NOSO regions disrupting agriculture and transport corridors.
    \item \textbf{NIC Reference:} Vietnam (Doi Moi 1986), Rwanda (Vision 2020/2050), Ethiopia (industrial parks), Morocco (automotive/aerospace clusters) — all combined \emph{state strategic guidance} with \emph{export-oriented FDI attraction} and \emph{agricultural productivity revolution}.
\end{itemize}

\courssub{Tasks}
Work in groups of 4--5. Each group acts as a \emph{Team of Economic Advisers}. Produce a \textbf{one-page Policy Brief} (max 500 words + 1 chart/table) addressing the four tasks below. You have 60 minutes preparation + 5 minutes presentation + 5 minutes Q\&A.

\begin{enumerate}
    \item \textbf{Diagnosis (Business Cycle \& Development Stage):} Using the data, identify the current phase of the \cle{trade cycle} Cameroon is in. Justify using at least three indicators. Classify the economy's development stage using \cle{Rostow's stages} and \cle{HDI} criteria. Explain the concept of \cle{premature deindustrialisation} evident in the data.
    \item \textbf{Prioritisation (Two Most Urgent Macroeconomic Problems):} Select the \textbf{two} most binding constraints from the dashboard. For each, classify its nature (e.g., \cle{demand-pull vs. cost-push inflation}, \cle{structural vs. cyclical vs. frictional unemployment}, \cle{balance of payments structural deficit}). Explain \emph{why} they are more urgent than the fiscal deficit or debt level.
    \item \textbf{Policy Package (Coherent Strategy):} Propose a package of \textbf{four} specific policy measures — one from each pillar:
        \begin{enumerate}
            \item \textbf{Industrialisation / NIC Lesson:} A concrete measure to reverse deindustrialisation drawing on NIC experience (e.g., Special Economic Zones, targeted FDI, local content).
            \item \textbf{Second-Generation Agriculture / Agribusiness:} A measure to raise agricultural productivity and create rural jobs (mechanisation, value chains, irrigation, youth agripreneurs).
            \item \textbf{SOE Reform (Privatisation / Nationalisation / Concession):} Decide the fate of \textbf{SONARA} and \textbf{CAMTEL} — justify using \cle{contestable markets}, \cle{natural monopoly}, \cle{strategic sector} arguments.
            \item \textbf{Population / Human Capital:} A policy to harness the \cle{demographic dividend} and reduce the dependency ratio (education reform, family planning, TVET, diaspora engagement).
        \end{enumerate}
    \item \textbf{Policy Conflict \& Management:} Identify \textbf{one major conflict} between two objectives in your package (e.g., \emph{Inflation targeting vs. Growth stimulation via credit to agriculture}; \emph{Privatisation proceeds vs. Regional equity}; \emph{Exchange rate stability vs. Export competitiveness}). Explain the transmission mechanism and propose a \emph{mitigation strategy} (sequencing, compensatory measures, institutional framework).
\end{enumerate}

\begin{exemplebox}[Model Policy Brief: Unlocking Structural Transformation — Priorities for 2034--2035]
\textbf{To:} Minister of Economy, Planning and Regional Development \hfill \textbf{From:} Advisory Team \#7 (Upper Sixth Economics, GCE A-Level)
\par\medskip\noindent
\textbf{Date:} January 2034 \hfill \textbf{Classification:} Policy Priority — Immediate Action

\textbf{1. Diagnosis: Stagnation Phase with Premature Deindustrialisation}
Cameroon is in the \textbf{stagnation / slowdown phase} of the trade cycle, bordering on a \emph{growth recession} (growth $<$ potential $\approx$ 5--6\% needed for Vision 2035). Evidence: (i) Real GDP growth decelerating to 3.8\% (below 10-yr average 4.0\%); (ii) Investment rate stagnant $\approx$ 18\% GDP (ICOR high); (iii) Manufacturing value added fell 0.4\% y/y in Q3 2023.
\par
Development stage: \textbf{Rostow's \emph{Pre-conditions for Take-off} / \emph{Drive to Maturity} transition stalled}. HDI 0.576 (Medium) confirms \emph{low human development} despite Lower Middle Income status. The economy exhibits \textbf{premature deindustrialisation} (Rodrik 2016): manufacturing share peaked at 9.2\% GDP at per capita income $\ll$ historical peaks (20--25\% in East Asia). Causes: Dutch Disease (oil), weak institutions, energy deficit, import competition without infant industry protection.

\textbf{2. Two Most Urgent Macroeconomic Problems}
\begin{enumerate}
    \item \textbf{Imported Cost-Push Inflation (7.4\%, Food +12.1\%).} \emph{Nature:} Supply-side shock transmitted via CFA franc appreciation vs. trading partners (Euro strength) + global commodity prices + domestic supply bottlenecks (insecurity, poor roads). \emph{Urgency:} Erodes real wages (poverty increase), breaches BEAC convergence criterion ($\le$ 3\%), forces BEAC tight monetary policy (high TIAO rates) crowding out private investment. More urgent than fiscal deficit because it directly hits the poor and constrains monetary sovereignty.
    \item \textbf{Structural Youth Unemployment \& Underemployment (13.2\% open, 68\% underemployed).} \emph{Nature:} \textbf{Structural} — skills mismatch (62\% graduates in low-demand fields) + \textbf{disguised} in rural agriculture (marginal product $\approx$ 0). \emph{Urgency:} Demographic dividend window (55\% working age) closing by 2040; social stability risk (NOSO crisis roots); wastes human capital. Cyclical component small (output gap modest).
\end{enumerate}

\textbf{3. Coherent Policy Package (Four Pillars)}
\begin{center}
\begin{tabularx}{\linewidth}{|l|X|}\hline
\rowcolor{popblueL}
\textbf{Pillar} & \textbf{Measure (2024--2025 Action)} \\ \hline
\textbf{1. Industrialisation (NIC Lesson: Vietnam/Morocco)} & \textbf{Launch 3 \emph{Agro-Industrial Special Economic Zones (SEZs)}:} Douala-Bassa (cocoa grinding, packaging), Ngaoundéré (cotton/textile, leather), Kribi (fish processing, timber). \emph{Incentives:} 10-yr tax holiday, duty-free capital goods import, streamlined \emph{Guichet Unique}, 50\% electricity tariff subsidy for first 5 yrs. \emph{Conditionality:} Export share $\ge$ 70\%, local raw material sourcing $\ge$ 60\%, TVET apprenticeship quota 15\%. \emph{Anchor investors:} Target Malaysian (cocoa), Turkish (textile), Chinese (leather) FDI via bilateral treaties. \\ \hline
\textbf{2. 2nd-Gen Agriculture / Agribusiness} & \textbf{\emph{Programme National de Mécanisation et d'Irrigation (PNMI)}:} State guarantees 50\% of commercial bank loans for youth cooperatives acquiring tractors, solar pumps, processing units (cassava flour, plantain chips, palm oil). \emph{Land reform:} Fast-track \emph{Certificats Fonciers} for 50,000 ha in \emph{Zones d'Expansion Agricole} (ZEA) with 25-yr leaseholds. \emph{Extension:} Deploy 2,000 \emph{Conseillers Agricoles} (digital advisory via USSD/App) — Rwanda \emph{Twitezimbere} model. \\ \hline
\textbf{3. SOE Reform: SONARA \& CAMTEL} & \textbf{SONARA: \emph{Strategic Concession (not full privatisation)}.} 25-yr concession to international operator (TotalEnergies / Vitol / Dangote consortium) with \emph{PSO (Public Service Obligation)}: maintain strategic stock (30 days), supply domestic market at regulated price formula (import parity + transport), invest \$1.2bn modernisation (hydrocracker). State retains 20\% golden share. \textbf{CAMTEL: \emph{Partial Privatisation (49\%) + Open Access.}} List 30\% on Douala Stock Exchange (retail + institutional), sell 19\% to strategic telecom operator (Orange/MTN excluded — competition). Mandate \emph{Open Access} on fibre backbone (wholesale price regulation by ART) to lower broadband costs. \\ \hline
\textbf{4. Population / Human Capital} & \textbf{\emph{Compétences 2035 — TVET Revolution:}} Convert 50\% of public university enrolment growth to \emph{Instituts Universitaires de Technologie (IUT)} and \emph{Centres de Formation Professionnelle (CFP)} aligned with SEZ sectors (mechatronics, food tech, welding, coding). \emph{Financing:} \emph{Skills Levy} (0.5\% payroll) + World Bank \emph{Skills Development Project}. \emph{Demographic:} Scale up \emph{Plan National de Développement du Capital Humain} — free contraceptives in health centres, girls' secondary completion cash transfers (Sahel Women's Empowerment model). \\ \hline
\end{tabularx}
\end{center}

\textbf{4. Policy Conflict \& Management Strategy}
\textbf{Conflict:} \textbf{Inflation Targeting (BEAC tight money) vs. Agro-Industrial Credit Expansion (PNMI / SEZ financing).}
\par
\textbf{Transmission:} BEAC raises TIAO rate $\to$ commercial bank prime rate $\uparrow$ $\to$ cost of working capital for agribusiness $\uparrow$ $\to$ investment in tractors/processing $\downarrow$ $\to$ food supply response $\downarrow$ $\to$ food inflation persists $\to$ real exchange rate appreciates $\to$ SEZ export competitiveness $\downarrow$.
\par
\textbf{Mitigation (Sequencing \& Institutional Design):}
\begin{enumerate}
    \item \textbf{Refinancing Window:} BEAC creates \emph{Fenêtre de Refinancement Agricole et Industriel (FRAI)} at TIAO $-200$ bps for loans certified by \emph{Fonds de Garantie} (backed by AfDB guarantee). \emph{Decouples} policy rate from priority sector cost.
    \item \textbf{Fiscal-Monetary Coordination:} MINEPAT commits to \textbf{zero net domestic financing} of budget (rely on concessional external + privatisation proceeds) $\to$ reduces crowding-out, allows BEAC to ease stance faster.
    \item \textbf{Supply-Side Anchor:} PNMI \emph{in-kind} subsidies (seeds, fertilizer vouchers via e-voucher) bypass credit channel — immediate food supply boost $\to$ lowers CPI $\to$ creates space for rate cuts.
    \item \textbf{Exchange Rate Buffer:} Advocate at BEAC/CEMAC for \emph{asymmetric adjustment} — temporary suspension of convergence criteria for inflation (Art. 14 Treaty) while structural reforms take hold (precedent: 2016--2017 oil shock).
\end{enumerate}

\medskip\noindent
\textbf{Conclusion:} Cameroon stands at a \emph{critical juncture}. The \emph{demographic dividend}, \emph{AfCFTA market access}, and \emph{commodity endowment} offer a narrow window for structural transformation. This package — \emph{SEZ-led industrialisation, mechanised agribusiness, strategic SOE reform, TVET revolution} — addresses binding constraints coherently. The inflation-growth conflict is manageable through \emph{targeted refinancing} and \emph{fiscal discipline}. We recommend immediate Cabinet adoption and a \emph{Quarterly Delivery Unit} in the Prime Minister's Office to monitor KPIs.

\vspace{1em}
\noindent\textbf{Annex: Key Monitoring Indicators (Quarterly Dashboard)}
\begin{center}
\begin{tabular}{|c|l|c|c|}\hline
\rowcolor{popblueL}
\# & KPI & Baseline (2023) & Target (Dec 2025) \\ \hline
1 & Manufacturing Value Added / GDP & 9.2\% & 12.0\% \\ \hline
2 & Youth NEET Rate & 28\% & 20\% \\ \hline
3 & Food Inflation (y/y) & 12.1\% & $\le$ 5.0\% \\ \hline
4 & Agro-processed Exports / Total Exports & 4\% & 12\% \\ \hline
5 & TVET Enrolment Share (Post-Secondary) & 18\% & 35\% \\ \hline
6 & SONARA Capacity Utilisation & 45\% & 85\% \\ \hline
\end{tabular}
\end{center}
\end{exemplebox}

\begin{attention}[Examiner's Tips for the Oral Defence]
\begin{itemize}
    \item \textbf{Know your numbers:} Be ready to quote 3.8\%, 7.4\%, 13.2\%, 45.8\% without notes.
    \item \textbf{Use the syllabus vocabulary:} ``\emph{Premature deindustrialisation}, \emph{imported cost-push inflation}, \emph{structural unemployment}, \emph{demographic dividend}, \emph{contestable market}, \emph{policy conflict}, \emph{transmission mechanism}''.
    \item \textbf{Contextualise:} Mention \emph{CEMAC constraints} (no independent monetary policy), \emph{NOSO insecurity}, \emph{AfCFTA}. This shows \emph{application} (AO3), not just knowledge (AO1).
    \item \textbf{Admit trade-offs:} ``\emph{Privatisation of CAMTEL may raise prices short-term but open access regulation mitigates this}''— this is \emph{evaluation} (AO4).
    \item \textbf{Time management:} 4 minutes presentation, 1 minute for the chart. Leave 5 minutes for Q\&A.
\end{itemize}
\end{attention}

\newpage

\coursec{Methods and worked examples}

\courssub{Method 1: Distinguishing Economic Growth from Economic Development in an Essay}

\begin{methode}[Handling a ``growth versus development'' essay]
This is one of the most frequent GCE Advanced Level essay themes. To score full marks:
\begin{enumerate}
\item \textbf{Define both concepts precisely.} Economic growth is the sustained increase in \cle{real national output} (real GDP) over time. Economic development is a wider, qualitative process: rising living standards, better health and education, reduced poverty and inequality, and structural transformation of the economy.
\item \textbf{State the relationship.} Growth is a \emph{necessary but not sufficient} condition for development: growth provides the resources, but development depends on how they are distributed and used.
\item \textbf{Give the indicators.} Growth is measured by the real GDP growth rate; development by GNI per capita, the Human Development Index (HDI) — combining life expectancy, education (mean and expected years of schooling) and GNI per capita — literacy rates, life expectancy, and poverty incidence.
\item \textbf{Illustrate with examples.} A country may grow rapidly through oil exports (e.g. an oil enclave economy) while most citizens remain poor — growth without development.
\item \textbf{Conclude with a judgement.} Explain under what conditions growth is transformed into development (investment in human capital, good governance, equitable distribution).
\end{enumerate}
\textbf{Reflex:} always write ``\cle{real} GDP'' (never nominal) when defining growth, and always mention at least one \emph{non-income} indicator of development.
\end{methode}

\begin{exoresolu}[Essay: growth without development]
\textbf{Statement (GCE A-Level style).} ``A country can experience rapid economic growth without economic development.'' Discuss this statement with reference to developing countries such as Cameroon.
\tcblower
\textbf{Detailed solution.}

\textbf{Step 1 — Definitions.} Economic growth is the sustained increase in a country's \cle{real gross domestic product} over time. Economic development is a broader concept involving not only higher income but also improvements in health, education, nutrition, reduction of poverty and inequality, and modernisation of institutions and production structures.

\textbf{Step 2 — Why growth may not bring development.}
\begin{itemize}
\item \textbf{Enclave growth.} If growth is driven by a capital-intensive extractive sector (oil, mining) employing few workers, the income accrues to foreign firms and a small elite. GDP rises, but employment and mass living standards do not.
\item \textbf{Unequal distribution.} If the benefits of growth are concentrated in a few hands, average income rises while the majority remains in poverty. A rising Gini coefficient alongside rising GDP illustrates this.
\item \textbf{Neglect of human capital.} If government revenue from growth is not channelled into schools, hospitals and clean water, life expectancy and literacy stagnate.
\item \textbf{Environmental and social costs.} Growth based on deforestation or pollution may reduce welfare even as measured output rises.
\end{itemize}

\textbf{Step 3 — Illustration.} Several oil-exporting African economies recorded high GDP growth during oil booms, yet retained low HDI rankings because oil rents were concentrated and diversification was weak. In Cameroon, growth driven by oil, timber and cocoa exports coexists with significant rural poverty and underemployment in the informal sector — evidence that growth alone does not guarantee development.

\textbf{Step 4 — Conditions for growth to become development.} Growth translates into development when it is \emph{inclusive}: investment in education and health, support to agriculture and small enterprises, infrastructure in rural areas, progressive taxation, and transparent management of public revenue.

\textbf{Step 5 — Conclusion.} The statement is correct: growth is quantitative, development is qualitative. Growth creates the \emph{possibility} of development, but only deliberate policies of redistribution, human-capital formation and structural transformation convert rising output into rising welfare for the whole population.
\end{exoresolu}

\courssub{Method 2: Computing and Interpreting Growth, Unemployment and Inflation Rates}

\begin{methode}[Mastering the key percentage calculations]
Examiners regularly test four numerical skills. Memorise the formulas and the units:
\begin{enumerate}
\item \textbf{Real GDP growth rate:}
\[ g = \frac{\text{Real GDP}_{t} - \text{Real GDP}_{t-1}}{\text{Real GDP}_{t-1}} \times 100 \]
\item \textbf{Unemployment rate:}
\[ u = \frac{\text{Number unemployed}}{\text{Labour force}} \times 100, \qquad \text{Labour force} = \text{employed} + \text{unemployed} \]
\item \textbf{Inflation rate (from CPI):}
\[ \pi = \frac{\text{CPI}_{t} - \text{CPI}_{t-1}}{\text{CPI}_{t-1}} \times 100 \]
\item \textbf{Converting nominal to real values (base year = 100):}
\[ \text{Real value} = \frac{\text{Nominal value}}{\text{Price index}} \times 100 \]
\end{enumerate}
\textbf{Reflexes:} always divide by the \emph{initial} (earlier) value; express the answer as a percentage with one decimal; state clearly what the result \emph{means} in words; the labour force excludes students, pensioners and others not seeking work.
\end{methode}

\begin{exoresolu}[Numerical problem: growth, unemployment and inflation in ``Kamerland'']
\textbf{Statement.} The economy of Kamerland records the following data. Real GDP was $2\,400$ billion FCFA in Year 1 and $2\,520$ billion FCFA in Year 2. The population of working age is $10$ million, of whom $6.2$ million are employed and $0.8$ million are unemployed and seeking work; the rest are not in the labour force. The consumer price index rose from $125$ in Year 1 to $135$ in Year 2. Nominal wage of a typical worker was $50\,000$ FCFA per month in Year 1 and $54\,000$ FCFA in Year 2.
\begin{enumerate}
\item Calculate the economic growth rate in Year 2.
\item Calculate the unemployment rate.
\item Calculate the inflation rate.
\item Calculate the worker's real wage in each year (at base-year prices) and comment.
\end{enumerate}
\tcblower
\textbf{Detailed solution.}

\textbf{1. Growth rate.}
\[ g = \frac{2520 - 2400}{2400} \times 100 = \frac{120}{2400} \times 100 = 5.0\% \]
Kamerland's real output grew by $5\%$ — a respectable rate of economic growth.

\textbf{2. Unemployment rate.} First find the labour force:
\[ \text{Labour force} = 6.2 + 0.8 = 7.0 \text{ million} \]
(The remaining $3$ million of working age are \emph{not} in the labour force — students, homemakers, discouraged workers — and are excluded.)
\[ u = \frac{0.8}{7.0} \times 100 \approx 11.4\% \]
About $11.4\%$ of the labour force is unemployed.

\textbf{3. Inflation rate.}
\[ \pi = \frac{135 - 125}{125} \times 100 = \frac{10}{125} \times 100 = 8.0\% \]
The general price level rose by $8\%$.

\textbf{4. Real wages.}
\[ \text{Real wage}_{\text{Y1}} = \frac{50\,000}{125} \times 100 = 40\,000 \text{ FCFA} \]
\[ \text{Real wage}_{\text{Y2}} = \frac{54\,000}{135} \times 100 = 40\,000 \text{ FCFA} \]

\textbf{Comment.} Although the \emph{nominal} wage rose by $8\%$ (from $50\,000$ to $54\,000$ FCFA), prices also rose by exactly $8\%$, so the \emph{real} wage — purchasing power — is unchanged at $40\,000$ FCFA at base-year prices. This illustrates \cle{money illusion}: workers feel richer because their pay packet is bigger, yet they can buy exactly the same basket of goods. It also shows why examiners insist on distinguishing nominal from real magnitudes.
\end{exoresolu}

\courssub{Method 3: Analysing the Trade (Business) Cycle}

\begin{methode}[Identifying and explaining the phases of the cycle]
\begin{enumerate}
\item \textbf{Draw the diagram first.} Plot real national output (vertical axis) against time (horizontal axis); draw a smooth upward trend line and a wave oscillating around it.
\item \textbf{Name the four phases in order:} \cle{boom} (peak), \cle{recession/downturn} (falling output), \cle{trough/depression} (lowest point), \cle{recovery/upturn} (rising output back towards the trend).
\item \textbf{Attach the symptoms to each phase.} Boom: low unemployment, rising prices, capacity pressure. Recession: falling investment and demand, rising unemployment, slowing inflation. Trough: high unemployment, spare capacity, low confidence. Recovery: rising investment, falling unemployment, improving business confidence.
\item \textbf{Explain the mechanism} using the multiplier and accelerator: a fall in investment reduces income (multiplier), which reduces consumption, which further reduces investment (accelerator) — and symmetrically in the upswing.
\item \textbf{Link to policy.} Governments use counter-cyclical fiscal and monetary policy: stimulate in recession, restrain in boom.
\end{enumerate}
\textbf{Reflex:} never confuse a \emph{recession} (output falling, often defined as two consecutive quarters of negative growth) with a \emph{depression} (a deep, prolonged trough).
\end{methode}

\begin{popfigure}
\begin{tikzpicture}
\begin{axis}[
    width=\linewidth,
    height=6cm,
    axis lines=middle,
    xlabel={Time},
    ylabel={Real GDP},
    xtick=\empty,
    ytick=\empty,
    xmin=0, xmax=10,
    ymin=0, ymax=10,
    clip=false
]
% Trend line
\addplot[popblue, thick, dashed, domain=0:10] {0.5*x + 2};
% Cycle
\addplot[poporange, very thick, smooth, domain=0:10] {0.5*x + 2 + 2*sin(deg(0.8*x))};
% Phase labels
\node[poporange, anchor=south] at (axis cs:2.5, 5.5) {Boom};
\node[poporange, anchor=north] at (axis cs:5, 3.5) {Recession};
\node[poporange, anchor=north] at (axis cs:6.5, 3.5) {Trough};
\node[poporange, anchor=south] at (axis cs:8.5, 6.5) {Recovery};
\end{axis}
\end{tikzpicture}
\legende{The trade (business) cycle: fluctuations around a long-run trend}
\end{popfigure}

\begin{exoresolu}[Question: reading a business-cycle diagram]
\textbf{Statement.} The diagram below (described in data form) shows an economy whose real GDP, in billions FCFA, evolves as follows: Year 1: $1\,000$; Year 2: $1\,080$; Year 3: $1\,120$ (peak); Year 4: $1\,060$; Year 5: $1\,000$ (trough); Year 6: $1\,050$; Year 7: $1\,110$.
\begin{enumerate}
\item Identify the phase of the trade cycle in each period.
\item Describe the likely behaviour of unemployment and inflation between Year 3 and Year 5.
\item Suggest two policy measures the government could adopt in Year 5.
\end{enumerate}
\tcblower
\textbf{Detailed solution.}

\textbf{1. Phases.}
\begin{itemize}
\item Years 1--3: output rising above trend — \textbf{upswing/recovery turning into boom}, peaking in Year 3 at $1\,120$ billion.
\item Years 3--5: output falling from $1\,120$ to $1\,000$ billion — \textbf{recession (downswing)}.
\item Year 5: lowest point — \textbf{trough}.
\item Years 5--7: output rising again — \textbf{recovery}.
\end{itemize}

\textbf{2. Behaviour of unemployment and inflation (Year 3 to Year 5).} As output contracts, firms cut production and lay off workers, so \cle{cyclical unemployment} rises. With demand falling, price increases slow down — \cle{disinflation} — and if the downturn is severe, prices may actually fall (deflation). Business confidence and investment collapse, reinforcing the decline through the multiplier--accelerator mechanism.

\textbf{3. Policies in Year 5 (the trough).} Any two of:
\begin{itemize}
\item \textbf{Expansionary fiscal policy:} increase government spending on infrastructure (roads, schools) or cut taxes to raise aggregate demand; the multiplier then amplifies the initial injection.
\item \textbf{Expansionary monetary policy:} lower interest rates and expand credit to encourage investment and consumption.
\end{itemize}
Such \emph{counter-cyclical} (reflationary) policies aim to shorten the recession and pull the economy back towards its trend growth path.
\end{exoresolu}

\courssub{Method 4: Diagnosing Types of Unemployment and Matching Policies}

\begin{methode}[Classifying unemployment and prescribing the cure]
\begin{enumerate}
\item \textbf{Ask: what is the underlying cause?} Then classify:
\begin{itemize}
\item \textbf{Frictional:} between jobs, searching — short term.
\item \textbf{Structural:} skills or location mismatch as the structure of the economy changes (e.g. decline of an industry).
\item \textbf{Cyclical (demand-deficient):} caused by a fall in aggregate demand during recession.
\item \textbf{Seasonal:} tied to seasons (agriculture, tourism).
\item \textbf{Technological:} machines replace labour.
\item \textbf{Disguised/underemployment:} workers nominally employed but with very low productivity — widespread in Cameroon's informal and rural sectors.
\end{itemize}
\item \textbf{Match the policy to the cause:} reflationary fiscal/monetary policy for cyclical unemployment; retraining, education and mobility schemes for structural unemployment; information and job-placement services for frictional unemployment; diversification of rural activities for seasonal unemployment.
\item \textbf{Evaluate:} note that using demand policy against structural unemployment causes inflation without creating durable jobs — a classic exam trap.
\end{enumerate}
\textbf{Reflex:} in Cameroonian answers, always mention the \cle{informal sector}, graduate unemployment and underemployment, and the mismatch between training and labour-market needs.
\end{methode}

\begin{exoresolu}[Guided work: diagnosing unemployment cases]
\textbf{Statement.} For each of the following situations in Cameroon, identify the type of unemployment and propose one appropriate policy:
\begin{enumerate}
\item[(a)] Workers of a cotton-processing factory in the North lose their jobs when the factory closes because of foreign competition.
\item[(b)] Thousands of university graduates in Yaounde spend months searching for their first job after graduation.
\item[(c)] During a recession, a construction company in Douala lays off workers because orders for new buildings collapse.
\item[(d)] Labourers on cocoa farms near Ebolowa have work only during the harvest season.
\end{enumerate}
\tcblower
\textbf{Detailed solution.}

\textbf{(a) Structural unemployment.} The jobs disappeared because the structure of the economy changed — the industry is no longer competitive. \emph{Policy:} retraining programmes to give workers skills demanded by growing sectors, plus incentives to attract new industries to the region (regional policy). Demand stimulation alone would not recreate these jobs.

\textbf{(b) Frictional unemployment.} Graduates are ``between'' education and employment; the search takes time because of imperfect information. \emph{Policy:} improve job-information and placement services (employment agencies, career guidance, internship schemes linking universities to firms). Note: if graduates remain jobless for years because their training does not match employers' needs, the problem becomes \emph{structural} — requiring reform of curricula towards technical and entrepreneurial skills.

\textbf{(c) Cyclical (demand-deficient) unemployment.} The cause is the fall in aggregate demand during the recession. \emph{Policy:} expansionary fiscal policy — for example, government investment in public works (roads, housing) which directly re-employs construction workers and, through the multiplier, revives demand throughout the economy; possibly complemented by lower interest rates.

\textbf{(d) Seasonal unemployment.} Demand for agricultural labour varies with the farming calendar. \emph{Policy:} promote off-season activities — irrigation allowing several cropping cycles (a goal of second generation agriculture), rural agro-processing industries, and rural crafts — so that labour is occupied all year round.

\textbf{Key lesson:} the correct diagnosis determines the correct cure; a single ``anti-unemployment policy'' cannot work for all types.
\end{exoresolu}

\courssub{Method 5: Explaining Inflation and Deflation — Causes, Effects, Remedies}

\begin{methode}[Building a complete inflation/deflation answer]
\begin{enumerate}
\item \textbf{Define.} Inflation: a sustained rise in the \emph{general} price level (fall in the value of money). Deflation: a sustained \emph{fall} in the general price level.
\item \textbf{Classify the causes of inflation:} \cle{demand-pull} (aggregate demand exceeds supply at full employment — ``too much money chasing too few goods''), \cle{cost-push} (rising wages, imported input and fuel prices, taxes), and \cle{monetary} (excessive money-supply growth).
\item \textbf{State the effects of inflation:} erosion of purchasing power of fixed incomes, redistribution from lenders to borrowers, menu and shoe-leather costs, uncertainty discouraging investment, loss of export competitiveness (if faster than abroad).
\item \textbf{State the dangers of deflation:} consumers delay purchases expecting lower prices, real debt burdens rise, firms cut output and jobs — a deflationary spiral.
\item \textbf{Match remedies:} demand-pull $\rightarrow$ contractionary fiscal/monetary policy; cost-push $\rightarrow$ supply-side measures, incomes policy; imported inflation $\rightarrow$ difficult for a small open economy, requiring productivity gains.
\end{enumerate}
\textbf{Reflex:} distinguish \cle{disinflation} (inflation rate falling but still positive) from \cle{deflation} (prices actually falling) — a frequent exam trap.
\end{methode}

\begin{exoresolu}[Data response: interpreting price data]
\textbf{Statement.} The annual inflation rate in a CEMAC country is: Year 1: $6\%$; Year 2: $4\%$; Year 3: $2\%$; Year 4: $-1\%$.
\begin{enumerate}
\item In which years is the country experiencing disinflation? Deflation?
\item Are prices higher or lower in Year 3 than in Year 1? Explain.
\item Explain two economic dangers of the situation in Year 4.
\item Suggest one monetary and one fiscal measure appropriate to Year 4.
\end{enumerate}
\tcblower
\textbf{Detailed solution.}

\textbf{1. Disinflation and deflation.} Years 2 and 3 show \textbf{disinflation}: prices are still rising, but at a decreasing rate ($4\%$, then $2\%$). Year 4 shows \textbf{deflation}: the price level actually falls by $1\%$.

\textbf{2. Price level in Year 3 versus Year 1.} Prices are \emph{higher} in Year 3. A positive inflation rate, however small, means the general price level is still rising. Taking Year 1 prices as $100$: Year 2 $\approx 104$, Year 3 $\approx 104 \times 1.02 \approx 106.1$. Falling inflation does \emph{not} mean falling prices.

\textbf{3. Dangers of deflation (Year 4).}
\begin{itemize}
\item \textbf{Postponed consumption:} households delay purchases expecting still lower prices, so aggregate demand falls, firms' sales collapse, and they cut production and employment — deepening the recession (deflationary spiral).
\item \textbf{Rising real debt burden:} with falling prices, the real value of debts increases, squeezing indebted firms and households, raising bankruptcies and bank failures.
\end{itemize}

\textbf{4. Remedies for Year 4.} \emph{Monetary:} the central bank lowers interest rates and expands credit to stimulate spending and investment. \emph{Fiscal:} the government increases public expenditure (e.g. infrastructure) and/or cuts taxes to inject demand into the economy. Both are \cle{reflationary} policies designed to reverse the deflation and restart growth.
\end{exoresolu}

\courssub{Method 6: Evaluating Privatisation and Nationalisation}

\begin{methode}[Writing a balanced evaluation essay]
\begin{enumerate}
\item \textbf{Define both terms.} Privatisation: transfer of ownership/control of state-owned enterprises to the private sector (sale of shares, concessions, management contracts). Nationalisation: transfer of private (often foreign) assets into state ownership, with or without compensation.
\item \textbf{List arguments for privatisation:} efficiency and profit motive, reduced budgetary burden of loss-making parastatals, competition and better service, attraction of private and foreign capital, wider share ownership.
\item \textbf{List arguments against (limits):} natural monopolies may become \emph{private} monopolies exploiting consumers; job losses from restructuring; sale of strategic national assets; foreign control; ``selling the family silver'' — undervaluation.
\item \textbf{List arguments for nationalisation:} control of strategic sectors, protection of employment, planning of essential services, retention of profits for the state.
\item \textbf{List arguments against:} inefficiency, political interference, fiscal burden of subsidies, discouragement of private investment.
\item \textbf{Conclude with a judgement} linked to context: privatisation works best with competition and strong regulation; nationalisation may be justified for natural monopolies and strategic services.
\end{enumerate}
\textbf{Reflex:} always include the role of \cle{regulation} — the success of privatisation depends less on ownership than on competition and regulatory quality.
\end{methode}

\begin{exoresolu}[Guided work: privatisation in Cameroon]
\textbf{Statement.} ``Privatisation is always beneficial to a developing economy.'' Discuss, with reference to the Cameroonian experience (e.g. water, electricity, telecommunications, agro-industrial parastatals).
\tcblower
\textbf{Detailed solution.}

\textbf{Step 1 — Definition and context.} Privatisation is the transfer of state-owned enterprises to private ownership or management. In Cameroon, as in many African countries, privatisation was encouraged from the late 1980s under structural adjustment programmes, because many parastatals were loss-making and drained the public budget.

\textbf{Step 2 — Potential benefits.}
\begin{itemize}
\item \textbf{Efficiency:} private managers face the profit motive and competition, encouraging cost control and better service quality.
\item \textbf{Fiscal relief:} the state stops subsidising loss-making firms and may earn revenue from the sale, freeing resources for health and education.
\item \textbf{Investment:} private (including foreign) capital can modernise networks — for example expanding electricity and telephone coverage faster than a cash-strapped state could.
\end{itemize}

\textbf{Step 3 — Limitations and risks.}
\begin{itemize}
\item \textbf{Private monopoly:} utilities such as water and electricity are natural monopolies; privatising them without regulation can mean higher tariffs and exclusion of poor households.
\item \textbf{Employment:} restructuring usually involves retrenchment, worsening unemployment in the short run.
\item \textbf{Strategic control:} key assets may pass into foreign hands, and profits may be repatriated rather than reinvested locally.
\item \textbf{Undervaluation and corruption:} if sales are not transparent, public assets may be sold below their true value.
\end{itemize}

\textbf{Step 4 — Evaluation.} The Cameroonian experience is mixed: some privatisations improved service coverage and reduced the fiscal burden, while others were criticised for job losses, tariff increases and weak regulatory oversight. The outcome depends on \emph{how} privatisation is carried out: competitive tendering, independent regulation, protection of consumers and workers, and retention of state responsibility for strategic and social objectives.

\textbf{Step 5 — Conclusion.} The statement is too sweeping. Privatisation is beneficial \emph{conditionally} — where competition is possible, regulation is strong and the process is transparent. For natural monopolies and essential social services, continued public ownership or tightly regulated concessions may serve development better.
\end{exoresolu}

\courssub{Method 7: Analysing Population Questions — Optimum Population and Population Policy}

\begin{methode}[Tackling population essays and diagrams]
\begin{enumerate}
\item \textbf{Define the key terms:} \cle{optimum population} — the population size which, given resources and technology, maximises output (income) per head; \textbf{under-population} — too few people to exploit resources fully (output per head below maximum, rising as population grows); \textbf{over-population} — too many people relative to resources (output per head falling as population grows).
\item \textbf{Draw the diagram:} output per head on the vertical axis, population size on the horizontal axis; an inverted-U curve peaking at the optimum population $P^{*}$.
\item \textbf{Stress that the optimum is dynamic:} technical progress, new resources or capital accumulation shift the curve upward and the optimum to the right.
\item \textbf{For growing/ageing populations,} analyse consequences: a young population (as in Cameroon) means a high dependency ratio, pressure on schools and jobs, but a future large labour force; an ageing population (developed countries) means rising pension and health costs and labour shortages.
\item \textbf{Link to population policy:} family planning and education (especially of girls) where population grows too fast; pro-natalist measures, immigration and later retirement where population ages.
\end{enumerate}
\textbf{Reflex:} optimum population is a \emph{theoretical} concept — it cannot be measured precisely, and it changes with technology; say so in your conclusion.
\end{methode}

\begin{popfigure}
\begin{tikzpicture}
\begin{axis}[
    width=\linewidth,
    height=6cm,
    axis lines=middle,
    xlabel={Population size},
    ylabel={Output per head},
    xtick=\empty,
    ytick=\empty,
    xmin=0, xmax=10,
    ymin=0, ymax=10,
    clip=false
]
% Inverted U curve
\addplot[popgreen, very thick, smooth, domain=0.5:9.5] {-0.1*(x-5)^2 + 7};
% Optimum line
\draw[popmuted, dashed] (axis cs:5,0) -- (axis cs:5,7);
\node[popgreen, anchor=south] at (axis cs:5,7.2) {$P^*$};
\node[popmuted, anchor=west] at (axis cs:2,4) {Under-population};
\node[popmuted, anchor=east] at (axis cs:8,4) {Over-population};
\end{axis}
\end{tikzpicture}
\legende{Optimum population: output per head peaks at $P^*$}
\end{popfigure}

\begin{exoresolu}[Essay: is Cameroon's rapid population growth a blessing or a burden?]
\textbf{Statement.} ``Rapid population growth is an obstacle to economic development.'' Discuss with reference to Cameroon, using the concept of optimum population.
\tcblower
\textbf{Detailed solution.}

\textbf{Step 1 — Concept.} The \cle{optimum population} is the size of population which, with given natural resources, capital and technology, yields the maximum output per head. Below it, resources are under-exploited (under-population); beyond it, diminishing returns set in and income per head falls (over-population).

\textbf{Step 2 — Arguments that rapid growth is a burden.}
\begin{itemize}
\item \textbf{High dependency ratio:} a very young population means many children per worker, consuming rather than producing, and straining household savings.
\item \textbf{Pressure on services:} schools, health centres and housing must expand faster than the population merely to maintain standards.
\item \textbf{Pressure on jobs and land:} the labour force grows faster than formal employment, swelling unemployment and the informal sector; in rural areas, land is subdivided into ever-smaller plots, reducing farm productivity.
\item \textbf{Low capital formation:} consumption absorbs income, limiting savings and investment — the economy must ``run to stand still''.
\end{itemize}

\textbf{Step 3 — Arguments that it can be a blessing.}
\begin{itemize}
\item \textbf{Larger market:} a growing population expands demand, allowing economies of scale and attracting investment (a large domestic market within AfCFTA).
\item \textbf{Future labour force:} today's children are tomorrow's workers; with education and job creation, Cameroon could reap a \emph{demographic dividend}.
\item \textbf{Under-exploited resources:} Cameroon has abundant land and natural resources; some regions may actually be \emph{under}-populated relative to their potential.
\end{itemize}

\textbf{Step 4 — Synthesis using the optimum concept.} Whether population growth is a burden depends on whether the economy is pushed beyond its optimum. Because the optimum \emph{moves} with technology, education and capital accumulation, policies that raise productivity — second generation agriculture, industrialisation, human-capital investment — shift the optimum outward and can convert population growth from a burden into an asset.

\textbf{Step 5 — Conclusion.} Rapid population growth is not automatically an obstacle: it becomes one only when it outpaces the growth of resources, capital and jobs. Cameroon's challenge is therefore twofold — moderate fertility through education and family-planning policy, and accelerate investment so that the growing population becomes a productive labour force rather than a dependent mass.
\end{exoresolu}

\newpage

\coursec{Exercises}

\courssub{I check my knowledge}

\begin{exercice}[Multiple choice questions]
For each question, choose the single correct answer.
\begin{enumerate}
\item Economic development differs from economic growth because development:
\begin{enumerate}
\item is measured only by real GDP;
\item includes improvements in living standards, health and education;
\item concerns only industrial output;
\item can occur without any growth in income.
\end{enumerate}
\item The phase of the trade cycle in which output, employment and prices are falling is called:
\begin{enumerate}
\item the boom; \item the recovery; \item the recession; \item the peak.
\end{enumerate}
\item Unemployment caused by a permanent decline of an industry such as a fall in world demand for a traditional export is:
\begin{enumerate}
\item frictional; \item seasonal; \item structural; \item cyclical.
\end{enumerate}
\item A fall in the general price level is called:
\begin{enumerate}
\item disinflation; \item deflation; \item reflation; \item stagflation.
\end{enumerate}
\end{enumerate}
\end{exercice}

\begin{exercice}[True or false — justify your answer]
State whether each statement is \textbf{true} or \textbf{false}, and justify in one or two sentences.
\begin{enumerate}
\item A country can experience economic growth without economic development.
\item Deindustrialisation means that a country's total output of goods and services must be falling.
\item The multiplier and the accelerator work in opposite directions and therefore cancel each other out.
\item Privatisation of a state monopoly automatically creates competition.
\item An ageing population tends to raise the dependency ratio.
\end{enumerate}
\end{exercice}

\begin{exercice}[Key definitions]
Define precisely each of the following terms, giving one example for each:
\begin{enumerate}
\item economic development;
\item reindustrialisation;
\item cyclical (demand-deficient) unemployment;
\item disinflation;
\item nationalisation;
\item optimum population.
\end{enumerate}
\end{exercice}

\begin{exercice}[Direct calculations]
\begin{enumerate}
\item The consumer price index (CPI) of a country rises from $120$ in year 1 to $132$ in year 2. Calculate the rate of inflation in year 2.
\item In a economy, the labour force is $10\,000\,000$ persons and $1\,200\,000$ persons are unemployed. Calculate the unemployment rate.
\item The marginal propensity to consume is $0.8$. Calculate the value of the multiplier, and the final change in national income resulting from an autonomous increase in investment of $500$ million FCFA.
\item A country has a population of $25$ million, of which $45\%$ are aged under 15 and $5\%$ are aged 65 and over. Calculate the dependency ratio.
\end{enumerate}
\end{exercice}

\courssub{I practise}

\begin{exercice}[Inflation rates from a CPI table]
The table below gives the consumer price index of a hypothetical CEMAC country over five years.
\begin{center}
\begin{tabularx}{\linewidth}{|l|X|X|X|X|X|}
\hline
\rowcolor{popblueL} Year & 1 & 2 & 3 & 4 & 5 \\
\hline CPI & 100 & 104 & 110 & 115 & 115 \\
\hline
\end{tabularx}
\end{center}
\begin{enumerate}
\item Calculate the rate of inflation for years 2, 3, 4 and 5.
\item Identify the year in which the country experienced disinflation, justifying your answer.
\item Explain why a constant CPI in year 5 does not mean that prices have fallen back to their year 1 level.
\item Suggest one monetary policy and one fiscal policy that could have been used to reduce inflation between years 3 and 4.
\end{enumerate}
\end{exercice}

\begin{exercice}[The multiplier--accelerator at work]
In the economy of a Central African country, the marginal propensity to consume is $0.75$. The government increases its spending on road construction by $200$ million FCFA.
\begin{enumerate}
\item Calculate the multiplier and the total change in national income.
\item Using the accelerator principle, explain how the rise in national income may induce further investment by firms.
\item Explain, using the interaction of the multiplier and the accelerator, how a \emph{fall} in investment of the same amount could push the economy into a recession.
\item State two limitations of the simple multiplier model in the Cameroonian context.
\end{enumerate}
\end{exercice}

\begin{exercice}[Classifying unemployment]
For each of the following situations in Cameroon, identify the type of unemployment involved and justify your choice:
\begin{enumerate}
\item a worker dismissed by a textile factory in Douala whose products can no longer compete with imported second-hand clothing;
\item a university graduate who spends six months searching for a first job after graduation;
\item casual labourers in the cocoa sector who are out of work between harvest seasons;
\item employees of a bank laid off during a severe recession;
\item a secretary replaced by computerised office software.
\end{enumerate}
\end{exercice}

\begin{exercice}[Reading a population pyramid]
A population pyramid for Cameroon shows a very broad base (large cohorts aged 0--14), a narrowing middle and a very small top (persons aged 65 and over). The youth dependency ratio is estimated at about $80\%$.
\begin{enumerate}
\item Describe the shape of the pyramid and state what it reveals about the birth rate and life expectancy.
\item Calculate the total dependency ratio if the old-age dependency ratio is $8\%$.
\item Explain two economic opportunities and two economic challenges created by this age structure.
\item Suggest two policies the government could adopt to transform this youthful population into a ``demographic dividend''.
\end{enumerate}
\end{exercice}

\courssub{I prepare for the GCE}

\begin{exercice}[Essay — Growth and development (25 marks)]
``Economic growth is necessary but not sufficient for economic development.'' Discuss this statement with reference to Cameroon. In your answer you should: define both concepts precisely; explain the indicators used to measure each; analyse why growth may fail to translate into development; and evaluate policies that could make growth more inclusive in Cameroon. \hfill (25 marks)
\end{exercice}

\begin{exercice}[Data response — Inflation in Cameroon (20 marks)]
Study the following data and answer the questions.
\begin{center}
\begin{tabularx}{\linewidth}{|l|X|X|X|X|X|}
\hline
\rowcolor{popblueL} Year & 1 & 2 & 3 & 4 & 5 \\
\hline CPI (base 100, year 1) & 100 & 103 & 108 & 116 & 121 \\
\hline
\end{tabularx}
\end{center}
\emph{Extract:} ``Over the period, the prices of imported food and fuel rose sharply following the depreciation of the currency and higher world prices, while domestic demand remained weak.''
\begin{enumerate}
\item Calculate the rate of inflation for each year from year 2 to year 5. \hfill (4 marks)
\item Using the extract, identify the type of inflation affecting the country and explain its causes. \hfill (4 marks)
\item Explain two economic costs of the inflation observed. \hfill (4 marks)
\item Evaluate two policies the government could use to control this type of inflation, indicating the limitations of each. \hfill (8 marks)
\end{enumerate}
\end{exercice}

\begin{exercice}[Structured question — The trade cycle and unemployment (20 marks)]
\begin{enumerate}
\item Describe the four phases of the trade (business) cycle, indicating the behaviour of output, employment and the price level in each phase. \hfill (6 marks)
\item Using the multiplier--accelerator model, explain how an initial fall in investment can trigger a cumulative recession. \hfill (6 marks)
\item Define cyclical unemployment and explain why it rises during the downswing of the cycle. \hfill (4 marks)
\item Assess the effectiveness of reflationary fiscal policy as a remedy for cyclical unemployment in a country like Cameroon. \hfill (4 marks)
\end{enumerate}
\end{exercice}

\begin{exercice}[Essay — Privatisation (25 marks)]
``Privatisation always improves economic efficiency.'' Critically examine this statement using Cameroonian examples such as SONEL and CAMTEL. Your answer should distinguish productive from allocative efficiency, consider the cases of natural monopolies and strategic industries, and reach a justified conclusion. \hfill (25 marks)
\end{exercice}

\begin{exercice}[Essay — Industrialisation strategies (25 marks)]
Compare import-substitution industrialisation and export-oriented industrialisation. Drawing lessons from the experience of the Newly Industrialised Countries (NICs), assess which strategy is better suited to Cameroon, taking into account the promotion of second-generation agriculture and agribusiness. \hfill (25 marks)
\end{exercice}

\begin{exercice}[Data response — Population and development (20 marks)]
The table below summarises demographic indicators for a hypothetical African economy.
\begin{center}
\begin{tabularx}{\linewidth}{|l|X|X|}
\hline
\rowcolor{popblueL} Indicator & Country A & Country B \\
\hline Population growth rate (\% per year) & 2.6 & 0.4 \\
\hline Share of population under 15 (\%) & 43 & 16 \\
\hline Share of population aged 65 and over (\%) & 4 & 19 \\
\hline
\end{tabularx}
\end{center}
\begin{enumerate}
\item Identify which country has a youthful population and which has an ageing population, justifying your answer. \hfill (4 marks)
\item Explain the economic challenges faced by each country. \hfill (6 marks)
\item With reference to the concept of optimum population, explain why rapid population growth may either help or hinder development in Country A. \hfill (5 marks)
\item Evaluate two population policies that Country A could adopt. \hfill (5 marks)
\end{enumerate}
\end{exercice}

\begin{exercice}[Short-answer set — Deflation and reflation (15 marks)]
\begin{enumerate}
\item Distinguish clearly between deflation and disinflation. \hfill (3 marks)
\item Explain two economic costs of a prolonged deflation. \hfill (4 marks)
\item With the help of an AD/AS diagram, show how a reflationary policy can move an economy out of a deflationary gap, and state one risk of such a policy. \hfill (5 marks)
\item Explain why nationalisation is sometimes proposed as a response to the failure of strategic private firms, giving one possible advantage and one possible disadvantage. \hfill (3 marks)
\end{enumerate}
\end{exercice}

\newpage

\coursec{Solutions to the exercises}

\begin{corrige}[Exercise 1 — Multiple choice questions]
\begin{enumerate}
\item \textbf{Correct answer: (b)} — Economic development is a multidimensional concept: it includes not only the rise in real income per head (growth) but also improvements in living standards, health, education, equity and structural transformation of the economy. Option (a) describes growth only; (c) confuses development with industrialisation; (d) is wrong because sustained development in practice requires rising real income.
\item \textbf{Correct answer: (c)} — The \emph{recession} (downswing/contraction) is the phase in which output, employment and the price level are falling. The boom and the peak are phases of high activity; the recovery is the upswing after the trough.
\item \textbf{Correct answer: (c)} — A permanent decline of an industry (e.g. a traditional export facing falling world demand) destroys jobs that do not reappear even when demand recovers: this is \emph{structural} unemployment. Cyclical unemployment is caused by deficient aggregate demand, frictional by job search, and seasonal by the rhythm of seasons.
\item \textbf{Correct answer: (b)} — \emph{Deflation} is a sustained fall in the general price level (negative inflation). Disinflation is a fall in the \emph{rate} of inflation (prices still rising, but more slowly); reflation is a policy of stimulating demand; stagflation combines stagnation with inflation.
\end{enumerate}
\end{corrige}

\begin{corrige}[Exercise 2 — True or false]
\begin{enumerate}
\item \textbf{True.} Growth (rising real GDP) may be concentrated in an enclave sector (e.g. oil extraction) with the income accruing to a small elite or foreign firms, while health, education and living standards of the majority stagnate. Cameroon experienced episodes of oil-driven growth without commensurate improvements in human development indicators.
\item \textbf{False.} Deindustrialisation means a fall in the \emph{share} of industry (especially manufacturing) in output and/or employment. Total output may still be rising if services expand faster, as happened in many advanced economies.
\item \textbf{False.} The multiplier and the accelerator do not cancel out; they \emph{reinforce} each other. An initial investment raises income by a multiple (multiplier); the rise in income induces further investment (accelerator), which again multiplies income. This interaction explains the cumulative nature of booms and recessions.
\item \textbf{False.} Privatising a state monopoly merely transfers ownership; if the market structure is unchanged, a \emph{private} monopoly replaces the public one. Competition requires liberalisation of entry and/or regulation. The case of electricity in Cameroon (SONEL/ENEO) illustrates this.
\item \textbf{True.} The dependency ratio $= \dfrac{\text{population} 0\text{--}14 + \text{population} 65+}{\text{population} 15\text{--}64}$. Ageing raises the number of persons aged 65+ relative to the working-age population, so the old-age and total dependency ratios rise.
\end{enumerate}
\end{corrige}

\begin{corrige}[Exercise 3 — Key definitions]
\begin{enumerate}
\item \textbf{Economic development:} a sustained, multidimensional process involving rising real income per head, structural transformation of the economy, and improvements in health, education, equity and living standards of the whole population. \emph{Example:} a country whose GDP per head rises while literacy, life expectancy and access to clean water also improve.
\item \textbf{Reindustrialisation:} a deliberate policy of rebuilding the industrial (especially manufacturing) base of an economy after a period of deindustrialisation, through investment, incentives and support to new industries. \emph{Example:} Cameroon's efforts to develop local transformation of cocoa and agro-industry under its import-substitution drive.
\item \textbf{Cyclical (demand-deficient) unemployment:} unemployment caused by a general deficiency of aggregate demand during the downswing of the trade cycle; firms cut output and lay off workers. \emph{Example:} bank employees laid off during a severe recession.
\item \textbf{Disinflation:} a fall in the \emph{rate} of inflation, with the general price level still rising but more slowly. \emph{Example:} inflation falling from $6\%$ to $3\%$.
\item \textbf{Nationalisation:} the transfer of ownership and control of a private firm or industry to the State, usually with compensation. \emph{Example:} the re-nationalisation of electricity distribution in Cameroon in 2014 (creation of ENEO's public shareholding structure after AES-SONEL).
\item \textbf{Optimum population:} the population size which, given the stock of natural resources, capital and technology, maximises output (income) per head. \emph{Example:} a country where adding one more worker would lower average product is beyond its optimum (over-populated relative to resources).
\end{enumerate}
\end{corrige}

\begin{corrige}[Exercise 4 — Direct calculations]
\begin{enumerate}
\item \textbf{Inflation rate (year 2):}
\[ \pi = \frac{132 - 120}{120} \times 100\% = \frac{12}{120} \times 100\% = 10\%. \]
\item \textbf{Unemployment rate:}
\[ u = \frac{1\,200\,000}{10\,000\,000} \times 100\% = 12\%. \]
\item \textbf{Multiplier and change in income:}
\[ k = \frac{1}{1 - 0.8} = 5; \qquad \Delta Y = k \times \Delta I = 5 \times 500 = 2\,500 \text{ million FCFA}. \]
\item \textbf{Dependency ratio:} dependents $= (0.45 + 0.05) \times 25 = 12.5$ million; working-age population $= 25 - 12.5 = 12.5$ million.
\[ \text{Dependency ratio} = \frac{12.5}{12.5} \times 100\% = 100\%. \]
Each person of working age supports, on average, one dependent.
\end{enumerate}
\end{corrige}

\begin{corrige}[Exercise 5 — Inflation rates from a CPI table]
\begin{enumerate}
\item Applying $\pi_t = \dfrac{\text{CPI}_t - \text{CPI}_{t-1}}{\text{CPI}_{t-1}} \times 100\%$:
\begin{align*}
\text{Year 2: } & \frac{104-100}{100}\times 100\% = 4\% \\
\text{Year 3: } & \frac{110-104}{104}\times 100\% \approx 5.8\% \\
\text{Year 4: } & \frac{115-110}{110}\times 100\% \approx 4.5\% \\
\text{Year 5: } & \frac{115-115}{115}\times 100\% = 0\%
\end{align*}
\item \textbf{Disinflation occurred in year 4}: the inflation rate fell from about $5.8\%$ to about $4.5\%$, but remained positive — prices were still rising, only more slowly.
\item A constant CPI in year 5 means prices \emph{stopped rising} (zero inflation), not that they fell. The price level remains at $115$, i.e. $15\%$ above its year 1 level. Prices would return to the year 1 level only with \emph{deflation} (a fall of the CPI back to 100).
\item \textbf{Monetary policy:} the central bank (BEAC in the CEMAC zone) could raise its key interest rate or tighten reserve requirements to reduce credit and aggregate demand. \textbf{Fiscal policy:} the government could reduce public expenditure or raise taxes to withdraw purchasing power and ease demand pressure on prices.
\end{enumerate}
\end{corrige}

\begin{corrige}[Exercise 6 — The multiplier--accelerator at work]
\begin{enumerate}
\item \textbf{Multiplier:} $k = \dfrac{1}{1 - 0.75} = \dfrac{1}{0.25} = 4$. \textbf{Total change in income:} $\Delta Y = 4 \times 200 = 800$ million FCFA.
\item \textbf{Accelerator:} as national income and demand for goods rise by 800 million FCFA, firms must expand capacity to meet the extra demand. With a capital-output ratio $v$, induced investment is $I_t = v\,(Y_t - Y_{t-1})$: the rise in output induces firms to order new machines, vehicles and buildings, generating a second wave of investment.
\item A \emph{fall} in investment of 200 million FCFA reduces income by $4 \times 200 = 800$ million FCFA (multiplier working in reverse). The fall in income and demand then induces firms to cut investment (accelerator in reverse), which reduces income further, and so on. The cumulative downward interaction of multiplier and accelerator pushes the economy into a recession — the downswing of the trade cycle.
\item \textbf{Limitations in Cameroon:}
\begin{itemize}
\item \textbf{High marginal propensity to import:} much of the extra income is spent on imported goods (fuel, manufactured products), so demand leaks abroad and the effective multiplier is far below 4.
\item \textbf{Large informal sector and supply rigidities:} poor infrastructure and low domestic productive capacity mean extra demand may translate into \emph{price rises} (inflation) rather than extra real output; also, part of consumption is self-produced and outside the monetary circuit.
\end{itemize}
\end{enumerate}
\end{corrige}

\begin{corrige}[Exercise 7 — Classifying unemployment]
\begin{enumerate}
\item \textbf{Structural unemployment.} The textile industry is in permanent decline because of competition from imported second-hand clothing; the worker's skills no longer match the structure of demand, and the job will not return even if aggregate demand recovers.
\item \textbf{Frictional unemployment.} The graduate is between education and a first job; the six-month search period reflects imperfect information and the time needed to match job seekers with vacancies.
\item \textbf{Seasonal unemployment.} Cocoa work is concentrated in harvest periods; between seasons, demand for casual labour falls regularly with the agricultural calendar.
\item \textbf{Cyclical (demand-deficient) unemployment.} The lay-offs result from a general fall in aggregate demand during a severe recession, not from any permanent change in the banking industry.
\item \textbf{Technological unemployment (a form of structural unemployment).} Computerised office software substitutes capital for labour; the secretary's specific skills have been made obsolete by a change in technology.
\end{enumerate}
\end{corrige}

\begin{corrige}[Exercise 8 — Reading a population pyramid]
\begin{enumerate}
\item The pyramid has a \textbf{broad base, narrowing middle and very small top} — the classic shape of a young, expanding population. It reveals a \textbf{high birth rate} (large cohorts aged 0--14) and a \textbf{relatively low life expectancy} (few people survive to ages 65+), typical of a developing country such as Cameroon.
\item \textbf{Total dependency ratio} $=$ youth dependency ratio $+$ old-age dependency ratio $= 80\% + 8\% = 88\%$. There are about 88 dependents for every 100 persons of working age.
\item \textbf{Opportunities:} (i) a large future labour force — a potential ``demographic dividend'' that can raise output and attract labour-intensive investment; (ii) a large and growing domestic market for goods and services, stimulating demand and entrepreneurship. \textbf{Challenges:} (i) heavy pressure on education, health and employment creation — many jobs must be created each year; (ii) a high dependency burden reduces savings and investment per head, and public resources are absorbed by consumption rather than productive investment.
\item \textbf{Policies for a demographic dividend:} (i) massive investment in \textbf{education and vocational/technical training} aligned with labour-market needs (agribusiness, ICT, construction), so that young people become productive workers; (ii) \textbf{job-creating growth policies}: support to second-generation agriculture, agro-industry and SMEs, plus health (including reproductive health/family planning) to progressively lower fertility while raising human capital.
\end{enumerate}
\end{corrige}

\begin{corrige}[Exercise 9 — Essay: Growth and development (25 marks)]
\textbf{Indicative mark scheme:} definitions (4); indicators of each (6); reasons why growth may not bring development (8); evaluation of inclusive-growth policies with Cameroonian application (5); structure and conclusion (2).

\textbf{Introduction and definitions.} \emph{Economic growth} is the sustained increase in a country's real national output (real GDP) over time. \emph{Economic development} is a broader, multidimensional process: rising real income per head \emph{plus} structural transformation, better health and education, reduced poverty and inequality, and wider economic freedom. The statement claims growth is a precondition of development but does not guarantee it.

\textbf{Indicators.} Growth is measured by the rate of change of real GDP (or real GDP per head). Development is measured by wider indicators: GDP per head at PPP, the Human Development Index (income, life expectancy, schooling), poverty and inequality measures (poverty headcount, Gini coefficient), and access to basic services (water, electricity, health).

\textbf{Why growth may fail to translate into development.}
\begin{itemize}
\item \textbf{Enclave growth:} in Cameroon, oil and mineral extraction raise GDP but employ few people and much of the income accrues to foreign companies or a small elite.
\item \textbf{Inequality:} if the gains of growth are concentrated, average income rises while the majority's living standards stagnate.
\item \textbf{Jobless growth:} capital-intensive growth creates few jobs; underemployment in the informal sector remains massive.
\item \textbf{Neglect of human capital:} growth without investment in health and education does not improve HDI.
\item \textbf{Environmental degradation and regional imbalance} may accompany growth, reducing welfare.
\end{itemize}

\textbf{However}, growth is still \emph{necessary}: without rising output, a poor country lacks the resources to finance schools, hospitals and infrastructure; development without any growth is practically impossible.

\textbf{Policies for inclusive growth in Cameroon (evaluation).} (i) Investment in \textbf{education, vocational training and health} raises human capital — effective but slow and budget-costly. (ii) \textbf{Second-generation agriculture and agribusiness} create rural jobs where poverty is concentrated — but requires roads, credit and extension services. (iii) \textbf{Redistribution} (progressive taxation, targeted transfers) — limited by a narrow tax base and a large informal sector. (iv) \textbf{Good governance and anti-corruption} ensure growth revenues finance development — essential but institutionally difficult.

\textbf{Conclusion.} The statement is correct: growth provides the means, development is the end. Cameroon must pursue growth \emph{and} deliberately channel its fruits into human development and structural transformation (as aimed at by Vision 2035) for growth to become development.
\end{corrige}

\begin{corrige}[Exercise 10 — Data response: Inflation in Cameroon (20 marks)]
\begin{enumerate}
\item \textbf{Inflation rates} (4 marks — 1 per correct calculation):
\begin{align*}
\text{Year 2: } & \frac{103-100}{100}\times 100\% = 3\% \\
\text{Year 3: } & \frac{108-103}{103}\times 100\% \approx 4.9\% \\
\text{Year 4: } & \frac{116-108}{108}\times 100\% \approx 7.4\% \\
\text{Year 5: } & \frac{121-116}{116}\times 100\% \approx 4.3\%
\end{align*}
\item \textbf{Type of inflation (4 marks).} This is \textbf{cost-push inflation}, more precisely \emph{imported} cost-push inflation. The depreciation of the currency raises the FCFA price of imported food and fuel; higher world prices raise input costs further. Firms pass these costs into prices. Since ``domestic demand remained weak'', demand-pull inflation is excluded — the pressure comes from the cost/supply side, not from excess demand.
\item \textbf{Two economic costs (4 marks — 2 each).} (i) \textbf{Fall in real purchasing power:} money incomes buy less, hitting poor households hardest since food and fuel weigh heavily in their budgets — poverty and inequality worsen. (ii) \textbf{Uncertainty and menu/shoe-leather costs:} unpredictable prices discourage saving and investment, distort planning by firms, and force frequent price adjustments; the currency's external value and competitiveness may also deteriorate.
\item \textbf{Two policies, evaluated (8 marks — 4 each).}
\begin{itemize}
\item \textbf{Tight monetary policy / defending the exchange rate:} higher interest rates or tighter credit reduce demand and support the currency, limiting imported inflation. \emph{Limitation:} demand is already weak, so tightening may deepen stagnation; moreover, within CEMAC, monetary policy is conducted by BEAC at regional level, so Cameroon alone cannot set interest rates.
\item \textbf{Supply-side and fiscal measures:} subsidies or tax reductions on essential imports (fuel, food), and support to domestic food production (second-generation agriculture) to reduce import dependence. \emph{Limitation:} subsidies strain the budget and are hard to target; raising domestic food supply takes years and requires infrastructure and credit.
\end{itemize}
A good answer notes that cost-push inflation is \emph{hard to control} with demand policies, which risk stagflation.
\end{enumerate}
\end{corrige}

\begin{corrige}[Exercise 11 — Structured question: The trade cycle and unemployment (20 marks)]
\begin{enumerate}
\item \textbf{The four phases (6 marks):}
\begin{itemize}
\item \textbf{Recovery (upswing):} output and employment begin to rise from the trough; demand picks up; prices stabilise then start rising; business confidence returns.
\item \textbf{Boom / peak (prosperity):} output is at or above full-employment capacity; employment is high (near full employment); demand-pull pressure pushes prices up; investment and profits are at their maximum.
\item \textbf{Recession (downswing/contraction):} output falls, firms lay off workers so unemployment rises; demand weakens and price rises slow or prices fall; investment collapses.
\item \textbf{Trough (depression/slump):} output and employment reach their lowest point; spare capacity is large; prices are stagnant or falling; pessimism prevails, before recovery begins.
\end{itemize}
\item \textbf{Multiplier--accelerator explanation of a cumulative recession (6 marks).} An initial fall in investment reduces incomes of workers in the capital-goods sector; they cut consumption, so other firms lose sales and incomes fall again — income falls by a \emph{multiple} of the initial fall in investment: $\Delta Y = k\,\Delta I$ (multiplier in reverse). The fall in income and demand then reduces the need for capacity: with $I_t = v\,(Y_t - Y_{t-1})$, firms slash investment (accelerator in reverse). This further fall in investment multiplies income downward again. The two mechanisms \emph{reinforce} each other, producing a cumulative downward spiral into recession.
\item \textbf{Cyclical unemployment (4 marks).} Cyclical (demand-deficient) unemployment is unemployment caused by insufficient aggregate demand. During the downswing, falling consumption and investment reduce firms' sales; firms cut production and lay off workers across the whole economy, so cyclical unemployment rises. It disappears only when demand recovers.
\item \textbf{Reflationary fiscal policy in Cameroon — assessment (4 marks).} Increased public spending or tax cuts raise aggregate demand; through the multiplier, output and employment recover, reducing cyclical unemployment. \emph{Effectiveness is limited} in Cameroon by: a high marginal propensity to import (leakages reduce the multiplier); supply bottlenecks (poor infrastructure may turn extra demand into inflation); a large informal sector less responsive to demand stimulus; and financing constraints (budget deficits, debt, limited monetary autonomy within CEMAC). Conclusion: useful but must be combined with supply-side measures.
\end{enumerate}
\end{corrige}

\begin{corrige}[Exercise 12 — Essay: Privatisation (25 marks)]
\textbf{Indicative mark scheme:} definition of privatisation and of the two efficiencies (5); theoretical case for privatisation (5); counter-arguments: natural monopoly, strategic industries, regulation (6); Cameroonian evidence — SONEL, CAMTEL (6); judgement and conclusion (3).

\textbf{Introduction.} \emph{Privatisation} is the transfer of ownership of state-owned enterprises to the private sector (sale of shares, concessions, management contracts). \emph{Productive efficiency} means producing at minimum average cost; \emph{allocative efficiency} means producing the output where price equals marginal cost, matching consumer wants.

\textbf{The case that privatisation improves efficiency.} Private owners face the profit motive and the discipline of capital markets: incentives to cut costs (productive efficiency), innovate, respond to consumers and end political interference and soft budget constraints. Competition introduced with privatisation pushes prices towards costs (allocative efficiency). In Cameroon, the liberalisation of \textbf{mobile telephony} (MTN, Orange operating alongside CAMTEL) dramatically expanded access, lowered prices and improved service quality — clear efficiency gains where competition was possible.

\textbf{The case against ``always''.}
\begin{itemize}
\item \textbf{Natural monopolies} (electricity networks, fixed-line infrastructure, water): a single firm minimises cost, so privatisation without competition merely creates a \emph{private monopoly} that may restrict output and charge monopoly prices — allocative efficiency may \emph{worsen}.
\item \textbf{Strategic industries and public-service obligations:} private firms may neglect unprofitable rural areas, undermining equity and development objectives.
\item \textbf{Regulatory weakness:} where the regulator is weak or captured, contracts are renegotiated opportunistically.
\end{itemize}

\textbf{Cameroonian evidence.} \textbf{SONEL}, privatised in 2001 to AES (AES-SONEL), suffered disputes over investment and tariffs; the arrangement broke down and the State regained control in 2014 (now ENEO). This shows privatisation of a natural monopoly with weak regulation does not guarantee efficiency. \textbf{CAMTEL}'s partial privatisation was reversed, while competition in the contestable mobile segment succeeded — confirming that \emph{competition and regulation}, not ownership change alone, drive efficiency.

\textbf{Conclusion.} The statement is too strong. Privatisation improves efficiency \emph{when} markets are competitive or contestable and regulation is strong and independent. For natural monopolies and strategic sectors, privatisation without regulation can replace a public monopoly with a private one. Efficiency depends on market structure and institutions, not ownership alone.
\end{corrige}

\begin{corrige}[Exercise 13 — Essay: Industrialisation strategies (25 marks)]
\textbf{Indicative mark scheme:} definitions of both strategies (5); advantages and limits of each (10); lessons from the NICs (4); application to Cameroon with second-generation agriculture/agribusiness (4); conclusion (2).

\textbf{Definitions.} \emph{Import-substitution industrialisation (ISI)} promotes domestic production of previously imported manufactured goods behind protective barriers (tariffs, quotas). \emph{Export-oriented industrialisation (EOI)} promotes industries producing manufactures for world markets, using openness, competitiveness and often incentives to exporters.

\textbf{ISI — advantages:} protects infant industries; saves foreign exchange; builds domestic capacity and reduces dependence on imports; easier to start since the domestic market already exists. \textbf{Limits:} protection breeds inefficiency and rent-seeking; the domestic market is small (limited economies of scale); it requires imported capital goods, so the foreign-exchange saving is often disappointing; Latin American and many African experiences (including Cameroon's earlier attempts) ended in inefficient, high-cost industries.

\textbf{EOI — advantages:} access to huge world markets permits economies of scale; foreign competition enforces efficiency and technological upgrading; exports earn foreign exchange; it is labour-intensive at first, creating jobs. \textbf{Limits:} requires competitiveness (infrastructure, skills, stable macroeconomic environment); exposes the economy to world-market shocks; initial reliance on low wages may be hard to sustain.

\textbf{Lessons from the NICs} (South Korea, Taiwan, Singapore, later Malaysia, Vietnam): they often began with a \emph{brief} ISI phase, then moved decisively to EOI, investing heavily in education, infrastructure and export discipline, with the State guiding but not shielding firms from competition permanently. Their success came from combining openness with strong human capital and capable institutions.

\textbf{Application to Cameroon.} Pure EOI is difficult today: Cameroon lacks the infrastructure, skills and cost-competitiveness of the Asian NICs. Pure ISI risks repeating past inefficiency. A \textbf{pragmatic mix} is best suited: \emph{selective} import substitution anchored on \textbf{second-generation agriculture and agribusiness} — modernising agriculture (mechanisation, improved inputs, irrigation) and developing local processing of cocoa, coffee, cotton and food crops. This builds on Cameroon's comparative advantage, creates rural jobs, saves foreign exchange on food imports, and can progressively become \emph{export-oriented agro-industry} (processed cocoa, textiles) within CEMAC and the AfCFTA markets.

\textbf{Conclusion.} For Cameroon, neither textbook strategy suffices alone. The NIC experience suggests sequencing: build competitiveness through agriculture-based industrialisation and human-capital investment (a targeted, time-limited ISI), then open progressively to export markets — an agribusiness-led path toward EOI.
\end{corrige}

\begin{corrige}[Exercise 14 — Data response: Population and development (20 marks)]
\begin{enumerate}
\item \textbf{Country A has a youthful population} (4 marks): $43\%$ under 15, only $4\%$ aged 65+, and a high growth rate of $2.6\%$ per year — signs of high fertility and a broad-based pyramid. \textbf{Country B has an ageing population}: only $16\%$ under 15, $19\%$ aged 65+, and near-stagnant growth of $0.4\%$ — low fertility and rising life expectancy.
\item \textbf{Economic challenges (6 marks).} \emph{Country A (youthful):} very high youth dependency ratio; heavy spending needed on schools, health and child services; the need to create enormous numbers of jobs each year; pressure on food, housing and savings (low saving per head limits investment). \emph{Country B (ageing):} rising old-age dependency — pension and health-care costs explode; shrinking labour force may slow growth and innovation; higher taxes on a smaller working population; possible labour shortages.
\item \textbf{Optimum population analysis (5 marks).} The optimum population is the size that maximises output per head given resources, capital and technology. Rapid growth in Country A may \emph{hinder} development if population grows faster than capital and jobs: diminishing returns set in, income per head stagnates, dependency rises — the country is then \emph{over-populated relative to its capital}. But it may \emph{help} if the growing labour force is educated and employed: a larger market allows economies of scale, and the bulge of young workers can yield a \emph{demographic dividend}, raising output per head — moving the country toward a new, higher optimum as technology and capital accumulate.
\item \textbf{Two population policies for Country A, evaluated (5 marks).}
\begin{itemize}
\item \textbf{Family planning and reproductive-health education} to lower fertility voluntarily: reduces the dependency burden over time; \emph{limitation} — results are slow and may face cultural or religious resistance.
\item \textbf{Investment in education (especially of girls) and job creation}: educated women marry later and have fewer children, while skills convert youth into productive workers; \emph{limitation} — costly, requires sustained budgets and complementary job-creating growth, otherwise it produces educated unemployed.
\end{itemize}
\end{enumerate}
\end{corrige}

\begin{corrige}[Exercise 15 — Short-answer set: Deflation and reflation (15 marks)]
\begin{enumerate}
\item \textbf{Deflation vs disinflation (3 marks).} \emph{Deflation} is a sustained \emph{fall in the general price level} — the inflation rate is negative (e.g. $-2\%$). \emph{Disinflation} is a \emph{fall in the rate of inflation} while prices still rise — inflation is positive but lower (e.g. from $6\%$ to $3\%$). In deflation the price level falls; in disinflation it rises more slowly.
\item \textbf{Two costs of prolonged deflation (4 marks).} (i) \textbf{Postponement of spending:} expecting lower prices tomorrow, consumers and firms delay purchases and investment, reducing aggregate demand, output and employment — a deflationary spiral. (ii) \textbf{Rise in the real burden of debt:} with falling prices and money incomes, the real value of fixed nominal debts increases, causing bankruptcies, bank losses and financial instability (debt deflation).
\item \textbf{AD/AS analysis of reflation (5 marks).} A reflationary policy (higher government spending, lower taxes, easier money) shifts aggregate demand to the right, from $AD_1$ to $AD_2$. The economy moves from the deflationary equilibrium ($Y_1$, below full-employment output $Y_f$, with falling prices) toward full employment at $Y_f$, closing the deflationary gap and restoring a stable, mildly positive price level.
\begin{popfigure}
\begin{tikzpicture}[scale=0.9]
\draw[->, thick] (0,0) -- (6.4,0) node[below left, font=\small] {Real output $Y$};
\draw[->, thick] (0,0) -- (0,4.6) node[above left, font=\small] {Price level};
\draw[popblue, very thick, domain=0.6:5.4] plot (\x,{4.3-0.65*\x}) node[right, font=\small] {$AD_1$};
\draw[poporange, very thick, domain=1.4:6.2] plot (\x,{5.1-0.65*\x}) node[right, font=\small] {$AD_2$};
\draw[popgreen, very thick, domain=0.5:5.8] plot (\x,{0.45+0.6*\x}) node[right, font=\small] {$AS$};
\draw[dashed, popmuted] (4.9,0) node[below, font=\small] {$Y_f$} -- (4.9,3.4);
\draw[dashed, popmuted] (2.9,0) node[below, font=\small] {$Y_1$} -- (2.9,2.2);
\draw[<->, popdark, thick] (2.9,0.9) -- (4.9,0.9) node[midway, below, font=\small] {deflationary gap};
\end{tikzpicture}
\legende{Reflationary policy shifts AD rightward, closing the deflationary gap.}
\end{popfigure}
\textbf{Risk:} if the stimulus is excessive or badly timed, demand overshoots capacity and the policy generates \emph{demand-pull inflation} (and possibly a larger budget deficit and public debt).
\item \textbf{Nationalisation of failing strategic firms (3 marks).} When a strategic private firm (e.g. electricity, rail, a major bank) fails, nationalisation may be proposed to keep an essential service running and protect jobs and the wider economy. \emph{Advantage:} the State can guarantee continuity of an essential service and pursue public-interest goals (universal access) rather than profit alone. \emph{Disadvantage:} state ownership may bring inefficiency, political interference and a heavy burden on the public budget (subsidies to loss-making firms).
\end{enumerate}
\end{corrige}

\newpage

\coursec{Summary sheet}

\begin{popfigure}
\begin{center}\tcbox[colback=popblue,colframe=popblue,coltext=white,arc=9pt,boxsep=4pt,left=6pt,right=6pt]{\bfseries\small Economic Growth \& Development (II)}\end{center}
\vspace{-2pt}
\begin{tcbraster}[raster columns=3,raster equal height=rows,raster column skip=6pt,raster row skip=6pt,enhanced,blankest]
\begin{tcolorbox}[enhanced,colback=popblue,colframe=popblue,coltext=white,boxrule=0.9pt,arc=3pt,left=4pt,right=4pt,top=3pt,bottom=3pt,boxsep=1pt,halign=left]\bfseries\scriptsize Economic Development \textit{Meaning, HDI, HPI, SDGs, Characteristics}\end{tcolorbox}
\begin{tcolorbox}[enhanced,colback=poporange,colframe=poporange,coltext=white,boxrule=0.9pt,arc=3pt,left=4pt,right=4pt,top=3pt,bottom=3pt,boxsep=1pt,halign=left]\bfseries\scriptsize Industrialisation Dynamics \textit{Process, De/Re-industrialisation, NICs Strategies}\end{tcolorbox}
\begin{tcolorbox}[enhanced,colback=popgreen,colframe=popgreen,coltext=white,boxrule=0.9pt,arc=3pt,left=4pt,right=4pt,top=3pt,bottom=3pt,boxsep=1pt,halign=left]\bfseries\scriptsize Agricultural Transformation \textit{2nd Gen Agriculture, Agribusiness, Cameroon}\end{tcolorbox}
\begin{tcolorbox}[enhanced,colback=poppurple,colframe=poppurple,coltext=white,boxrule=0.9pt,arc=3pt,left=4pt,right=4pt,top=3pt,bottom=3pt,boxsep=1pt,halign=left]\bfseries\scriptsize Business Cycle \& Govt Objectives \textit{Phases, Growth, Employment, Stability, External Balance}\end{tcolorbox}
\begin{tcolorbox}[enhanced,colback=poppink,colframe=poppink,coltext=white,boxrule=0.9pt,arc=3pt,left=4pt,right=4pt,top=3pt,bottom=3pt,boxsep=1pt,halign=left]\bfseries\scriptsize Unemployment \textit{Types, Causes, Measurement, Cameroon Context}\end{tcolorbox}
\begin{tcolorbox}[enhanced,colback=popgold,colframe=popgold,coltext=white,boxrule=0.9pt,arc=3pt,left=4pt,right=4pt,top=3pt,bottom=3pt,boxsep=1pt,halign=left]\bfseries\scriptsize Price Stability \textit{Inflation \& Deflation: Types, Causes, Effects, Policies}\end{tcolorbox}
\begin{tcolorbox}[enhanced,colback=popmuted,colframe=popmuted,coltext=white,boxrule=0.9pt,arc=3pt,left=4pt,right=4pt,top=3pt,bottom=3pt,boxsep=1pt,halign=left]\bfseries\scriptsize Structural Policies \& Population \textit{Privatisation, Nationalisation, Regional Policy, Optimum Population}\end{tcolorbox}
\end{tcbraster}

\legende{Chapter Mind Map: Economic Growth and Development (Continued)}
\end{popfigure}

\begin{bilan}

\courssub{Key Definitions}
\begin{itemize}
    \item \cle{Economic Development}: A multidimensional process involving major changes in social structures, popular attitudes, and national institutions, as well as the acceleration of economic growth, the reduction of inequality, and the eradication of absolute poverty (Todaro). \textbf{Not just GDP growth.}
    \item \cle{Human Development Index (HDI)}: Geometric mean of normalised indices for Life Expectancy, Education (Mean/Expected years schooling), and GNI per capita (PPP \$). Range $[0,1]$.
    \item \cle{Industrialisation}: Structural shift where manufacturing sector share of GDP and employment rises. \cle{Deindustrialisation}: Fall in manufacturing share (premature if at low income). \cle{Reindustrialisation}: Policy-driven revival of manufacturing (high-tech/green).
    \item \cle{NICs (Newly Industrialised Countries)}: Economies (e.g., South Korea, Singapore) that transitioned via export-oriented industrialisation (EOI), high savings/investment, and human capital focus.
    \item \cle{Second Generation Agriculture}: Modern, commercial, technology-driven agriculture focusing on value chains, processing (agribusiness), and export competitiveness (Cameroon's PDRA strategy).
    \item \cle{Business Cycle}: Short-run fluctuations in real GDP around potential output. Phases: Expansion $\to$ Peak $\to$ Contraction/Recession $\to$ Trough $\to$ Recovery.
    \item \cle{Unemployment Rate}: $U/L \times 100$ where $U$=unemployed (actively seeking), $L$=labour force. Types: Frictional, Structural, Cyclical, Seasonal, Disguised.
    \item \cle{Inflation}: Sustained rise in general price level (CPI). \cle{Demand-pull} (AD $>$ AS), \cle{Cost-push} (AS shifts left), \cle{Imported} (exchange rate). \cle{Deflation}: Sustained fall in prices (risk of debt deflation spiral).
    \item \cle{Privatisation}: Transfer of state-owned enterprises (SOEs) to private sector (sale, concession, liquidation). \cle{Nationalisation}: Reverse transfer (strategic sectors). \cle{Regional Policy}: Correcting spatial imbalances (growth poles vs. depressed areas).
    \item \cle{Optimum Population}: Population size maximising output per capita (or welfare) given resources/technology. \cle{Demographic Dividend}: Window where working-age population share peaks.
\end{itemize}

\courssub{Key Laws, Formulas \& Theorems}
\begin{itemize}
    \item \textbf{HDI Formula (2010+)}: $HDI = \sqrt[3]{I_{Health} \times I_{Education} \times I_{Income}}$. $I_{dim} = \frac{Actual - Min}{Max - Min}$.
    \item \textbf{Unemployment Rate}: $u = \frac{N_{unemployed}}{Labour Force} \times 100$. \textbf{Labour Force Participation Rate}: $LFPR = \frac{Labour Force}{Working Age Pop} \times 100$.
    \item \textbf{Inflation Rate (CPI)}: $\pi_t = \frac{CPI_t - CPI_{t-1}}{CPI_{t-1}} \times 100$. \textbf{Real Interest Rate} (Fisher): $r \approx i - \pi^e$.
    \item \textbf{Okun's Law} (Empirical): $\Delta u \approx -0.5 (\Delta Y - \Delta Y^*)$. 1\% rise in unemployment $\approx$ 2\% GDP gap.
    \item \textbf{Phillips Curve} (Short-run): Inverse relation between $\pi$ and $u$. Long-run: Vertical at NAIRU (Non-Accelerating Inflation Rate of Unemployment).
    \item \textbf{Harrod-Domar Growth}: $g = s/v$ (Savings rate / Capital-Output ratio). \textbf{Solow}: Diminishing returns to capital $\to$ steady state; TFP (technology) drives long-run growth.
    \item \textbf{Lewis Model}: Surplus labour in traditional sector $\to$ modern sector at constant wage $\to$ turning point $\to$ wages rise.
\end{itemize}

\courssub{Methods \& Procedures}
\begin{enumerate}
    \item \textbf{Calculating HDI}: Get Life Exp, Mean/Exp Schooling, GNI pc (PPP). Compute dimension indices using goalposts (e.g., Life Exp: 20/85). Take geometric mean. \textit{Exam: Show steps clearly.}
    \item \textbf{Analysing Business Cycle}: Identify phase using real GDP growth, unemployment, inflation, confidence. Prescribe policy: Expansionary (Recession) vs Contractionary (Overheating).
    \item \textbf{Distinguishing Inflation Types}: Plot AD/AS. Demand-pull: AD shifts right (Price $\uparrow$, Output $\uparrow$). Cost-push: AS shifts left (Price $\uparrow$, Output $\downarrow$ = Stagflation). Policy differs!
    \item \textbf{Evaluating Privatisation}: Check: Competitive market? Regulatory capacity? Strategic importance? Social safety nets? Method: Trade sale $\to$ efficiency; Voucher $\to$ speed; Concession $\to$ infrastructure.
    \item \textbf{Population Policy Diagnosis}: Compare Crude Birth/Death Rates, TFR, Dependency Ratio. High TFR + High Youth Dependency $\to$ Invest in education/family planning (Cameroon). Ageing $\to$ Raise retirement age, automation, migration.
\end{enumerate}

\courssub{Common Mistakes \& Traps}
\begin{itemize}
    \item \textbf{Confusing Growth \& Development}: "GDP rose 5\% $\therefore$ development occurred." \textbf{Wrong}. Check HDI, inequality (Gini), poverty headcount.
    \item \textbf{HDI vs HPI (Human Poverty Index)}: HDI measures \textit{achievements}; HPI measures \textit{deprivations} (not in current UNDP reports but syllabus relevant). Don't mix formulas.
    \item \textbf{Unemployment Measurement}: "Discouraged workers are unemployed." \textbf{Wrong}. They are \textit{inactive} (not seeking). This lowers measured $u$ rate artificially.
    \item \textbf{Inflation Causes}: Attributing \textit{all} inflation to money printing. In Cameroon/CEMAC, supply shocks (fuel, food) and imported inflation (CFA parity) are major drivers.
    \item \textbf{Privatisation $\neq$ Liberalisation}: Privatisation = ownership change. Liberalisation = entry deregulation. A privatised monopoly (e.g., SONATEL/Orange early days) may still need regulation.
    \item \textbf{Optimum Population $\neq$ Maximum Population}: Optimum maximises \textit{per capita} output/welfare. Beyond it, diminishing returns lower living standards (Malthusian trap risk).
    \item \textbf{NICs Strategy}: "They just opened up." \textbf{Wrong}. Strategic state intervention (credit allocation, chaebols/keiretsu, education) \textit{preceded} liberalisation. Sequencing matters.
\end{itemize}

\courssub{GCE Advanced Level Tips}
\begin{itemize}
    \item \textbf{Essay Structure (25 marks)}: \textbf{Intro}: Define key terms, context (Cameroon/CEMAC), thesis. \textbf{Body}: 3-4 developed arguments (Point $\to$ Explain $\to$ Example $\to$ Link). Use AD/AS, Lorenz, Phillips diagrams. \textbf{Evaluation}: "However...", "Depends on...", "Short-run vs Long-run". \textbf{Conclusion}: Judgement + Policy recommendation.
    \item \textbf{Cameroon Context is Mandatory}: Cite \cle{PDRA} (Agriculture), \cle{SONARA} (Privatisation/Subsidies), \cle{CEMAC} convergence criteria (Inflation $<3\%$, Debt/GDP $<70\%$), \cle{AfCFTA} (Industrialisation/Trade), \cle{Demographic Dividend} potential (median age $\sim$18).
    \item \textbf{Diagram Checklist}: Business Cycle (label axes: Time, Real GDP), AD/AS for Inflation types, Lorenz Curve (Gini), Phillips Curve (SR/LR), Production Function (Solow/Lewis).
    \item \textbf{Command Words}: "Assess" / "Evaluate" $\to$ Need 2-sided analysis + Judgement. "Analyse" $\to$ Step-by-step causal chain. "Distinguish" $\to$ Table of differences. "Explain" $\to$ Mechanism.
    \item \textbf{Time Management}: 3 hours / 4 questions $\approx$ 45 min each. Plan 5 min, Write 35 min, Check 5 min.
    \item \textbf{Current Affairs}: Mention post-COVID recovery, Russia-Ukraine impact on imported inflation/food security, CFA Franc reform debates (Eco currency), Climate change adaptation in agriculture.
\end{itemize}

\end{bilan}

\findecours

\end{document}
